% paralleles.tex — f165v, vv. 252–293 : loci paralleli. Généré par .travail/f165v_lp_redac2/rediger.py
% à partir des collectes vérifiées (paralleles_252-272.json, paralleles_273-293.json) et des parallèles
% d'Austin–Olson 2004 retrouvés dans les sources ; règles : edition/CONVENTIONS.md § 4
\lp{252 λαμβάνετε καὶ χρῆσθ᾽]}{Ar.~\textit{Thesm.}~219–220 χρῆσόν τί νυν ἡμῖν ξυρόν. — αὐτὸς λάμβανε / ἐντεῦθεν ἐκ τῆς ξυροδόκης ; Ar.~\textit{Ran.}~735 χρῆσθε τοῖς χρηστοῖσιν αὖθις}
\lp{252 λαμβάνετε … οὐ φθονῶ (κοινὰ τὰ φίλων)]}{Alciphr.~2.12 Schepers (=~3.15) εἰ οὖν δέοιο, προθύμως λάμβανε. τὸ γὰρ κοινὰ τὰ τῶν φίλων οὐχ ἥκιστα τοῖς ἀγροῖς ἐμφιλοχωρεῖν ὀφείλει ; Eur.~\textit{Andr.}~376–377 φίλων γὰρ οὐδὲν ἴδιον, οἵτινες φίλοι / ὀρθῶς πεφύκασ᾽, ἀλλὰ κοινὰ χρήματα ; Eur.~\textit{Or.}~735 συγκατασκάπτοις ἂν ἡμᾶς· κοινὰ γὰρ τὰ τῶν φίλων ; Zen.~4.79 Κοινὰ τὰ φίλων: Τίμαιος φησὶν ὅτι προσιόντας Πυθαγόρᾳ μαθητὰς περὶ τὴν Ἰταλίαν ἔπειθεν ὁ φιλόσοφος κοινὰς τὰς οὐσίας ποιεῖσθαι ; Tosi n°~433 Communia esse amicorum inter se omnia}
\lp{252 οὐ φθονῶ]}{Ar.~\textit{Thesm.}~249–250 Ἀγάθων, ἐπειδὴ σαυτὸν ἐπιδοῦναι φθονεῖς, / ἀλλ᾽ ἱμάτιον γοῦν χρῆσον ἡμῖν τουτῳὶ ; Eur.~\textit{HF}~333 κοσμεῖσθ᾽ ἔσω μολόντες· οὐ φθονῶ πέπλων ; Eur.~\textit{Hec.}~238 ἔξεστ᾽, ἐρώτα· τοῦ χρόνου γὰρ οὐ φθονῶ ; \textit{Od.}~17.400 δός οἱ ἑλών· οὔ τοι φθονέω· κέλομαι γὰρ ἐγώ γε ; Pl.~\textit{Prt.}~320c ἀλλ᾽, ὦ Σώκρατες, ἔφη, οὐ φθονήσω· ἀλλὰ πότερον ὑμῖν, ὡς πρεσβύτερος νεωτέροις, μῦθον λέγων ἐπιδείξω ; Eur.~\textit{Hipp.}~20 τούτοισι μέν νυν οὐ φθονῶ· τί γάρ με δεῖ ; Aesch.~\textit{PV}~626 ἀλλ᾽ οὐ μεγαίρω τοῦδέ σοι δωρήματος ; Soph.~\textit{Ant.}~553 σῶσον σεαυτήν· οὐ φθονῶ σ᾽ ὑπεκφυγεῖν ; Pl.~\textit{La.}~200b διδάξω καὶ σέ, καὶ οὐ φθονήσω· δοκεῖς γάρ μοι καὶ μάλα σφόδρα δεῖσθαι μαθεῖν ; X.~\textit{Cyr.}~8.4.16 ὧν ἐγώ σοι οὐ φθονήσω, ἢν τὴν θυγατέρα μου γυναῖκα λαμβάνῃς ; Philo \textit{Abr.}~204 κεκτημένος δ᾽ ὅμως οὐ φθονῶ τοῖς ἀξίοις χρῆσθαι ; Hsch.~μ~451 μεγαίρειν· φθονεῖν, ζηλοῦν. στερίσκειν ; Hsch.~μ~452 μεγαίρω· φθονῶ ; Hsch.~φ~424 φθονεῖς· ζηλοῖς, βασκαίνεις ; Suid.~φ~511 Φθονῶ σοι· δοτικῇ. φθονῶ σοι τοῦδε ; Eur.~\textit{Ba.}~1005 τὸ σοφὸν οὐ φθονῶ ; Eur.~\textit{Med.}~312 καὶ νῦν τὸ μὲν σὸν οὐ φθονῶ καλῶς ἔχειν ; \textit{Od.}~1.346 μῆτερ ἐμή, τί τ᾽ ἄρα φθονέεις ἐρίηρον ἀοιδὸν ; Nonn.~\textit{D.}~27.127 οὐ φθονέω Μορρῆι ; X.~\textit{Smp.}~3.5 τότε κἀγὼ οὐ φθονήσω εἰπεῖν τὴν τέχνην ; D.~42.22 καὶ οὐ φθονῶ· δεῖ γὰρ τοὺς εὐπόρους χρησίμους αὑτοὺς παρέχειν ; Hsch.~ν~283 νεμεσίζομαι· μέμφομαι. φθονῶ ; Ammon.~\textit{Diff.}~211 Nickau ζῆλος φθόνου διαφέρει. ζῆλος μὲν γάρ ἐστιν ἡ δι᾽ ἐπιθυμίαν γινομένη μίμησις δοκοῦντός τινος καλοῦ, φθόνος δὲ βασκανία τις τῶν ἄλλοις μὲν προσόντων ἀγαθῶν, ἡμῖν δ᾽ οὔ}
\lp{253 ὅ τι;]}{Ar.~\textit{Av.}~1640 τί δαὶ ποιῶμεν; — ὅ τι; διαλλαττώμεθα ; Ar.~\textit{Ach.}~106 ὅ τι; χαυνοπρώκτους τοὺς Ἰάονας λέγει ; Ar.~\textit{Ach.}~959 Δικαιόπολι. — τίς ἔστι; τί με βωστρεῖς; — ὅ τι ; Ar.~\textit{Nub.}~753 τί δῆτα τοῦτ᾽ ἂν ὠφελήσειέν σ᾽; — ὅ τι ; Ar.~\textit{Vesp.}~957 ὅ τι; σοῦ προμάχεται καὶ φυλάττει τὴν θύραν ; Ar.~fr.~305.4~K.-A. (=~294~Kock) σοὶ δ᾽ ὄνομα δὴ τί ἐστιν; — ὅ τι; Γεωργία ; Ar.~\textit{Eq.}~742 ὅ τι ; Ar.~\textit{Av.}~164 τί σοι πιθώμεσθ᾽; — ὅ τι πίθησθε; πρῶτα μὲν}
\lp{253 τὸν κροκωτόν]}{Ar.~\textit{Lys.}~44 αἳ καθήμεθ᾽ ἐξηνθισμέναι, / κροκωτοφοροῦσαι καὶ κεκαλλωπισμέναι ; Ar.~\textit{Thesm.}~138 λαλεῖ κροκωτῷ; τί δὲ λύρα κεκρυφάλῳ ; Ar.~\textit{Thesm.}~941 ἵνα μὴ ᾽ν κροκωτοῖς καὶ μίτραις γέρων ἀνὴρ / γέλωτα παρέχω ; Ar.~\textit{Thesm.}~945 ἰαππαπαιάξ· ὦ κροκώθ᾽ οἷ᾽ εἴργασαι ; Ar.~\textit{Thesm.}~1044 ὃς ἔμ᾽ ἀπεξύρησε πρῶτον, / ὃς ἐμὲ κροκόεντ᾽ ἐνέδυσεν ; Ar.~\textit{Ran.}~46 ἀλλ᾽ οὐχ οἷός τ᾽ εἴμ᾽ ἀποσοβῆσαι τὸν γέλων / ὁρῶν λεοντῆν ἐπὶ κροκωτῷ κειμένην ; Ar.~\textit{Eccl.}~332 ἀλλὰ τῆς γυναικὸς ἐξελήλυθα / τὸ κροκωτίδιον ἀμπισχόμενος οὑνδύεται ; Cratin.~fr.~40~K.-A. (=~38~Kock ; Ath.~11.474e–475a) στολὴν δὲ δὴ τίν᾽ εἶχεν; τοῦτό μοι φράσον. / θύρσον, κροκωτόν, ποικίλον, καρχήσιον ; Arar.~fr.~4~Kock (Poll.~7.65) παρθένος δ᾽ εἶναι δοκεῖ, / φορῶν κροσωτοὺς καὶ γυναικείαν στολήν ; Eur.~\textit{Ba.}~828–830 τίνα στολήν; ἦ θῆλυν; ἀλλ᾽ αἰδώς μ᾽ ἔχει. / … / στολὴν δὲ τίνα φῂς ἀμφὶ χρῶτ᾽ ἐμὸν βαλεῖν ; Verg.~\textit{Aen.}~9.614–616 Vobis picta croco et fulgenti murice vestis, / desidiae cordi, iuvat indulgere choreis, / et tunicae manicas et habent redimicula mitrae ; Juv.~6, O~22–23 distinctus croceis et reticulatus adulter, / suspectus tibi sit quanto vox mollior ; D.Chr.~32.94 τὸν δὲ Ἡρακλέα τοιοῦτον ὁρῶσι γελοῖον δοκεῖ, παραφερόμενον, καὶ καθάπερ εἰώθασιν, ἐν κροκωτῷ ; Ael.~\textit{VH}~7.9 οὐ κροκωτοῦ οὐ ταραντινιδίου οὐκ ἀναβολῆς οὐκ ἐγκύκλου οὐ κεκρυφάλου οὐ καλύπτρας ; Ath.~4.155c Πολυσπέρχοντά φησιν εἰ μεθυσθείη καίτοι πρεσβύτερον ὄντα ὀρχεῖσθαι, … καὶ ἐνδυόμενον αὐτὸν κροκωτὸν καὶ ὑποδούμενον Σικυώνια διατελεῖν ὀρχούμενον ; Apul.~\textit{Met.}~8.27 facie caenoso pigmento delita et oculis obunctis graphice prodeunt; mitellis et crocotis et carbasinis et bombycinis iniecti ; Σ Ar.~\textit{Ran.}~46 Διονυσιακὸν φόρεμα ὁ κροκωτός. … γυναικείῳ ἐκ πήνης ἐνδύματι ; Poll.~7.55 αὗται μὲν ἀνδρῶν, γυναικῶν δὲ κροκωτὸς κροκώτιον, παραλουργίς, ὀμφάκινον ; Suid.~α~2657 οὗτος ὑποδήμασι Μιλησίοις πρῶτος ἐχρήσατο καὶ κροκωτὸν ἐν τῷ Κωμαστῇ περιεβάλλετο ἱμάτιον ; Suid.~κ~2461 Κροκωτός: εἶδος χιτῶνος. ὁ δὲ ἤσθητο κροκωτὸν χιτῶνα, ὑπὲρ τοῦ πλείονας εἰς αὐτὸν ἐπιστρέφειν. ὅτι ὁ μὲν κροκωτὸς ἔνδυμά ἐστι, τὸ δὲ ἔγκυκλον ἱμάτιον ; Ar.~\textit{Lys.}~47 τὰ κροκωτίδια καὶ τὰ μύρα χαἰ περιβαρίδες ; Ar.~\textit{Lys.}~51 κροκωτὸν ἄρα νὴ τὼ θεὼ ᾽γὼ βάψομαι ; Ar.~\textit{Lys.}~219–220 κροκωτοφοροῦσα καὶ κεκαλλωπισμένη,— ; Ar.~\textit{Lys.}~645 κᾆτ᾽ ἔχουσα τὸν κροκωτὸν ἄρκτος ἦ Βραυρωνίοις ; Ar.~\textit{Thesm.}~1220 κροκῶτ᾽ ἔκοντο τὴ γέροντο ; Ar.~\textit{Eccl.}~318–319 τουτὶ τὸ τῆς γυναικὸς ἡμιδιπλοίδιον, / καὶ τὰς ἐκείνης Περσικὰς ὑφέλκομαι ; Ar.~\textit{Eccl.}~374 τὸ τῆς γυναικὸς δ᾽ ἀμπέχει χιτώνιον ; Ar.~\textit{Eccl.}~879 ἕστηκα καὶ κροκωτὸν ἠμφιεσμένη ; Aesch.~\textit{Ag.}~239 κρόκου βαφὰς δ᾽ ἐς πέδον χέουσα ; Plu.~\textit{Mor.}~785e ἐν Ὀμφάλης κροκωτοφόρον ; D.Chr.~4.110 μύρου καὶ οἴνου ἀποπνέων ἐν κροκωτῷ μετὰ πολλοῦ καὶ ἀτάκτου γέλωτος ; Luc.~\textit{Bacch.}~2 νάρθηκι ἐπερειδόμενον, ἐπ᾽ ὄνου τὰ πολλὰ ἱππεύοντα, ἐν κροκωτῷ καὶ τοῦτον ; Luc.~\textit{Hist.Conscr.}~10 αὐτὸν δὲ ἐν κροκωτῷ καὶ πορφυρίδι ἔρια ξαίνοντα ; Philostr.~\textit{VA}~4.21 κροκωτοὶ δὲ ὑμῖν καὶ ἁλουργία καὶ κοκκοβαφία τοιαύτη πόθεν ; Ath.~5.198c χιτῶνα πορφυροῦν ἔχον διάπεζον καὶ ἐπ᾽ αὐτοῦ κροκωτὸν διαφανῆ ; Ath.~12.519c ἐπόμπευον ἔχοντες κροκωτοὺς ἐπὶ τοῖς θώραξιν ; Clem.Al.~\textit{Paed.}~2.10 τόν τε μυροβαφῆ ἐκεῖνον κροκωτὸν ; Apul.~\textit{Apol.}~13 uti me consuesse tragoedii syrmate, histrionis crocota ; Plot.~3.2.17 ὥσπερ οἱ τῇδε ὑποκριταὶ τὰ προσωπεῖα, τὴν ἐσθῆτα, τοὺς κροκωτοὺς καὶ τὰ ῥάκη ; Polyaen.~5.3 αὐτὸς δὲ ἤδη μεθυόντων ἐς μέσους παρελθὼν, κροκωτὸν ἐνδὺς, Ταραντῖνον περιβαλόμενος ηὔλησεν, ἐκιθάρισεν, ὡρχήσατο ; Σ Ar.~\textit{Lys.}~51 κροκωτὸν ἄρα: Εἰ ἄρα οὕτως ἔχει, καλλωπισθήσομαι καὶ γραῦς οὖσα ; Σ Ar.~\textit{Lys.}~645 αἱ ἀρκτευόμεναι δὲ τῇ θεῷ κροκωτὸν ἠμφιέννυντο ; Σ Ar.~\textit{Nub.}~70 ξύστις λέγεται τὸ κροκωτὸν ἱμάτιον, ὃ οἱ ἡνίοχοι μέχρι τοῦ νῦν φοροῦσι πομπεύοντες. χρῶνται δὲ αὐτῷ καὶ οἱ τραγικοὶ βασιλεῖς ; \textit{idem} Suid.~ξ~169 (ξυστὶς ; μέχρι νῦν) ; Σ Ar.~\textit{Thesm.}~689 ἀντὶ δὲ παιδίου ἀσκὸν οἴνου λαμβάνει, καὶ αὐτὸν ἐνδύει κροκωτὸν ; Σ Ar.~\textit{Thesm.}~1044 ἀντὶ τοῦ κροκωτὸν ἀμφέδυσεν ; Σ Soph.~\textit{OC}~684 ὅτι κἀν τῇ Νιόβῃ ὁ Σοφοκλῆς τὸν κρόκον ἄντικρυς τῇ Δήμητρι ἀνατίθεται ; Hsch.~α~184 ἀβρανίδας· κροκωτούς. Λάκωνες ; Hsch.~τ~1507 τρουφωνίδαι· εἶδος κροκωτοῦ ; Poll.~4.117 ὁ δὲ κροκωτὸς, ἱμάτιον. Διόνυσος δὲ αὐτῷ ἐχρῆτο, καὶ μασχαλιστῆρι ἀνθεινῷ, καὶ θύρσῳ ; Suid.~α~3958 ἀρκτευόμεναι γυναῖκες τῇ Ἀρτέμιδι ἑορτὴν ἐτέλουν, κροκωτὸν ἠμφιεσμέναι ; Eur.~\textit{Hec.}~468 ἐν κροκέῳ πέπλῳ ; Pi.~\textit{N.}~1.38 κροκωτὸν σπάργανον ἐγκατέβα ; Plaut.~\textit{Epid.}~231 indusiatam, patagiatam, caltulam aut crocotulam ; Lib.~\textit{Decl.}~33.5 καὶ διαφανὲς χιτώνιον καὶ κροκωτὸν καὶ θέριστρον ; Σ Eur.~\textit{Hec.}~471 ἀνθοβαφέσι καὶ κροκωτοῖς ; Hsch.~κ~4172 κροκόπεπλος· κροκοειδῆ ἱμάτια ἔχουσα}
\lp{253 τὸν κροκωτόν (γαλῇ κροκωτόν)]}{Zen.~2.93 Γαλῇ χιτώνιον: ὁμοία ἡ παροιμία αὕτη τῇ· Οὐ πρέπει γαλῇ κροκωτόν. Ἐπειδὴ γαλῇ κατὰ πρόνοιαν Ἀφροδίτης γυνὴ γενομένη ἐν χιτῶνι κροκωτῷ οὖσα ἐπέδραμε μυί. Μέμνηται ταύτης Στράττις (= Stratt. fr.~71 Kock) ; \textit{idem} Plu.~\textit{Prov.Alex.}~101 (Γαλῇ χιτώνιον, καὶ, Οὐ πρέπει γαλῇ κροκωτός pro Γαλῇ χιτώνιον … κροκωτόν ; μυΐ ; om. Μέμνηται ταύτης Στράττις) ; Diogenian.~3.82 Γαλῇ χιτώνιον: ἐπὶ τῶν τὰ μὴ ἑαυτῶν περιβλήματα ἐνδυόντων ; \textit{idem} Apostol.~5.25 (om. μὴ ; ἐσθιόντων) ; Macar.~6.65 Οὐ πρέπει γαλῇ κροκωτὸς οὔτε πορφύρα: πρὸς τοὺς ὑπὲρ τὴν τύχην κοσμεῖσθαι θέλοντας ; Prov. Coisl.~64 Γαλῆ χιτών: Ἐπὶ τῶν ἀπρεπῶν, ὁμοία τῇ Οὐ πρέπει γαλῇ κροκωτόν. Τὸ δὲ Ἀπέδυ τὸν κροκωτὸν ἡ γαλῆ, ἐπὶ τῶν ἀφεμένων καὶ ἀνακαλυψάντων ἅπερ πρότερον ὑπεκρίνοντο· φασὶ γὰρ οἱ μυθολόγοι, γαλῆν κροκωτῷ κοσμηθεῖσαν καὶ καλλωπιζομένην τῷ ἐνδύματι, καὶ μετὰ σεμνότητος καὶ αἰδοῦς προσισταμένην, εὐθὺς ταῦτα ἀπορρίψαι, καὶ πρὸς μῦν ἐκδραμεῖν ἀναφανέντα ; Aesop.~76 Chambry (3) Γαλῆ ἠράσθη ποτὲ ἀνδρὸς εὐπρεποῦς καὶ τὴν Ἀθηνᾶν ἐδυσώπει ταύτην μεταμεῖψαι εἰς γυναῖκα … ῞Οτι, κἂν πρὸς βραχύ τις ἐν ὑποκρίσει μορφῶται καὶ κρύπτηται, ἡ φύσις τοῦτον διὰ τῶν ἐργων ἐξελέγχει ; Diogenian.~6.98 Ἐπὶ τῶν οὖν ἀναξίως τι βασταζόντων λαμβάνεται. Ὁμοία, Γαλῇ κροκωτόν· καὶ, Πίθηκος ἐν πορφύρᾳ ; Suid.~κ~2460 καὶ παροιμία· Τὸν κροκωτὸν ἡ γαλῆ ; Suid.~ο~929 Οὐ πρέπει γαλῇ κροκωτός: παροιμία ἐστὶν ὁμοία τῇ γαλῇ χιτών ; Erasm.~\textit{Adag.}~1.2.72 Feli crocoton. Γαλῇ κροκωτόν, id est Feli crocoton, subaudiendum das aut addis ; Suid.~γ~34 Γαλῇ χιτών: αὕτη ἐπὶ τῶν ἀδίκων τάσσεται καὶ μηδὲν ἀνιέντων. ὁμοία ἐστὶ τῇ, οὐ πρέπει γαλῇ κροκωτός ; Babr.~fr. (Crusius, append.~3.1.39) Ἀνὴρ γαλῆν γυναῖκα πρὸς δόμους ἄγεν· / παρῆν δὲ Κύπρις εἰς ἑορτὴν τοῦ γάμου. / νύμφη δὲ μῦν βλέψασα συντόνῳ τάχει / δίωκε τοῦτον, μὴ τραπεῖσα τὴν φύσιν}
\lp{253 τὸν κροκωτόν (Clodius)]}{Plu.~\textit{Caes.}~10.1 ὁ Κλώδιος οὔπω γενειῶν καὶ διὰ τοῦτο λήσειν οἰόμενος ἐσθῆτα καὶ σκευὴν ψαλτρίας ἀναλαβὼν ἐχώρει, νέᾳ γυναικὶ τὴν ὄψιν ἐοικώς ; Cic.~\textit{Har.}~44 P. Clodius a crocota, a mitra, a muliebribus soleis purpureisque fasceolis, a strophio, a psalterio, a flagitio, a stupro est factus repente popularis}
\lp{253 πρῶτον ἐνδύου λαβών]}{Theopomp. Com. fr.~52~Kock (Suid.~π~1080) ὑποδοῦ λαβὼν τὰς περιβαρίδας ; Eur.~\textit{Ba.}~836 οὐκ ἂν δυναίμην θῆλυν ἐνδῦναι στολήν ; Ar.~\textit{Thesm.}~212 ἐμοὶ δ᾽ ὅ τι βούλει χρῶ λαβών ; Ar.~\textit{Thesm.}~380 περίθου νυν τόνδε πρῶτον πρὶν λέγειν ; Ar.~\textit{Thesm.}~476 ἐγὼ γὰρ αὐτὴ πρῶτον, ἵνα μἄλλην λέγω ; Suid.~τ~523 Τὴν λεοντῆν ἐνδύου: ἴσον τῷ γενναΐζου}
\lp{254 νὴ τὴν Ἀφροδίτην … ποσθίου]}{\textit{Test.} Suid.~π~2114 Πόσθιον: τὸ αἰδοῖον. … Ἀριστοφάνης· νὴ τὴν Ἀφροδίτην, ἡδύ γ᾽ ὄζει ποσθίου ; Crates fr.~2~K.-A. (=~2~Kock ; Ath.~15.690e) γλυκύτατον δ᾽ ὦζε βασιλείου μύρου}
\lp{254 νὴ τὴν Ἀφροδίτην]}{Ar.~\textit{Lys.}~208 μὰ τὴν Ἀφροδίτην οὔκ, ἐάν γε μὴ λάχῃς ; Ar.~\textit{Eccl.}~158 μὰ Δί᾽ ἀλλ᾽ ἀνὴρ ὢν τὼ θεὼ κατώμοσας ; Ar.~\textit{Eccl.}~189–191 νὴ τὴν Ἀφροδίτην εὖ γε ταυταγὶ λέγεις. / τάλαιν᾽ Ἀφροδίτην ὤμοσας; χαρίεντά γ᾽ ἂν / ἔδρασας, εἰ τοῦτ᾽ εἶπας ἐν τἠκκλησίᾳ ; Gell.~11.6.1 IN veteribus scriptis neque mulieres Romanae per Herculem deiurant neque viri per Castorem ; Ar.~\textit{Lys.}~858 νὴ τὴν Ἀφροδίτην ; Ar.~\textit{Lys.}~939 νὴ τὴν Ἀφροδίτην ἤν τε βούλῃ γ᾽ ἤν τε μή ; Ar.~\textit{Thesm.}~517–518 νὴ τὴν Ἄρτεμιν / ἡμεῖς γε ; Ar.~\textit{Eccl.}~558 νὴ τὴν Ἀφροδίτην μακαρία γ᾽ ἄρ᾽ ἡ πόλις ; Ar.~\textit{Plut.}~1069 μὰ τὴν Ἀφροδίτην οὐκ ἐμοῦ γ᾽ ὦ βδελυρὲ σύ ; Phoenicid.~fr.~4.1~K.-A. (=~4~Kock, ex Stob.) μὰ τὴν Ἀφροδίτην οὐκ ἂν ὑπομείναιμ᾽ ἔτι, / Πυθιάς, ἑταιρεῖν ; Ar.~\textit{Thesm.}~225 οὐ γὰρ μὰ τὴν Δήμητρά γ᾽ ἐνταυθοῖ μενῶ ; Ar.~\textit{Thesm.}~383 φιλοτιμίᾳ μὲν οὐδεμιᾷ μὰ τὼ θεὼ ; Ar.~\textit{Thesm.}~533 οὔ τοι μὰ τὴν Ἄγλαυρον ὦ γυναῖκες εὖ φρονεῖτε ; Ar.~fr.~325~K.-A. (Phot.~α~3404) Ἀφρόδιτος· ὁ Ἑρμαφρόδιτος. παραπλήσιοι δὲ τούτῳ καὶ ἄλλοι δαίμονες … ὧν καὶ Ἀριστοφάνης μέμνηται Ἥρωσιν}
\lp{254 ἡδύ γ᾽ ὄζει]}{Ar.~\textit{Vesp.}~1058–1059 ὑμῖν δι᾽ ἔτους τῶν ἱματίων / ὀζήσει δεξιότητος ; Ar.~\textit{Pax}~525–526 οἷον δὲ πνεῖς, ὡς ἡδὺ κατὰ τῆς καρδίας, / γλυκύτατον ὥσπερ ἀστρατείας καὶ μύρου ; Ar.~\textit{Pax}~529–532 τοῦ μὲν γὰρ ὄζει κρομμυοξυρεγμίας, / ταύτης δ᾽ ὀπώρας, ὑποδοχῆς, Διονυσίων, / αὐλῶν, τραγῳδῶν, Σοφοκλέους μελῶν, κιχλῶν, / ἐπυλλίων Εὐριπίδου— ; Ar.~\textit{Lys.}~206 καὶ μὰν ποτόδδει γ᾽ ἁδὺ ναὶ τὸν Κάστορα ; Ar.~\textit{Lys.}~662–663 ὡς τὸν ἄνδρα δεῖ / ἀνδρὸς ὄζειν εὐθύς ; Ar.~\textit{Lys.}~942–943 οὐχ ἡδὺ τὸ μύρον μὰ τὸν Ἀπόλλω τουτογί, / εἰ μὴ διατριπτικόν γε κοὐκ ὄζον γάμων ; Ar.~\textit{Ran.}~338 ὦ πότνια πολυτίμητε Δήμητρος κόρη, / ὡς ἡδύ μοι προσέπνευσε χοιρείων κρεῶν ; Ar.~\textit{Eccl.}~524 εἰ τῆς κεφαλῆς ὄζω μύρου ; Ar.~\textit{Plut.}~1020 ὄζειν τε τῆς χρόας ἔφασκεν ἡδύ μου ; Plaut.~\textit{Most.}~273 quia ecastor mulier recte olet, ubi nihil olet ; \textit{idem} Tosi n°~305 (om. quia ecastor ; Mulier) ; X.~\textit{Smp.}~2.3–4 ὥσπερ γάρ τοι ἐσθὴς ἄλλη μὲν γυναικί, ἄλλη δὲ ἀνδρὶ καλή, οὕτω καὶ ὀσμὴ ἄλλη μὲν ἀνδρί, ἄλλη δὲ γυναικὶ πρέπει ; Thphr.~\textit{Char.}~4.1 καὶ τὸ μύρον φάσκειν οὐδὲν τοῦ θύμου ἥδιον ὄζειν, καὶ μείζω τοῦ ποδὸς τὰ ὑποδήματα φορεῖν ; Plu.~\textit{Alex.}~4.2 ὅτι δὲ τοῦ χρωτὸς ἥδιστον ἀπέπνει καὶ τὸ στόμα κατεῖχεν εὐωδία καὶ τὴν σάρκα πᾶσαν, ὥστε πληροῦσθαι τοὺς χιτωνίσκους, ἀνέγνωμεν ἐν ὑπομνήμασιν Ἀριστοξενείοις ; Hsch.~ο~115 ὄδωδεν· ὄζει, πονηρὰν καὶ βαρυτάτην ἀποφορὰν πνεῖ, ἢ εὐῶδές ἐστιν ; Ar.~\textit{Ach.}~192–193 ὄζουσι χαὖται πρέσβεων ἐς τὰς πόλεις / ὀξύτατον ; Ar.~\textit{Ach.}~196 αὗται μὲν ὄζουσ᾽ ἀμβροσίας καὶ νέκταρος ; Ar.~\textit{Ach.}~852–853 ὄζων κακὸν τῶν μασχαλῶν / πατρὸς Τραγασαίου ; Ar.~\textit{Eq.}~892 οὐκ ἐς κόρακας ἀποφθερεῖ βύρσης κάκιστον ὄζων ; Ar.~\textit{Eq.}~1332 οὐ χοιρινῶν ὄζων ἀλλὰ σπονδῶν, σμύρνῃ κατάλειπτος ; Ar.~\textit{Nub.}~50 ὄζων τρυγὸς τρασιᾶς ἐρίων περιουσίας ; Ar.~\textit{Vesp.}~38 ὄζει κάκιστον τοὐνύπνιον βύρσης σαπρᾶς ; Pherecr.~fr.~67~Kock (Ath.~4.159e) ἢν γὰρ τράγῃ τις, τοῦ στόματος ὄζει κακόν ; Hermipp.~fr.~82.8~Kock (Ath.~1.29e) ὄζει ἴων, ὄζει δὲ ῥόδων, ὄζει δ᾽ ὑακίνθου ; Theoc.~1.149 ἠνίδε τοι τὸ δέπας· θᾶσαι φίλος, ὡς καλὸν ὄσδει ; Mart.~2.12.4 Postume, non bene olet qui bene semper olet ; Plu.~\textit{Dem.}~8.3–4 Πυθέας ἐπισκώπτων ἐλλυχνίων ἔφησεν ὄζειν αὐτοῦ τὰ ἐνθυμήματα ; \textit{idem} Tosi n°~2095 (Ἐλλυχνίων ὄζειν τὰ ἐνθυμήματα) ; Aesch.~\textit{Ag.}~1310 τόδ᾽ ὄζει θυμάτων ἐφεστίων ; \textit{Od.}~9.210 ὀδμὴ δ᾽ ἡδεῖα ἀπὸ κρητῆρος ὀδώδει ; Hsch.~κ~2442 κεῶεν ὄζει· εὐωδεῖ}
\lp{254 ποσθίου]}{\textit{Test.} Hsch.~π~3096 πόσθ[ο]ιον· αἰδοῖον ἀνδρός ; Phot.~π~1113 πόσθιον· αἰδοῖον· ἢ τὴν Ἀφροδίτην ; Ar.~\textit{Thesm.}~515 τά τ᾽ ἄλλ᾽ ἁπαξάπαντα καὶ τὸ πόσθιον ; Σ Ar.~\textit{Thesm.}~254 δέον εἰπεῖν μύρου εἶπε ποσθίου. πόσθιον δέ ἐστι τὸ αἰδοῖον τοῦ ἀνδρός. διαβάλλει δὲ αὐτὸν ὡς μετ᾽ αὐτοῦ ἑταιρίζοντα ; Phot.~π~1114 πόσθων· κυρίως λέγεται παιδάριον· οὕτως γὰρ ὑποκοριζόμενοι ἔλεγον ἀπὸ τοῦ αἰδοίου· πόσθιον γὰρ τοῦτο καλοῦσιν· δηλοῖ δὲ καὶ τὸν παιδαριώδη καὶ ἄφρονα κατὰ μεταφοράν ; \textit{Test.} Suid.~π~2113 Πόσθη: τὸ αἰδοῖον τοῦ ἀνδρός ; Ar.~\textit{Nub.}~1014 πυγὴν μεγάλην, πόσθην μικράν ; Ar.~\textit{Pax}~1300 εἰπέ μοι ὦ πόσθων ; Ar.~\textit{Thesm.}~1188 εἶεν· καλὴ τὸ σκῆμα περὶ τὸ πόστιον ; Hsch.~π~3095 πόσθη· τοῦ αἰδοίου τὸ δέρμα ; Hsch.~π~3097 πόσθων· πόσθην τὸ ἀνδρεῖον αἰδοῖον λέγουσι ; Hsch.~σ~1274 σμορδονεύς· ὑποκοριστικῶς ἀπὸ τῶν μορίων· ὡς πόσθωνες ; Ar.~\textit{Thesm.}~62 τουτὶ τὸ πέος χοανεῦσαι ; Hsch.~α~2621 ἀκροπόσθια· τὰ ἄκρα τοῦ αἰδοίου ; Synag.~π~829 πόσθων· ὁ παιδαριώδης καὶ ἀνάγωγος. Μένανδρος ἐν Ὑμνίδι}
\lp{255 σύζωσον]}{Ar.~\textit{Lys.}~536 κᾆτα ξαίνειν ξυζωσάμενος ; Ar.~\textit{Thesm.}~656 ξυζωσαμένας εὖ κἀνδρείως τῶν θ᾽ ἱματίων ἀποδύσας ; Eur.~\textit{Ba.}~935–936 ζῶναί τέ σοι χαλῶσι κοὐχ ἑξῆς πέπλων / στολίδες ὑπὸ σφυροῖσι τείνουσιν σέθεν ; Hes.~\textit{Op.}~72–76 ζῶσε δὲ καὶ κόσμησε θεὰ γλαυκῶπις Ἀθήνη· / … / [πάντα δέ οἱ χροῒ κόσμον ἐφήρμοσε Παλλὰς Ἀθήνη.] ; Hes.~\textit{Th.}~573–575 ζῶσε δὲ καὶ κόσμησε θεὰ γλαυκῶπις Ἀθήνη / ἀργυφέῃ ἐσθῆτι· κατὰ κρῆθεν δὲ καλύπτρην / δαιδαλέην χείρεσσι κατέσχεθε, θαῦμα ἰδέσθαι ; Ar.~\textit{Lys.}~72 μόλις γὰρ ηὗρον ἐν σκότῳ τὸ ζώνιον ; Aesch.~\textit{Supp.}~457–462 ἔχω στρόφους ζώνας τε, συλλαβὰς πέπλων. / … / τί σοι περαίνει μηχανὴ συζωμάτων ; Hsch.~ζ~236 ζώνιον· γυναικεῖον ζῶμα ; Hsch.~ζ~258 ζώσατο· τὴν ζώνην περὶ τὴν ὀσφὺν ἔσφιγξε ; Soph.~\textit{El.}~452 καὶ ζῶμα τοὐμὸν οὐ χλιδαῖς ἠσκημένον}
\lp{255 ἀνύσας]}{Ar.~\textit{Nub.}~635 ἀνύσας τι κατάθου καὶ πρόσεχε τὸν νοῦν. — ἰδού ; Ar.~\textit{Vesp.}~1158 ἄγε νυν ὑπολύου τὰς καταράτους ἐμβάδας, / τασδὶ δ᾽ ἀνύσας † ὑπόδυθι † τὰς Λακωνικάς ; Ar.~\textit{Lys.}~920 ἰδοὺ κατάκεισ᾽ ἀνύσας τι, κἀγὼ ᾽κδύομαι ; Hsch.~α~5574 ἀνύσας· συνελών ; Ar.~\textit{Vesp.}~30 λέγε νυν ἀνύσας τι τὴν τρόπιν τοῦ πράγματος ; Ar.~\textit{Eccl.}~1058 ἕπου μαλακίων δεῦρ᾽ ἀνύσας καὶ μὴ λάλει ; Σ Ar.~\textit{Nub.}~635 ἀνύσας: σπεύσας ; Σ Ar.~\textit{Vesp.}~30 λέγε νῦν ἀνύσας: Τὸ ἀνύσας Ἀττικοὶ δασύνουσι ; Ar.~\textit{Ach.}~570–571 βοηθησάτω / τις ἀνύσας ; Ar.~\textit{Nub.}~181 ἄνοιγ᾽ ἄνοιγ᾽ ἀνύσας τὸ φροντιστήριον}
\lp{255 αἶρε νῦν στρόφιον. — ἰδού]}{Ar.~\textit{Pax}~1–2 αἶρ᾽ αἶρε μᾶζαν ὡς τάχιστα κανθάρῳ. — ἰδού}
\lp{255 αἶρε]}{Ar.~\textit{Pax}~1227 ἀλλ᾽ αἶρέ μοι τοῦτόν γε τῆς ἰσωνίας ; Ar.~\textit{Av.}~759 αἶρε πλῆκτρον, εἰ μαχεῖ ; Soph.~\textit{Aj.}~545 αἶρ᾽ αὐτόν, αἶρε δεῦρο ; \textit{Il.}~6.264 μή μοι οἶνον ἄειρε μελίφρονα πότνια μῆτερ}
\lp{255 στρόφιον]}{Ar.~\textit{Lys.}~931 τὸ στρόφιον ἤδη λύομαι ; Ar.~\textit{Thesm.}~139 τί λήκυθος καὶ στρόφιον; ὡς οὐ ξύμφορον ; Ar.~\textit{Thesm.}~638 χάλα ταχέως τὸ στρόφιον ὦναίσχυντε σύ ; Ar.~fr.~332~K.-A. (=~320~Kock ; Poll.~7.95 ; Clem. Al. Paed.~2.12) μίτρας, ἀναδήματα, νίτρον, κίσηριν, στρόφιον, ὀπισθοσφενδόνην, κάλυμμα, φῦκος, περιδέραια, ὑπογράμματα, τρυφοκαλάσιριν, ⟨ἐλλέβορον,⟩ κεκρύφαλον, ζῶμα, ἀμπέχονον, τρύφημα, παρυφές, ξυστίδα, κιθῶνα, βάραθρον, ἔγκυκλον ; Ar.~fr.~664~K.-A. ἀλλὰ τὸ στρόφιον λυθὲν / τὰ κάρυά μου ἐξέπιπτεν ; \textit{idem} Poll.~7.67 (praem. τὸ μέντοι γυναικεῖον ζώνιον οὕτω τ᾽ ἐκάλουν καὶ στρόφιον ὠνόμαζον, ὡς Ἀριστοφάνης) ; Catull.~64.65 non tereti strophio lactentis vincta papillas ; Prop.~4.9.49 mollis et hirsutum cepit mihi fascia pectus ; Ov.~\textit{Fast.}~2.319–324 dat tenuis tunicas Gaetulo murice tinctas, / dat teretem zonam, qua modo cincta fuit. / ventre minor zona est; tunicarum vincla relaxat, / … / scindebant magni vincula parva pedes ; Hsch.~σ~2044 στρόφιον· ἡ στρογγύλη ζώνη ; Hsch.~τ~802 Τηνία· ἔχιδνα. ἄλλοι δὲ τὴν στηθοδέσμην ; Ar.~\textit{Thesm.}~251 καὶ στρόφιον· οὐ γὰρ ταῦτά γ᾽ ὡς οὐκ ἔστ᾽ ἐρεῖς ; Hsch.~α~6298 ἀπόδεσμος· κόσμιόν τι γυναικεῖον, περικεφάλαιον ; Hsch.~κ~2618 κίδαρις· … ἢ στρόφιον, ὃ οἱ ἱερεῖς φοροῦσιν ; Synag.~σ~261 στρόφιον· τὸ στρογγύλον ζωνάριον ; \textit{idem} Phot.~σ~639 (add. Ἀριστοφάνης· „τὸ στρόφιον ἤδη λύομαι“), Suid.~σ~1221 (add. Ἀριστοφάνης· τὸ στρόφιον ἤδη λυόμενον) ; Suid.~σ~1222 Στρόφιον: ὃ οἱ ἱερεῖς φοροῦσι}
\lp{255 ἰδού]}{Ar.~\textit{Ach.}~434 μεταξὺ τῶν Ἰνοῦς. — ἰδοὺ ταυτὶ λαβέ ; Eur.~\textit{Ba.}~934 ἰδού, σὺ κόσμει· σοὶ γὰρ ἀνακείμεσθα δή ; Ar.~\textit{Thesm.}~25 βάδιζε δευρὶ καὶ πρόσεχε τὸν νοῦν. — ἰδού ; Hsch.~ι~170 ἴδε· ἰδού}
\lp{256 ἴθι νυν]}{Ar.~\textit{Nub.}~497 ἴθι νυν κατάθου θοἰμάτιον ; Ar.~\textit{Ran.}~494–496 ἴθι νυν ἐπειδὴ ληματίας κἀνδρεῖος εἶ, / σὺ μὲν γενοῦ ᾽γὼ τὸ ῥόπαλον τουτὶ λαβὼν / καὶ τὴν λεοντῆν ; Ar.~\textit{Nub.}~237 ἴθι νυν κατάβηθ᾽ ὦ Σωκρατίδιον ὡς ἐμέ ; Ar.~\textit{Vesp.}~843 ἴθι νυν ἄγ᾽ αὐτὼ δεῦρο ; Ar.~\textit{Lys.}~861 ἴθι νυν κάλεσον αὐτήν ; Σ Ar.~\textit{Nub.}~497 ἴθι νῦν κατάθου: Ἀντὶ τοῦ ἄγε ; Ar.~\textit{Eq.}~105 ἴθι νυν ἄκρατον ἐγκάναξόν μοι πολὺν}
\lp{256 κατάστειλον]}{Ar.~\textit{Av.}~1567–1568 οὗτος τί δρᾷς; ἐπ᾽ ἀριστέρ᾽ οὕτως ἀμπέχει; / οὐ μεταβαλεῖς θοἰμάτιον ὧδ᾽ ἐπιδέξια ; Eur.~\textit{Ba.}~821–822 στεῖλαί νυν ἀμφὶ χρωτὶ βυσσίνους πέπλους. / τί δὴ τόδ᾽; ἐς γυναῖκας ἐξ ἀνδρὸς τελῶ ; Eur.~\textit{Ba.}~827 ἐγὼ στελῶ σε δωμάτων ἔσω μολών ; Eur.~\textit{Ba.}~932–933 ἀλλ᾽ αὐτὸν ἡμεῖς, οἷς σε θεραπεύειν μέλει, / πάλιν καταστελοῦμεν ; Eur.~\textit{Ba.}~938 κἀμοὶ δοκοῦσι παρά γε δεξιὸν πόδα· / τἀνθένδε δ᾽ ὀρθῶς παρὰ τένοντ᾽ ἔχει πέπλος ; Stat.~\textit{Ach.}~1.325–329 iniecitque sinus; tum colla rigentia mollit / submittitque graves umeros et fortia laxat / bracchia et impexos certo domat ordine crines / ac sua dilecta cervice monilia transfert ; Hsch.~κ~1433 καταστέλλει· περικαλύπτει ; Ar.~\textit{Vesp.}~1132 τηνδὶ δὲ χλαῖναν ἀναβαλοῦ τριβωνικῶς ; Eur.~\textit{Med.}~1161–1166 λαμπρῷ κατόπτρῳ σχηματίζεται κόμην, / … / δώροις ὑπερχαίρουσα, πολλὰ πολλάκις / τένοντ᾽ ἐς ὀρθὸν ὄμμασι σκοπουμένη ; Plu.~\textit{Mor.}~69c καὶ γὰρ αἱ τίτθαι τοῖς παιδίοις πεσοῦσιν οὐ λοιδορησόμεναι προστρέχουσιν, ἀλλ᾽ ἤγειραν καὶ περιέπλυναν καὶ κατέστειλαν ; Plu.~\textit{Per.}~5.1 πρᾳότης πορείας καὶ καταστολὴ περιβολῆς ; Hsch.~κ~1441 καταστολήν· περιβολήν ; Apollod. Com. fr.~19.2~K.-A. χρόνον γὰρ εἰς τὰ πράγματ᾽ ἂν λάβηις βραχύν, / ἅπαντα λήξει καὶ κατασταλήσεται ; NT~1 Ep. Tim.~2.9 γυναῖκας ἐν καταστολῇ κοσμίῳ μετὰ αἰδοῦς καὶ σωφροσύνης κοσμεῖν ἑαυτάς, μὴ ἐν πλέγμασιν καὶ χρυσίῳ}
\lp{256 τὰ περὶ τὼ σκέλει]}{Ar.~\textit{Thesm.}~246 αἰθὸς γεγένημαι πάντα τὰ περὶ τὴν τράμιν ; Hsch.~σ~896 σκέλεαι· τὰ τῶν σκελῶν σκεπάσματα ; Harp. s.v.~ἀμφιδέαι ἈΜΦΙΔΕΑΙ εἰσὶ περισκελίδες τινές· Ἀριστοφάνης Θεσμοφοριαζούσαις ; \textit{idem} Suid.~α~1715 (Ἀμφιδέας ; om. Ἀριστοφάνης Θεσμοφοριαζούσαις) ; Suid.~π~1283 Περισκέλια: ἃ ἐφόρουν οἱ Ῥωμαῖοι περὶ τὰ σκέλη. ἧσαν δὲ κροκοβαφῆ ; Ar.~\textit{Vesp.}~438 ὦ Κέκροψ ἥρως ἄναξ τὰ πρὸς ποδῶν Δρακοντίδη ; Ar.~\textit{Av.}~66 καὶ μὴν ἐροῦ τὰ πρὸς ποδῶν ; Ar.~\textit{Av.}~1254 πρώτης ἀνατείνας τὼ σκέλει διαμηριῶ ; Ar.~\textit{Thesm.}~24 ἔτι προσμάθοιμι χωλὸς εἶναι τὼ σκέλει ; Ar.~\textit{Eccl.}~265 εἰθισμέναι γάρ ἐσμεν αἴρειν τὼ σκέλει ; Plu.~\textit{Mor.}~142c τῶν δὲ πλείστων γυναικῶν ἂν ὑποδήματα διάχρυσα περιέλῃς καὶ ψέλλια καὶ περισκελίδας καὶ πορφύραν καὶ μαργαρίτας, ἔνδον μένουσιν ; Hsch.~π~1865 περισκελῆ· β⟨ρ⟩ακ[κ]ία, φιμινάλια ; \textit{idem} Suid.~π~1282 (βρακία, φημινάλια) ; Hsch.~σ~910 σκέλος· τάττεται ἐπὶ τοῦ φορτικοῦ ; Ar.~\textit{Pax}~820 ἔγωγέ τοι πεπόνηκα κομιδῇ τὼ σκέλει ; Longus~1.5.3 μίτρα διάχρυσος, ὑποδήματα ἐπίχρυσα, περισκελίδες χρυσαῖ ; Hsch.~σ~909 σκελοῖν· ποδῶν}
\lp{257–258 ἡδὶ μὲν οὖν κεφαλὴ περίθετος, ἣν ἐγὼ νύκτωρ φορῶ]}{\textit{Test.} Poll.~2.35 ἔλεγον δὲ καὶ κεφαλὴν περίθετον, ὡς Ἀριστοφάνης, καὶ ὡς ἂν εἴποιεν οἱ νῦν περιθέτην, καὶ πάλιν τὴν περίθετον κόμην ; Phot.~π~690 περίθετος κεφαλή· ἡ περιθετή. Ἀριστοφάνης· „κεφαλὴ περίθετος, ἣν ἐγὼ νύκτωρ φορῶ“ ; \textit{idem} Suid.~π~1165 (Περίθετος· πρόσωπον ; ἥδιον μὲν οὖν κεφαλὴ … ; add. διαβάλλεται δὲ Ἀγάθων ἐπὶ μαλακίᾳ, ἵνα λανθάνῃ νυκτὸς ὡς γυνὴ καὶ πάσχῃ), Eust.~1280.63 (ὃν περίθετον κεφαλὴν ἐκάλουν. Ἀριστοφάνης Θεσμοφόροις)}
\lp{257 ἡδὶ μὲν οὖν]}{Ar.~\textit{Av.}~292 ἦ ᾽πὶ τὸν δίαυλον ἦλθον; — ὥσπερ οἱ Κᾶρες μὲν οὖν ; Ar.~\textit{Av.}~341 ἵνα μὲν οὖν κλάοιμι μεγάλα}
\lp{257 κεκρυφάλου]}{Ar.~\textit{Thesm.}~138 λαλεῖ κροκωτῷ; τί δὲ λύρα κεκρυφάλῳ ; Antiph.~fr.~117~Kock (Ath.~15.681c) οὐκ ἐφύσων οἱ Λάκωνες ὡς ἀπόρθητοί ποτε; νῦν δ᾽ ὁμηρεύουσ᾽ ἔχοντες πορφυροῦς κεκρυφάλους ; \textit{Il.}~22.469 τῆλε δ᾽ ἀπὸ κρατὸς βάλε δέσματα σιγαλόεντα, / ἄμπυκα κεκρύφαλόν τε ἰδὲ πλεκτὴν ἀναδέσμην ; Juv.~2.96 reticulumque comis ; D.H.~7.9.4 κομᾶν τε γὰρ τοὺς ἄρρενας ὥσπερ τὰς παρθένους ἐκέλευσεν ἐξανθιζομένους καὶ βοστρυχιζομένους καὶ κεκρυφάλοις τὰς πλοκαμίδας ἀναδοῦντας ἐνδύεσθαί τε ποικίλους καὶ ποδήρεις χιτωνίσκους ; Luc.~\textit{Merc.Cond.}~33 κίναιδόν τινα τῶν πεπιττωμένων τὰ σκέλη καὶ τὸν πώγωνα περιεξυρημένων … καὶ εἴ γε μὴ πολλὰ δεηθῆναι αὐτοῦ, καὶ τὸν κεκρύφαλον ἔχοντα ἐπὶ τῇ κεφαλῇ ἂν συγκαθίζεσθαι ; Clem.Al.~\textit{Paed.}~3.3 μικροῦ δεῖν ἀναδούμενοι κεκρυφάλοις τὰς κόμας ; Hsch.~α~7291 ἄρκυν· μεσάγκυλον. καὶ γυναικεῖον κεκρύφαλον ; Hsch.~κ~2110 κεκρύφαλος· εἶδος κοιλίας τῶν μηρυκαζομένων ζώων. σαβακάθιον. δεσμότριχον ; \textit{AP}~16.78 βάσκανος ἡ γραφίς ἐσσι, καὶ εἰσορόωσι μεγαίρεις / χρύσεα κεκρυφάλοις βόστρυχα κρυψαμένη ; Ael.~\textit{VH}~7.9 οὐ κροκωτοῦ οὐ ταραντινιδίου οὐκ ἀναβολῆς οὐκ ἐγκύκλου οὐ κεκρυφάλου οὐ καλύπτρας οὐ βαπτῶν χιτωνίσκων ; Clem.Al.~\textit{Paed.}~3.4 αἳ μὲν τὰ κάτοπτρα, αἳ δὲ τοὺς κεκρυφάλους, ἄλλαι τοὺς κτένας ; IG II²~1524,~5 κεκρύφαλον λευκὸν πεζίδα ἔ[χοντα] ; Hsch.~α~3817 ἄμπυκα, κεκρύφαλον· κόσμον γυναικῶν ; Hsch.~κ~768 κάραννος· κεκρύφαλος. κρήδεμνον ; Hsch.~κ~3703 κορυμβάσι· περιδρόμοις, δι᾽ ὧν συσπᾶται γύργαθος καὶ κεκρύφαλος καὶ δεσμοί ; Hsch.~κ~3730 κορυφιστής· κόσμου γυναικείου τὸ περὶ τὴν κεφαλὴν χρυσίον, καὶ κεκρυφάλου τὸ μέσον ῥάμμα ; Poll.~7.66 φάριον δὲ τὸν ἐρεοῦν κεκρύφαλον ὠνόμαζον ; Poll.~7.179 κεκρυφαλοπλόκος, ὡς Κριτίας ἔφη, κεκρύφαλος κεκρυφάλιον ; Poll.~10.192 ὅταν δὲ Δημοσθένης εἴπῃ σακχυφάντας, τοὺς πλέκοντας ταῖς γυναιξὶ τοὺς κεκρυφάλους ἀκούουσιν ; Synag.~κ~260 κεκρύφαλον· κεκρύφαντον, σαβακάθιον, σουδάριον ; \textit{idem} Phot.~κ~541, Suid.~κ~1273 (Κεκρύφαλον· κεκρύφατον ; add. ἐν Ἐπιγράμμασι· κεκρύφαλοι σφίγγουσι τεὴν τρίχα) ; Eup.~fr.~170~Kock (Poll.~7.192) κεκρύφαλοί τε καὶ τύλη ; Antiph.~fr.~187~K.-A. (=~189~Kock ; Phot.~σ~560) στίμμιν, κάτοπτρα, κροβύλους, κεκρυφάλους ; Plu.~\textit{Sept.\,sap.\,conv.}~154b παίζουσαν καὶ διαπλέκουσαν ὥσπερ ἕτεραι ζώνια καὶ κεκρυφάλους ; Alciphr.~1.6 Schepers κεκρυφάλους Μιλησίους καὶ Σικελικὸν ἱμάτιον καὶ ἐπ᾽ αὐτῷ χρυσίον εἰσπέμπεις ; Clem.Al.~\textit{Paed.}~3.2 τὸν κεκρύφαλον λέγω, τὴν βαφήν, τὴν ἐσθῆτα, τὸ χρυσίον, τὸ φῦκος ; Hsch.~κ~2111 κεκρυφάλους· τοὺς τῶν ἵππων κορυφαστῆρας, καὶ κροκυφάντους ; Hsch.~π~2411 καὶ πλαγγόνες· κεκρύφαλοι}
\lp{257 μίτρας]}{Ar.~\textit{Thesm.}~163 ἐμιτροφόρουν τε ; Ar.~\textit{Thesm.}~941 ἵνα μὴ ᾽ν κροκωτοῖς καὶ μίτραις γέρων ἀνὴρ ; Pherecr.~fr.~106~K.-A. (Phot.~ο~743) ὄχθοιβος· λῶμα. Φερεκράτης Λήροις· „μίτραν ἁλουργῆ, στρόφιον, ὄχθοιβον, κτένα“ ; Eur.~\textit{Ba.}~833 πέπλοι ποδήρεις· ἐπὶ κάρᾳ δ᾽ ἔσται μίτρα ; Eur.~\textit{Ba.}~929 ἀλλ᾽ ἐξ ἕδρας σοι πλόκαμος ἐξέστηχ᾽ ὅδε, / οὐχ ὡς ἐγώ νιν ὑπὸ μίτρᾳ καθήρμοσα ; Verg.~\textit{Aen.}~4.215–217 Et nunc ille Paris cum semiviro comitatu, / Maeonia mentum mitra crinemque madentem / subnexus ; Prop.~4.2.32 cinge caput mitra, speciem furabor Iacchi ; Ov.~\textit{Met.}~14.654–656 Ille etiam picta redimitus tempora mitra, / innitens baculo, positis per tempora canis, / adsimulavit anum cultosque intravit in hortos ; Juv.~2.84–85 paulatim qui longa domi redimicula sumunt / frontibus et toto posuere monilia collo ; Plu.~\textit{Mor.}~304c διὰ τί παρὰ Κῴοις ὁ τοῦ Ἡρακλέους ἱερεὺς ἐν Ἀντιμαχείᾳ γυναικείαν ἐνδεδυμένος ἐσθῆτα καὶ τὴν κεφαλὴν ἀναδούμενος μίτρᾳ κατάρχεται τῆς θυσίας ; Plu.~\textit{Sol.}~8.5 τῶν δὲ νεωτέρων τοὺς μηδέπω γενειῶντας ἐνδύμασι καὶ μίτραις καὶ ὑποδήμασι τοῖς ἐκείνων σκευασαμένους καὶ λαβόντας ἐγχειρίδια κρυπτὰ παίζειν καὶ χορεύειν προσέταξε ; Polyaen.~\textit{Exc.}~28 αὐτὸς δὲ νέους στολίσας γυναικείαν στολὴν μετὰ ἐγχειριδίων προῆγεν ὡς πρὸς τὴν ἑορτὴν τῆς Δήμητρος ; Suid.~ε~989 ἀντὶ τοῦ μίτραις τουτέστι ζώναις διαζωσάμενοι ἁβρῶς περιεπάτουν ; Eur.~\textit{Ba.}~1115 καὶ προσπίτνει νιν· ὃ δὲ μίτραν κόμης ἄπο ; Eur.~\textit{Hec.}~924 ἐγὼ δὲ πλόκαμον ἀναδέτοις / μίτραισιν ἐρρυθμιζόμαν ; Catull.~64.63 non flavo retinens subtilem vertice mitram ; Phoenix Coloph. fr.~1 (Ath.~12.530e–f) σποδὸς δὲ πολλὴ χὠ μιτρηφόρος κεῖμαι ; Hdt.~1.195.1 κομῶντες δὲ τὰς κεφαλὰς μίτρῃσι ἀναδέονται ; Luc.~\textit{Bacch.}~2 μίτρᾳ τὴν κόμην ἀναδεδεμένον, ἐν πορφυρίδι καὶ χρυσῇ ἐμβάδι ; Tatian.~\textit{Ad\,Gr.}~10 τὸν μὲν καρήξανθον Μενέλαον καταλιποῦσα, τῷ δὲ μιτρηφόρῳ καὶ πολυχρύσῳ Πάριδι κατακολουθοῦσα ; Hsch.~α~198 ἁβρομίτρας· λαμπροζώνους ; Hsch.~α~4262 ἀναδέσμη· μίτρα, διάδημα. οἱ δὲ εἶδος κόσμου ἐπὶ κεφαλῆς ; Hsch.~κ~3912 κράδεμνον· κάλυμμα. στεφάνη. μίτρα ; Hsch.~μ~1475 μίτρα· ἡ χαλκῆ λεπίς. διάδημα. ζώνη. θώραξ. ταινία ; Hsch.~μ~1476 μίτρα χρυσᾶ· ὡς λεπὶς χρυσᾶ, ἣ τῷ μετώπῳ φορεῖται ; Suid.~α~1850 Ἀναδέσμην: ὅμοιόν τι διαδήματι· ἢ μίτραν κεφαλῆς ; Suid.~μ~1136 Μίτρα: διάδημα, ἢ ζώνη ; Diogenes \textit{Trag.,\,Semele} (Ath.~14.636a) καίτοι κλύω μὲν Ἀσιάδος μιτρηφόρους Κυβέλας γυναῖκας, παῖδας ὀλβίων Φρυγῶν ; Mosch.~\textit{Eur.}~73 οὐδ᾽ ἄρα παρθενίην μίτρην ἄχραντον ἔρυσθαι ; Lucr.~4.1129 et bene parta patrum fiunt anademata, mitrae ; Hdt.~7.62.2 ἀντὶ δὲ τῶν πίλων μιτρηφόροι ἦσαν ; Arr.~\textit{Ind.}~7.9 μιτρηφορέειν τε ἀναδεῖξαι, καὶ μύρων ἀλοιφὰς ἐκδιδάξαι ; Hsch.~ζ~100 ζείρη· μίτρα. ταινία. διάδημα}
\lp{258 κεφαλὴ περίθετος]}{Ar.~\textit{Ach.}~439 κἀκεῖνά μοι δὸς τἀκόλουθα τῶν ῥακῶν, / τὸ πιλίδιον περὶ τὴν κεφαλὴν τὸ Μύσιον ; Ar.~\textit{Lys.}~532–533 παρ᾽ ἐμοῦ τουτὶ τὸ κάλυμμα λαβὼν / ἔχε καὶ περίθου περὶ τὴν κεφαλήν ; Ar.~\textit{Eccl.}~99–100 τὸν πώγωνά τε / ὅταν καθῶμεν ὃν περιδησόμεσθ᾽ ἐκεῖ ; Ar.~fr.~181~Kock (Poll.~10.170) καὶ πηνήκην δὲ καὶ προκόμιον, καὶ ὡς Ἀριστοφάνης ἐν Γηρυτάδῃ περίθεσιν· καὶ περιθέτην δέ, ὡς ἐν Ἄμφιδος Ἀλκμαίωνι καὶ Μενάνδρου Ὀλυνθίᾳ (= Amphis fr.~2 Kock ; Men. fr.~359 Kock) ; Ar.~fr.~321~Kock (Σ \textit{Plut.}~159) ὅσα περίεργ᾽ αὐταῖσι τῶν φορημάτων / ὅσοις τε περιπέττουσιν αὐτὰς προσθέτοις ; Ar.~fr.~332.1–2~K.-A. (=~320~Kock ; Poll.~7.95) ξυρόν, κάτοπτρον, ψαλίδα, κηρωτήν, λίτρον, / προκόμιον, ὀχθοίβους, μίτρας, ἀναδήματα ; Ar.~fr.~898~Kock (Gramm. Darmst.) ἐγὼ διὰ ταῦτα, μὴ γέλων ὄφλων λάθω, / περὶ τὴν κεφαλὴν ἐξῆμμαι πηνίκην τινά ; Cratin.~fr.~282~Kock (EM~269,~31) τοῦτον μὲν οὖν καλῶς διεπηνίκισας λόγον ; Aristomen.~fr.~5~Kock (Σ Thesm.~258) παντευχίαν δὲ τοῦ θεοῦ ταύτην λαβεῖν / καὶ περίθετον πρόσωπον, ὃ λαβὼν ἕσταθι ; Eur.~\textit{Ba.}~831 κόμην μὲν ἐπὶ σῷ κρατὶ ταναὸν ἐκτενῶ ; Arist.~\textit{Oec.}~1348a τούς τε Λυκίους ὁρῶν ἀγαπῶντας τὸ τρίχωμα φορεῖν, ἔφησε γράμματα ἥκειν παρὰ βασιλέως, κόμας ἀποστεῖλαι εἰς προκόμια ; Luc.~\textit{DMeretr.}~12.5 νῦν δὲ καὶ τὴν πηνήκην ἐπέθετο ; Ath.~12.523a ὥστε πρῶτοι τὸ πρόσωπον ἐντριψάμενοι καὶ προκόμια περιθετά λαβόντες στολὰς μὲν ἀνθινὰς ; Clem.Al.~\textit{Paed.}~3.11 ἀλλοτρίων δὲ αἱ προσθέσεις τριχῶν τέλεον ἔκβλητοι … οὐ τὴν γυναῖκα τὴν κεκοσμημένην, ἀλλὰ τὰς ἀλλοτρίας τρίχας καὶ δι᾽ αὐτῶν ἄλλην κεφαλήν ; Plb.~3.78.2–3 κατεσκευάσατο περιθετὰς τρίχας, ἁρμοζούσας ταῖς κατὰ τὰς ὁλοσχερεῖς διαφορὰς τῶν ἡλικιῶν ἐπιπρεπείαις, / καὶ ταύταις ἐχρῆτο συνεχῶς μετατιθέμενος· ὁμοίως δὲ καὶ τὰς ἐσθῆτας μετελάμβανε τὰς καθηκούσας ἀεὶ ταῖς περιθεταῖς ; \textit{idem} Suid.~ε~2534 (ὁ δὲ Ἀννίβας ὁ Καρχηδονίων στρατηγὸς κατεσκευάσατο περιθετὰς τρίχας ; cetera om.) ; Joseph.~\textit{Vit.}~47 φεύγει περιθετῇ χρησάμενος κόμῃ τοῦ μὴ κατάδηλος γενέσθαι ; Suet.~\textit{Cal.}~11 ganeas atque adulteria capillamento celatus et ueste longa noctibus obiret ; Polyaen.~5.42 Χαριμένης Μιλήσιος … περιθέμενος περιθετὴν πεζῇ διὰ τῆς Περικλέους χώρας ἐσώθη ; Σ Ar.~\textit{Thesm.}~258 ὅτι τὴν περίθετον οὕτως. καὶ τοῦτο δὲ εἰς μαλακίαν, ἵνα λανθάνῃ νυκτὸς καὶ ὡς γυνὴ πάσχῃ. λέγουσι δὲ περίθετον πρόσωπον, ὡς Ἀριστομένης ἐν Γόησιν ; Hypoth. Ar. Eccl. I πώγωνας περιθέτους ποιοῦνται καὶ ἀνδρείαν ἀναλαμβάνουσι στολήν ; Hsch.~δ~1220 διαπηνηκίσαι· διακεφαλαιῶσαι. πηνήκη γὰρ γένος περικεφαλαίας. οἱ δὲ τὸ ἀπατῆσαι ; Hsch.~ε~2304 ἐμβολίδες· αἱ περιθεταί ⟨ … ⟩ ; Hsch.~π~1729 περικεφαλαία· ἡ ἐκ τριχῶν γεγονυῖα περιθέτη ; Hsch.~π~2205 πηνηκίζειν· ἀπατᾶν ; Poll.~2.30–31 ἐκαλεῖτο δέ τι καὶ πηνήκη, καὶ προκόμιον προσθετόν, οὐ γυναιξὶ μόνον ἀλλὰ καὶ ἀνδρῶν τοῖς ἐνδεῶς εἰς κόμης λόγον πράττουσιν. τὸ δ᾽ αὐτὸ καὶ ἔντριχον ὠνόμαζον ; Phot.~π~862 πηνήκη· περικεφαλαία τριχωτή· καὶ περιθετὴ κόμη· ἔνθεν καὶ πηνηκίζειν τὸ ἀπατᾶν· τῆς δὲ περιθετῆς κόμης τὸ μὲν ἔντριχον, τὸ δὲ προκόμιον, τὸ δὲ πηνήκην ἐκάλουν ; \textit{idem} Suid.~π~1529 (Πηνίκη ; πηνικίζειν ; om. τῆς δὲ περιθετῆς … ἐκάλουν), Apostol.~14.29 (Πηνικίζειν μοι δοκεῖς· ἀντὶ τοῦ ἀπατᾶν. πηνίκη δὲ ἐστι … ; om. ἔνθεν … ἐκάλουν) ; Suid.~δ~698 Διπηνικίσαι: ἐξαπατῆσαι. πηνίκην γὰρ οἱ ἀρχαῖοι Ἀττικοὶ τὴν ἐκ τῶν τριχῶν περιθετὴν κόμην. καὶ τὸ φενακίσαι τινὲς ἀπὸ τούτου εἰρῆσθαι λέγουσιν ; Suid.~κ~978 Κατερρυηκυιῶν γυναικῶν καὶ οὕτως ἀπατουσῶν διὰ τῆς ἐπεισάκτου κόμης. ζήτει ἐν τῷ φενάκη ; Ar.~\textit{Eccl.}~25 ἔχουσι τοὺς πώγωνας, οὓς εἴρητ᾽ ἔχειν ; Ar.~\textit{Eccl.}~127 πώγωνα περιδήσειεν ἐσταθευμέναις ; Ar.~\textit{Eccl.}~273 περιδεῖσθε τοὺς πώγωνας ; Cratin.~fr.~319~Kock (Gramm. Darmst.) πονηροὺς ἀνθρώπους … πηνικίζων ἐξαπατᾷ ; Ov.~\textit{Ars}~3.165 Femina procedit densissima crinibus emptis ; Mart.~6.12.1 Iurat capillos esse, quos emit, suos ; Mart.~12.23.1 Dentibus atque comis—nec te pudet—uteris emptis ; X.~\textit{Cyr.}~1.3.2 καὶ ὁρῶν δὴ αὐτὸν κεκοσμημένον καὶ ὀφθαλμῶν ὑπογραφῇ καὶ χρώματος ἐντρίψει καὶ κόμαις προσθέτοις ; Luc.~\textit{Alex.}~59 ὅτεπερ καὶ ἐφωράθη φαλακρὸς ὤν, … μὴ οὐχὶ τῆς φενάκης ἀφῃρημένης ; Luc.~\textit{DMeretr.}~11.3 τὰ δὲ ἄλλα φενάκη βαθεῖα ; Ael.~\textit{VH}~1.26 περίθετον δὲ εἶχε κόμην καὶ λόφον ἐπὶ τῆς κεφαλῆς ; Ael.~fr.~155 οὔτε προσθέτους οὔτε ἐπακτοὺς κόμας ; Ath.~10.415a τὸ πομπικόν, περιθέτην ἔχουσα καὶ λόφον ἐπὶ τῆς κεφαλῆς, ὡς δηλοῖ Ποσείδιππος ἐν ἐπιγράμμασιν ; Hsch.~κ~2618 ἢ ἐκ τριχῶν ὕφασμα ; Hsch.~π~2206 πηνηκισμάτων· φενακισμάτων ; Hsch.~π~2211 πηνίκη· θρίξ. ἐμπλοκή ; Ar.~\textit{Ach.}~90 ταῦτ᾽ ἄρ᾽ ἐφενάκιζες σὺ δύο δραχμὰς φέρων ; Hsch.~π~2181 πηληκίσματα· φενακίσματα}
\lp{258 ἣν ἐγὼ νύκτωρ φορῶ]}{Ar.~\textit{Thesm.}~35 καὶ μὴν βεβίνηκας σύ γ᾽, ἀλλ᾽ οὐκ οἶσθ᾽ ἴσως ; Ar.~\textit{Thesm.}~204–205 δοκῶν γυναικῶν ἔργα νυκτερείσια / κλέπτειν ὑφαρπάζειν τε θήλειαν Κύπριν ; Eur.~\textit{Ba.}~485–487 τὰ δ᾽ ἱερὰ νύκτωρ ἢ μεθ᾽ ἡμέραν τελεῖς; / νύκτωρ τὰ πολλά· σεμνότητ᾽ ἔχει σκότος. / τοῦτ᾽ ἐς γυναῖκας δόλιόν ἐστι καὶ σαθρόν ; Juv.~6.117–121 ausa Palatino tegetem praeferre cubili, / sumere nocturnos meretrix Augusta cucullos / linquebat comite ancilla non amplius una. / sed nigrum flavo crinem abscondente galero / intravit calidum veteri centone lupanar ; Ov.~\textit{Ars}~1.249–250 Nocte latent mendae, vitioque ignoscitur omni, / horaque formosam quamlibet illa facit ; Hsch.~ν~709 νύκτωρ· νυκτός ; Plu.~\textit{Mor.}~144e (Coniug. praec.~46) πᾶσα γυνὴ τοῦ λύχνου ἀρθέντος ἡ αὐτή ἐστι ; \textit{idem} Tosi n°~303 (Λύχνου ἀρθέντος γυνὴ πᾶσα ἡ αὐτή)}
\lp{259 νὴ τὸν Δί᾽ ἀλλά]}{Ar.~\textit{Plut.}~202 νὴ τὸν Δί᾽ ἀλλὰ καὶ λέγουσι πάντες ὡς ; Ar.~\textit{Vesp.}~912 νὴ τὸν Δί᾽ ἀλλὰ δῆλός ἐστ᾽ ; Ar.~\textit{Lys.}~609 νὴ τὸν Δί᾽ ἀλλὰ τοῖς προβούλοις ἄντικρυς}
\lp{259 κἀπιτηδεία πάνυ]}{Ar.~\textit{Pax}~1254 ἔστιν γὰρ ἐπιτήδεια συρμαίαν μετρεῖν ; Ar.~\textit{Eccl.}~79 νὴ τὸν Δία τὸν σωτῆρ᾽ ἐπιτήδειός γ᾽ ἂν ἦν ; Σ Ar.~\textit{Thesm.}~259 πάντα ἔχουσα τὰ τοιαῦτα ἐν ἑαυτῇ ; Hsch.~ε~5332 ἐπιτήδειον· ἕτοιμον. ὠφέλιμον ; Hsch.~ε~5334 ἐπιτηδείως· ἐπιμελῶς, ἁρμοδίως ; Ar.~\textit{Pax}~1228 ἐναποπατεῖν γάρ ἐστ᾽ ἐπιτήδειος πάνυ— ; Ar.~\textit{Av.}~836 ὡς δ᾽ ὁ θεὸς ἐπιτήδειος οἰκεῖν ἐπὶ πετρῶν ; Ar.~\textit{Thesm.}~233 μὴ φροντίσῃ· ὡς εὐπρεπὴς φανεῖ πάνυ ; Ar.~\textit{Thesm.}~639 ὡς καὶ στιβαρά τις φαίνεται καὶ καρτερά ; Ar.~\textit{Ran.}~1307 πρὸς ἥνπερ ἐπιτήδεια ταῦτ᾽ ᾄδειν μέλη ; Hsch.~ε~5331 ἐπίτηδες· ἑκούσιον. ἐπιμελές ; Hsch.~ε~5333 ἐπιτήδειος· ὁ φίλος}
\lp{260 ἆρ᾽ ἁρμόσει μοι;]}{Ar.~\textit{Thesm.}~263 ἆρ᾽ ἁρμόσει μοι ; Eur.~\textit{Ba.}~929 οὐχ ὡς ἐγώ νιν ὑπὸ μίτρᾳ καθήρμοσα ; Pl.~\textit{Hp.Ma.}~294a ὥσπερ γε ἐπειδὰν ἱμάτιά τις λάβῃ ἢ ὑποδήματα ἁρμόττοντα, κἂν ᾖ γελοῖος, καλλίων φαίνεται ; X.~\textit{HG}~2.3.31 καὶ γὰρ ὁ κόθορνος ἁρμόττειν μὲν τοῖς ποσὶν ἀμφοτέροις δοκεῖ ; D.L.~2.41 εἰπόντος δὲ τοῦ Λυσίου, πῶς, εἰ καλός ἐστιν ὁ λόγος, οὐκ ἄν σοι ἁρμόττοι; ἔφη, οὐ γὰρ καὶ ἱμάτια καλὰ καὶ ὑποδήματα εἴη ἂν ἐμοὶ ἀνάρμοστα ; Σ Ar.~\textit{Ran.}~47 κόθορνος : Τινὲς ὅτι ὁ κόθορνος εἰς ἀμφοτέρους τοὺς πόδας ἁρμόζει, οἱ δὲ ὅτι ἀνδράσι καὶ γυναιξὶν ἁρμόττει ; Hsch.~κ~3214 κόθορνος· ὑπόδημα ἀμφοτέροις τοῖς ποσὶ πεποιημένον τινὲς δὲ καὶ ἀνδράσι καὶ γυναιξὶν ἐφαρμόττειν φασὶ τὸ ὑπόδημα τὸ κόθορνον ; Zen.~3.93 Εὐμεταβολώτερος κοθόρνου: κόθορνός ἐστιν ὑποδήματος εἶδος ἐφαρμόζοντος καὶ δεξιῷ καὶ ἀριστερῷ ποδί· ὅθεν καὶ Θηραμένην, τὸν ἐπὶ τῶν λ᾽, κόθορνον ἐκάλουν οἱ Ἀθηναῖοι ; X.~\textit{Mem.}~3.10.10 εὐρυθμοτέρους ποιῶ ; D.Chr.~40.11 καίτοι τὸ μὲν ὑπόδημα δεῖ πρὸς αὑτὸν ἁρμόττειν καὶ τὸν αὑτοῦ πόδα, καὶ ἂν δοκῇ μεῖζον, περιτεμεῖν ; Σ Ar.~\textit{Ran.}~541 ἐπειδὴ καὶ ὁ κόθορνος ἀνδράσι καὶ γυναιξὶ πρὸς τὰς ὑποδέσεις ἁρμόττει ; Suid.~π~1258 Περὶ πόδα: οἷον ἁρμόττον σφόδρα, ὡς τὰ ὑποδήματα τοῖς ποσί ; Diogenian.~4.72 Εὐμεταβολώτερος κοθόρνου: ἐπὶ τῶν πᾶσιν ἐφαρμοζόντων. Ἔστι δὲ εἶδος ὑποδήματος ἐφαρμόζον τοῖς δυσὶ ποσίν ; Erasm.~\textit{Adag.}~1.1.94 Cothurno versatilior. Εὐμεταβολώτερος κοθόρνου, id est, Versatilior cothurno ; Hsch.~η~810 ἡρμοσμένου· εὖ κατεσκευασμένου}
\lp{260 νὴ Δί᾽ ἀλλ᾽ ἄριστ᾽ ἔχει]}{Ar.~\textit{Lys.}~927 ἀλλ᾽ οὐδὲ δέομ᾽ ἔγωγε. — νὴ Δί᾽ ἀλλ᾽ ἐγώ ; Ar.~\textit{Thesm.}~207 ἀτὰρ ἡ πρόφασίς γε νὴ Δί᾽ εἰκότως ἔχει ; Ar.~\textit{Thesm.}~640 καὶ νὴ Δία τιτθούς γ᾽ ὥσπερ ἡμεῖς οὐκ ἔχει ; Ar.~\textit{Ran.}~1161 ἄνθρωπε ταὔτ᾽ ἔστ᾽, ἀλλ᾽ ἄριστ᾽ ἐπῶν ἔχον ; Ar.~\textit{Nub.}~522 καὶ ταύτην σοφώτατ᾽ ἔχειν τῶν ἐμῶν κωμῳδιῶν ; Ar.~\textit{Vesp.}~966 νὴ Δί᾽ ἀλλὰ ψεύδεται ; Pl.~\textit{Tht.}~176d μακρῷ ἄριστ᾽ ἔχει ; Aeschin.~1.17 εὑρήσετε ὅτι πάντων ἄριστα ἔχει}
\lp{261 φέρ᾽ ἔγκυκλον· τουτὶ λάβ᾽ ἀπὸ τῆς κλινίδος]}{\textit{Test.} Suid.~ε~135 Ἔγκυκλον: ἱμάτιον. Ἀριστοφάνης· ἐγώ γ᾽ ἂν κἂν εἴ με χρεῖ᾽ ᾖ τοὔγκυκλον τουτὶ καταθεῖσαν. καὶ ἀλλαχοῦ· φέρ᾽ ἔγκυκλον· τουτὶ λάμβαν᾽ ἀπὸ τῆς κλινίδος. καὶ δῆλον, ὡς τὸ μὲν ἔγκυκλόν ἐστιν ἱμάτιον, ὁ δὲ κροκωτὸς ἔνδυμα ; Suid.~κ~1809 Κλινίς: βάθριόν τι ἡ κλινὶς εἰστεταμένον, ἔγκυκλον. τουτὶ λάμβαν᾽ ἀπὸ τῆς κλινίδος}
\lp{261 ἔγκυκλον]}{Ar.~\textit{Lys.}~113–114 ἔγωγ᾽ ἂν ⟨οὖν⟩ κἂν εἴ με χρείη τοὔγκυκλον / τουτὶ καταθεῖσαν ἐκπιεῖν αὐθημερόν ; Ar.~\textit{Lys.}~1162–1163 ἁμές γε λῶμες, αἴ τις ἁμὶν τὤγκυκλον / λῇ τοῦτ᾽ ἀποδόμεν. — ποῖον ὦ τᾶν; — τὰν Πύλον ; Ar.~\textit{Thesm.}~499–500 ὡς ἡ γυνὴ δεικνῦσα τἀνδρὶ τοὔγκυκλον / † ὑπ᾽ αὐγὰς † οἷόν ἐστιν ; Ar.~\textit{Eccl.}~536 εἶτ᾽ οὐ τὸ σαυτῆς ἱμάτιον ἐχρῆν σ᾽ ἔχειν; / ἀλλ᾽ ἔμ᾽ ἀποδύσασ᾽ ἐπιβαλοῦσα τοὔγκυκλον ; Ar.~fr.~332~K.-A. (=~320~Kock ; Poll.~7.95) χιτῶνα, βάραθρον, ἔγκυκλον, κομμώτριον ; Σ Ar.~\textit{Lys.}~113 τοὔγκυκλον: τὸ περιβόλαιον, (ἢ) τὸ ζωνάριον ; Σ Ar.~\textit{Thesm.}~261 τῆς κλίνης. στρῶμα δίδωσιν αὐτῷ. βάθριόν τι ἡ κλινὶς εἰστεταμένον. δῆλον δὲ ὅτι τὸ ἔγκυκλον ἱμάτιον, ὁ δὲ κροκωτὸς ἔνδυμα ; Σ Ar.~\textit{Thesm.}~656 τῶν ἐγκύκλων ἀποδυσαμένας ; Hsch.~π~1767 περίνησον· ἱμάτιον ἔχον πορφύραν κύκλῳ ; Poll.~7.53 ἢ πορφύρα κύκλῳ τὰ τέλη τοῦ ὑφάσματος περιέρχεται, νήσου σχῆμα ποιοῦσα τῇ περιρροῇ τοῦ χρώματος … τοιοῦτον δ᾽ ἂν εἴη καὶ τὸ ἔγκυκλον ; Phot.~π~280 παράπηχυ· ἱμάτιον τὸ παρ᾽ ἑκάτερον μέρος ἔχον πορφύραν· τοῦτο δὲ καὶ παρυφὲς καλοῦσι· τὸ δὲ κύκλῳ τὴν πορφύραν ἔχον ἔγκυκλον ; Suid.~κ~2461 ὅτι ὁ μὲν κροκωτὸς ἔνδυμά ἐστι, τὸ δὲ ἔγκυκλον ἱμάτιον καὶ ζήτει ἐν τῷ ἔγκυκλον ; Eust.~976.13–15 καθάπερ ὁ ἑανός, οὕτω καὶ τὸ ἔκκυκλον ἢ ἔγκυκλον, οὗ καὶ ὁ κωμικὸς μέμνηται, γυναικεῖον ἦν ποτε φόρημα. φησὶ γοῦν Παυσανίας ὅτι ἔγκυκλον περιπόρφυρον ἱμάτιον καὶ χιτὼν γυναικεῖος, ὃν ἔνδοθεν ἐνδύονται γυναῖκες, εἶτα τὸ ἔνδυμα. λέγεται δέ, φησίν, ὑπ᾽ ἐνίων ἐπιγονατίς, διότι μέχρι γονάτων διήκει ; Anaxil.~fr.~34~K.-A. (=~35~Kock ; Poll.~7.53) καὶ πῶς γυνή, ὥσπερ θάλαττα, νῆσον ἀμφιέννυται ; Ael.~\textit{VH}~7.9 οὐ κροκωτοῦ οὐ ταραντινιδίου οὐκ ἀναβολῆς οὐκ ἐγκύκλου ; Σ Ar.~\textit{Eccl.}~536 τὸ γυναικεῖον ἱμάτιον. R ; Σ Ar.~\textit{Lys.}~1162 ἡμεῖς θέλομεν, ἐάν τις ἡμῖν θελήσῃ ἀποδοῦναι τὸ ἔγκυκλον, τουτέστι τὴν Πύλον, διὰ τὴν τῶν τειχῶν περιβολήν ; \textit{idem} Suid.~τ~812 (Τοὔγκυκλον· τὸ ἔγκυκλον pro ἡμεῖς θέλομεν … τὸ ἔγκυκλον) ; Hsch.~π~1765 περίνησσαια· περιβόλαια, οἷς ἐν κύκλῳ πορφύρα παράκειται ; Phot.~ε~63 ἔγκυκλον· ἱματίου εἶδος ; Phot.~ε~977 ἔγκυκλον· ἱμάτιον περιπόρφυρον ; Phot.~π~708 περίνησα· περιβόλαια περιφερῆ καὶ νησοειδῆ}
\lp{261 ἀπὸ τῆς κλινίδος]}{Ar.~\textit{Lys.}~916 φέρε νυν ἐνέγκω κλινίδιον νῷν ; Cratin.~fr.~137~Kock οἱ δ᾽ ἀλυσκάζουσιν ὑπὸ ταῖς κλινίσιν ; \textit{idem} Poll.~10.33 (ἐπὶ ; praem. ἢ κλινίδα ὡς ἐν Ὀδυσσεῦσι Κρατῖνος. ἐπὶ δὲ τῶν κλινῶν ἀκουστέον, ὅταν φῇ ; add. οὐ μέντοι ἀγνοῶ ὅτι κλινὶς ἐκαλεῖτο τὸ ἐπὶ τῆς ἁμάξης καταστορνύμενον, ὅταν μετίωσι τὰς νύμφας, ἐφ᾽ οὗ κάθηται ἡ νύμφη μεταξὺ τοῦ παρόχου τε καὶ τοῦ νυμφίου) ; Hsch.~κ~3001 κλινίς· ἐπὶ τῆς ἁμάξης νυμφικὴ καθέδρα ; Phot.~κ~787 κλινίς· κλίνη βραχεῖα, οἷον ἥ τε παράβυστος καὶ ἡ ἐπὶ τῷ ζεύγει τοῦ νυμφίου καὶ τῆς νύμφης τιθεμένη ; Ar.~\textit{Ach.}~1090 κλῖναι τράπεζαι προσκεφάλαια στρώματα ; Ar.~\textit{Eccl.}~840 κλῖναί τε σισυρῶν καὶ δαπίδων †νενασμέναι† ; Ar.~\textit{Plut.}~527 ἔτι δ᾽ οὐχ ἕξεις οὔτ᾽ ἐν κλίνῃ καταδαρθεῖν ; Hsch.~κ~3011 κλινίδιον· ὑποκοριστικῶς τὸ αὐτό}
\lp{261 τουτὶ λάβ᾽ (262 λάμβανε)]}{Ar.~\textit{Ach.}~1103–1104 ἔνεγκε δεῦρο τὼ πτερὼ τὼ ᾽κ τοῦ κράνους. / ἐμοὶ δὲ τὰς φάττας γε φέρε καὶ τὰς κίχλας ; Eur.~\textit{Hel.}~663–665 λέγ᾽· ὡς ἀκουστὰ πάντα δῶρα δαιμόνων. / … / ὅμως δὲ λέξον}
\lp{262 ὑποδημάτων δεῖ]}{Ar.~\textit{Vesp.}~1157–1158 ἄγε νυν ὑπολύου τὰς καταράτους ἐμβάδας, / τασδὶ δ᾽ ἀνύσας † ὑπόδυθι † τὰς Λακωνικάς ; Ar.~\textit{Thesm.}~142 καὶ ποῦ πέος; ποῦ χλαῖνα; ποῦ Λακωνικαί ; Ar.~\textit{Thesm.}~734 οἴνου πλέως καὶ ταῦτα Περσικὰς ἔχων ; Ar.~\textit{Eccl.}~269 ὑποδεῖσθε δ᾽ ὡς τάχιστα τὰς Λακωνικάς ; Ar.~\textit{Eccl.}~319 καὶ τὰς ἐκείνης Περσικὰς ὑφέλκομαι ; Ar.~\textit{Eccl.}~344–346 μὰ τὸν Διόνυσον οὐδ᾽ ἐγὼ γὰρ τὰς ἐμὰς / Λακωνικάς, ἀλλ᾽ ὡς ἔτυχον χεζητιῶν, / ἐς τὼ κοθόρνω τὼ πόδ᾽ ἐνθεὶς ἵεμαι ; Theopomp. Com. fr.~52~Kock (Suid.~π~1080) ὑποδοῦ λαβὼν τὰς περιβαρίδας ; \textit{idem} Suid.~π~1080 (praem. Περίβαρα· ὑποδήματα. καὶ Περιβαρίδες· Θεόπομπος ἐν Σειρῆσιν ; add. ἢ Περιβαρίδες, ὑποδήματα γυναικῶν) ; Hdt.~1.155.4 κέλευε δὲ σφέας κιθῶνάς τε ὑποδύνειν τοῖσι εἵμασι καὶ κοθόρνους ὑποδέεσθαι ; X.~\textit{Oec.}~9.6 ὑποδήματα γυναικεῖα, ὑποδήματα ἀνδρεῖα ; Plu.~\textit{Sol.}~8.5 ἐνδύμασι καὶ μίτραις καὶ ὑποδήμασι τοῖς ἐκείνων σκευασαμένους ; Ath.~4.155c καὶ ἐνδυόμενον αὐτὸν κροκωτὸν καὶ ὑποδούμενον Σικυώνια διατελεῖν ὀρχούμενον ; Σ Ar.~\textit{Nub.}~151 Περσικαὶ : Ἔστι μὲν δένδρα, ἔστι δὲ καὶ ὑποδημάτων εἶδος γυναικείων ; Hsch.~β~363 βαυκίδες· εἶδος ὑποδήματος γυναικείου ; Hsch.~λ~225 Λακωνικαί· ὑποδήματα ἀνδρεῖα, ἃ καλεῖται Ἀμυκλᾷδες ; Hsch.~π~1585 περιβαρίδες· ὑποδημάτων εἶδος γυναικεῖον ; Hsch.~ρ~50 ῥάιδια· σανδάλια. οἱ δὲ ὑποδήματα γυναικεῖα ; Poll.~7.92 ἴδια δὲ γυναικῶν ὑποδήματα Περσικαί· λευκὸν ὑπόδημα, μᾶλλον ἑταιρικόν. περιβαρίς· θεραπαινίδων μᾶλλον τὸ ὑπόδημα ; Poll.~7.93 τὰ δὲ πλοιάρια καὶ ἀκάτια ὀνομάζει Ἀριστοφάνης, εἶδος ὄντα γυναικείων ὑποδημάτων ; Poll.~7.94 αἱ δὲ βαυκίδες καὶ βαυκίδια ἐλέγοντο· πολυτελὲς δ᾽ ἦν ὑπόδημα κροκοειδές. γυναικεῖα δὲ καὶ τὰ ἄφρακτα, ὀπισθοκρηπῖδες, καννάβια, γυμνοπόδια, περίβαρα, νυκτιπήδηκες, μεσοπερσικαί, φιττάκια, Σελευκίδες, νοσσίδες, Ἀμβρακίδες, ἀμφίσφυρα, ἀκροσφύρια ; Ar.~\textit{Nub.}~151 κᾆτα ψυχείσῃ περιέφυσαν Περσικαί ; Ar.~\textit{Lys.}~53 κτήσομαι περιβαρίδας ; Ar.~\textit{Lys.}~229 οὐ πρὸς τὸν ὄροφον ἀνατενῶ τὼ Περσικά ; Ar.~\textit{Lys.}~657 τῷδέ γ᾽ ἀψήκτῳ πατάξω τῷ κοθόρνῳ τὴν γνάθον ; Ar.~\textit{Ran.}~47 τίς ὁ νοῦς; τί κόθορνος καὶ ῥόπαλον ξυνηλθέτην ; Ar.~\textit{Eccl.}~74 Λακωνικὰς γὰρ ἔχετε καὶ βακτηρίας ; Ar.~\textit{Eccl.}~542–544 αἱ δὲ δὴ Λακωνικαὶ / ᾤχοντο μετὰ σοῦ κατὰ τί χἠ βακτηρία; / ἵνα θοἰμάτιον σώσαιμι, μεθυπεδησάμην ; X.~\textit{Cyr.}~8.2.5 ὑποδήματα ποιεῖ ὁ μὲν ἀνδρεῖα, ὁ δὲ γυναικεῖα ; Luc.~\textit{Lex.}~10 ἐνδρομίδας ὑποδούμενος ἢ βαυκίδας ; Σ Ar.~\textit{Thesm.}~734 εἶδος δὲ ὑποδημάτων αἱ Περσικαί ; Hsch.~α~2301 ἀκάτια· ὑποδημάτων εἶδος ; Hsch.~α~3838 ἀμυκλᾶιδες· εἶδος ὑποδήματος πολυτελοῦς Λακωνικοῦ ; Hsch.~π~2005 περσικά· εἴδη καρύων. ἢ εὐτελῆ ὑποδήματα ; Hsch.~π~2609 πλοιάρια· πλοῖα μικρά (…). ἢ ὑποδήματα ποιά ; Zen.~3.36 Δεξιὸν εἰς ὑπόδημα, ἂριστερὸν εἰς ποδόνιπτρον: ἐπὶ τῶν ἁρμοδίως τοῖς πράγμασι κεχρημένων ; \textit{idem} Tosi n°~1625 (ποδάνιπτρα ; om. ἐπὶ τῶν ἁρμοδίως τοῖς πράγμασι κεχρημένων) ; Ar.~\textit{Vesp.}~103 εὐθὺς δ᾽ ἀπὸ δορπηστοῦ κέκραγεν ἐμβάδας ; Ar.~\textit{Av.}~994 τίς ἡ ᾽πίνοια, τίς ὁ κόθορνος τῆς ὁδοῦ ; Pherecr.~fr.~153.4–6~Kock (Ath.~8.364b) εἶτα γνούς πως τοῦθ᾽ ὑποδεῖται, κᾆτά τις εἶπε / τῶν ξυμπινόντων «ἤδη σύ; τί οὐχ ὑποπίνεις; / οὐχ ὑπολύσεις αὐτόν;» ; Eub.~fr.~30~Kock (Ath.~3.100a) ἐγὼ κεχόρτασμαι μέν, ἄνδρες, οὐ κακῶς, / ἀλλ᾽ εἰμὶ πλήρης, ὥστε καὶ μόλις πάνυ / ὑπεδησάμην ἅπαντα δρῶν τὰς ἐμβάδας ; Hsch.~α~7002 ἀρβύλαι· εἶδος ὑποδημάτων ; Hsch.~ε~2277 ἐμβάδια· τὰ ὑποδήματα ; Hsch.~π~1584 περίβαρα· ὑποδήματα}
\lp{262 τἀμὰ ταυτὶ λάμβανε]}{Ar.~\textit{Ach.}~434 ἰδοὺ ταυτὶ λαβέ ; Ar.~\textit{Thesm.}~219–220 αὐτὸς λάμβανε / ἐντεῦθεν ἐκ τῆς ξυροδόκης ; Hdt.~6.1.2 τοῦτο τὸ ὑπόδημα ἔρραψας μὲν σύ, ὑπεδήσατο δὲ Ἀρισταγόρης ; \textit{idem} Diogenian.~8.49 (ἔῤῥαψε μὲν Ἱστιαῖος ; Ἀρισταγόρας) ; Ar.~\textit{Eq.}~888–889 οὔκ, ἀλλ᾽ ὅπερ πίνων ἀνὴρ πέπονθ᾽ ὅταν χεσείῃ, / τοῖσιν τρόποις τοῖς σοῖσιν ὥσπερ βλαυτίοισι χρῶμαι ; Ar.~\textit{Vesp.}~1135 ἔχ᾽ ἀναβαλοῦ τηνδὶ λαβὼν καὶ μὴ λάλει ; Ar.~\textit{Thesm.}~759 τουτί. λαβέ}
\lp{263 χαλαρὰ γοῦν χαίροις φορῶν]}{\textit{Test.} Suid.~χ~7 Χαλαρά: χαῦνα. χαλαρὰ γοῦν χαίρεις φορῶν ὑποδήματα ; Ar.~\textit{Ach.}~104–106 οὐ λῆψι χρῦσο χαυνόπρωκτ᾽ Ἰαοναῦ. / … / ὅ τι; χαυνοπρώκτους τοὺς Ἰάονας λέγει ; Ar.~\textit{Eq.}~316–318 ὅστις ὑποτέμνων ἐπώλεις δέρμα μοχθηροῦ βοὸς / τοῖς ἀγροίκοισιν πανούργως, ὥστε φαίνεσθαι παχύ, / καὶ πρὶν ἡμέραν φορῆσαι μεῖζον ἦν δυοῖν δοχμαῖν ; Ar.~\textit{Eq.}~321 πρὶν γὰρ εἶναι Περγασῆσιν ἔνεον ἐν ταῖς ἐμβάσιν ; Ar.~\textit{Vesp.}~1495 στρέφεται χαλαρὰ κοτυληδών ; Ar.~\textit{Lys.}~414–419 ὦ σκυτοτόμε μου τῆς γυναικὸς τοῦ ποδὸς / τὸ δακτυλίδιον ξυμπιέζει τὸ ζυγὸν / ἅθ᾽ ἁπαλὸν ὄν· τοῦτ᾽ οὖν σὺ τῆς μεσημβρίας / ἐλθὼν χάλασον, ὅπως ἂν εὐρυτέρως ἔχῃ ; Ar.~\textit{Thesm.}~200–201 καὶ μὴν σύ γ᾽ ὦ κατάπυγον εὐρύπρωκτος εἶ / οὐ τοῖς λόγοισιν ἀλλὰ τοῖς παθήμασιν ; Ar.~\textit{Eccl.}~794 χαρίεντα γοῦν πάθοιμ᾽ ἄν, εἰ μὴ ᾽χοιμ᾽ ὅποι ; Pherecr.~fr.~145.3–5~Kock ([Plut.] Mus.~1141f) ἐμοὶ γὰρ ἦρξε τῶν κακῶν Μελανιππίδης, / ἐν τοῖσι πρῶτος ὃς λαβὼν ἀνῆκέ με / χαλαρωτέραν τ᾽ ἐποίησε χορδαῖς δώδεκα ; Hor.~\textit{Sat.}~1.3.31–32 rusticius tonso toga defluit et male laxus / in pede calceus haeret ; Ov.~\textit{Ars}~1.516 Nec vagus in laxa pes tibi pelle natet ; Hdt.~6.125.3 ἐνδὺς κιθῶνα μέγαν καὶ κόλπον βαθὺν καταλιπόμενος τοῦ κιθῶνος, κοθόρνους τε τοὺς εὕρισκε εὐρυτάτους ἐόντας ὑποδησάμενος ; Pl.~\textit{R.}~398e τίνες οὖν μαλακαί τε καὶ συμποτικαὶ τῶν ἁρμονιῶν; ἰαστί, ἦ δ᾽ ὅς, καὶ λυδιστὶ αὖ τινες χαλαραὶ καλοῦνται ; Thphr.~\textit{Char.}~2.7 καὶ συνωνούμενος ἐπικρηπῖδας τὸν πόδα φῆσαι εἶναι εὐρυθμότερον τοῦ ὑποδήματος ; Thphr.~\textit{Char.}~4.1 καὶ μείζω τοῦ ποδὸς τὰ ὑποδήματα φορεῖν ; Plu.~\textit{Mor.}~208c οὐδὲ σκυτοτόμον ἔφησεν εἶναι σπουδαῖον, ὃς μικρῷ ποδὶ ὑποδήματα μεγάλα περιτίθησι ; \textit{idem} Apostol.~15.57 (Σμικρῷ ποδὶ μεγάλα περιτίθης ὑποδήματα· λέγεται ἐπὶ τῶν ῥητόρων μᾶλλον τῶν τὰ μικρὰ αὐξόντων) ; D.Chr.~77/78.32 χιτῶνά τε ποδήρη καταζωσάμενον καὶ τὸν κόλπον ἐμπλήσαντα γυναικεῖον καὶ βαθὺν καὶ τὰ ὑποδήματα ἐξεπίτηδες μεγάλα καὶ κοῖλα ὑποδησάμενον ; Cic.~\textit{De\,or.}~1.231 ut, si mihi calceos Sicyonios attulisses, non uterer, quamvis essent habiles atque apti ad pedem, quia non essent viriles ; Σ Ar.~\textit{Thesm.}~263 χαῦνα. διαβάλλει δὲ πάλιν τὸν Ἀγάθωνα ὡς χαῦνον ; Hsch.~θ~804 θρύψιχος· τρυφερός, μαλακός. αἰσχρός. χαῦνος ; Hsch.~χ~38 χαλαρῶν· ἀραιῶν ; Σ Luc.~\textit{Pro\,imag.}~10 μηδὲ ὑπὲρ τὸν πόδα ἔσται τὸ ὑπόδημα: παροιμία ἐπὶ τῶν ὑπερτείνειν ἀξιούντων τὰ καθ᾽ ἑαυτοὺς πράγματα, ὃ ποιεῖν φιλοῦσιν οἱ αὐχηματίαι ; \textit{idem} Apostol.~11.34b (ἔστω ; ἤγουν μηδὲν ἔξω τοῦ μέτρου ποίει καὶ τῆς δυνάμεως pro παροιμία … αὐχηματίαι), Tosi n°~1626 (ἔστω ; cetera om.) ; Ar.~\textit{Eq.}~87 ἰδού γ᾽ ἄκρατον. περὶ πότου γοῦν ἐστί σοι ; Ar.~\textit{Thesm.}~845 ἀξία γοῦν εἶ τόκου τεκοῦσα τοιοῦτον τόκον ; Ar.~\textit{Plut.}~565 πάνυ γοῦν κλέπτειν κόσμιόν ἐστιν ; Hor.~\textit{Ep.}~1.10.42–43 cui non conveniet sua res, ut calceus olim, / si pede maior erit, subvertet, si minor, uret ; Plu.~\textit{Mor.}~770c ἔφη γὰρ Ἀγάθωνα τὸν καλὸν ἢδη γενειῶντα περιβάλλων καὶ κατασπαζόμενος, ὅτι τῶν καλῶν καὶ τὸ μετόπωρον ; Ael.~\textit{VH}~2.21 Ἀγάθωνος ἤρα τοῦ ποιητοῦ Παυσανίας ὁ ἐκ Κεραμέων ; Ael.~\textit{VH}~13.4 εἶτα συγκλιθέντα αὐτῷ Ἀγάθωνα τὸν τῆς τραγῳδίας ποιητὴν περιλαβὼν κατεφίλει ; Hsch.~β~657 βλαδαρόν· ἐκλελυμένον. χαῦνον ; Hsch.~δ~1284 διαρρέοντες· χαῦνοι, καὶ διαλελυμένοι ; Hsch.~κ~1365 καταπύγου· κιναίδου, ἤγουν ἀσελγοῦς ; Hsch.~μ~239 μανόν· ἀραιόν μικρόν. ὀλίγον. χαῦνον ; Epict.~\textit{Ench.}~39 μέτρον κτήσεως τὸ σῶμα ἑκάστῳ ὡς ὁ ποὺς ὑποδήματος. … καθάπερ καὶ ἐπὶ τοῦ ὑποδήματος, ἐὰν ὑπὲρ τὸν πόδα ὑπερβῇς ; Hsch.~χ~35 χαλαρά· ἅλυσις}
\lp{263 ἆρ᾽ ἁρμόσει μοι;]}{Ar.~\textit{Thesm.}~260 ἆρ᾽ ἁρμόσει μοι ; Ov.~\textit{Fast.}~2.324 scindebant magni vincula parva pedes ; Plu.~\textit{Mor.}~141a τὸν κάλτιον αὐτοῖς προτείνας; καὶ γὰρ οὗτος ἔφη καλὸς ἰδεῖν καὶ καινός, ἀλλ᾽ οὐδεὶς οἶδεν ὅπου με θλίβει ; \textit{idem} Tosi n°~2192 (om. τὸν κάλτιον αὐτοῖς προτείνας ; ἔφη) ; Plu.~\textit{Aem.}~5.2 προτείνας τὸ ὑπόδημα (κάλτιον αὐτὸ Ῥωμαῖοι καλοῦσιν) εἶπεν ; Luc.~\textit{Hist.Conscr.}~23 ἐοικέναι παιδίῳ, εἴ που Ἔρωτα εἶδες παίζοντα, προσωπεῖον Ἡρακλέους πάμμεγα ἢ Τιτᾶνος περικείμενον ; Erasm.~\textit{Adag.}~3.6.67 Herculis cothurnos aptare infanti}
\lp{264 σὺ τοῦτο γίγνωσκ᾽]}{Eur.~\textit{Hel.}~1257 αὐτὸς σὺ γίγνωσκ᾽· ἀρκέσει γὰρ ἃν διδῷς ; Eur.~\textit{Hel.}~1256 τίνος; σύ μοι σήμαινε, πείσομαι δ᾽ ἐγώ ; Aesch.~\textit{Sept.}~650–652 [σὺ δ᾽ αὐτὸς ἤδη γνῶθι τίνα πέμπειν δοκεῖ·] / … / μέμψῃ, σὺ δ᾽ αὐτὸς γνῶθι ναυκληρεῖν πόλιν ; Aesch.~\textit{PV}~311 γίγνωσκε σαυτὸν καὶ μεθάρμοσαι τρόπους ; Epict.~1.25.24 καὶ σὺ γίγνωσκε}
\lp{264 ἀλλ᾽ ἔχεις γὰρ ὧν δέει]}{Ar.~\textit{Ach.}~448 ἀτὰρ δέομαί γε πτωχικοῦ βακτηρίου ; Ar.~\textit{Ach.}~470 ἀπολεῖς μ᾽. ἰδού σοι. φροῦδά μοι τὰ δράματα ; Ar.~\textit{Ach.}~479 ἁνὴρ ὑβρίζει· κλῇε πηκτὰ δωμάτων ; Ar.~\textit{Plut.}~1135 εἴ του δέει γ᾽ ὧν δυνατός εἰμί σ᾽ ὠφελεῖν ; Soph.~\textit{OT}~1409–1411 ἀλλ᾽ οὐ γὰρ αὐδᾶν ἔσθ᾽ ἃ μηδὲ δρᾶν καλόν, / ὅπως τάχιστα πρὸς θεῶν ἔξω μέ που / καλύψατ᾽ ἢ φονεύσατ᾽ ; Hsch.~δ~415 δέει· χρῄζεις ; Ar.~\textit{Ach.}~464 ὦνθρωπ᾽ ἀφαιρήσει με τὴν τραγῳδίαν ; Ar.~\textit{Ach.}~1026 εἶτα νυνὶ τοῦ δέει ; Ar.~\textit{Eq.}~219 ἔχεις ἅπαντα πρὸς πολιτείαν ἃ δεῖ ; Ar.~\textit{Plut.}~827 ἔπειτα τοῦ δέει ; Eur.~\textit{Ion}~1018 ὦ φιλτάτη παῖ, πάντ᾽ ἔχεις ὅσων σε δεῖ ; Soph.~\textit{El.}~256 ἀλλ᾽ ἡ βία γὰρ ταῦτ᾽ ἀναγκάζει με δρᾶν, / σύγγνωτε ; Soph.~\textit{Phil.}~81 ἀλλ᾽ ἡδὺ γάρ τι κτῆμα τῆς νίκης λαβεῖν ; Σ Ar.~\textit{Plut.}~827 τοῦ δέει: Τίνος χρείαν ἔχεις. … τίνος χρῄζεις ; Σ Ar.~\textit{Plut.}~1135 δέει: Χρῄζεις. … δέῃ: Τινὸς χρείαν ἔχεις ; Ar.~\textit{Nub.}~798 ἀλλ᾽ οὐκ ἐθέλει γὰρ μανθάνειν. τί ἐγὼ πάθω ; Eur.~\textit{Alc.}~1132 πρόσειπ᾽. ἔχεις γὰρ πᾶν ὅσονπερ ἤθελες}
\lp{265 εἴσω τις ὡς τάχιστά μ᾽ ἐσκυκλησάτω]}{Ar.~\textit{Eq.}~1249 κυλίνδετ᾽ εἴσω τόνδε τὸν δυσδαίμονα ; Ar.~\textit{Thesm.}~238 ἐνεγκάτω τις ἔνδοθεν δᾷδ᾽ ἢ λύχνον ; Eur.~fr.~310 κομίζετ᾽ εἴσω τόνδε τὸν δυσδαίμονα ; Soph.~\textit{OT}~1429 ἀλλ᾽ ὡς τάχιστ᾽ ἐς οἶκον ἐσκομίζετε ; Ar.~\textit{Vesp.}~529 ἐνεγκάτω μοι δεῦρο τὴν κίστην τις ὡς τάχιστα ; Ar.~\textit{Vesp.}~860 ἀλλ᾽ ὡς τάχιστα πῦρ τις ἐξενεγκάτω ; Eur.~\textit{Heracl.}~698–699 ἀλλ᾽ εἴσιθ᾽ εἴσω κἀπὸ πασσάλων ἑλὼν / ἔνεγχ᾽ ὁπλίτην κόσμον ὡς τάχιστά μοι ; Eur.~fr.~671 κομίζετ᾽ εἴσω τήνδε· πιστεύειν δὲ χρὴ / γυναικὶ μηδὲν ὅστις εὖ φρονεῖ βροτῶν ; Soph.~\textit{Ant.}~578 κομίζετ᾽ εἴσω, δμῶες ; Ar.~\textit{Ach.}~805 ἐνεγκάτω τις ἔνδοθεν τῶν ἰσχάδων ; Ar.~\textit{Plut.}~228 τῶν ἔνδοθέν τις εἰσενεγκάτω λαβών}
\lp{265 ὡς τάχιστα]}{Ar.~\textit{Thesm.}~662 ἀλλὰ τὴν πρώτην τρέχειν χρῆν ὡς τάχιστ᾽ ἤδη κύκλῳ ; Ar.~\textit{Thesm.}~1214 ἀπότρεκ᾽ ὠς τάκιστα σύ}
\lp{265 ἐσκυκλησάτω]}{Ar.~\textit{Ach.}~408–409 ἀλλ᾽ ἐκκυκλήθητ᾽. — ἀλλ᾽ ἀδύνατον. — ἀλλ᾽ ὅμως. / ἀλλ᾽ ἐκκυκλήσομαι· καταβαίνειν δ᾽ οὐ σχολή ; Ar.~\textit{Vesp.}~1474–1475 νὴ τὸν Διόνυσον ἄπορά γ᾽ ἡμῖν πράγματα / δαίμων τις ἐσκεκύκληκεν ἐς τὴν οἰκίαν ; Ar.~\textit{Thesm.}~96 καὶ ποῖός ἐστιν; — οὗτος οὑκκυκλούμενος ; Σ Ar.~\textit{Ach.}~408 Εἰ μὴ σχολὴν ἔχεις κατελθεῖν, ἀλλ᾽ ἐκκυκλήθητι, τουτέστι συστράφηθι. ἐκκύκλημα δὲ λέγεται μηχάνημα ξύλινον τροχοὺς ἔχον, ὅπερ περιστρεφόμενον τὰ δοκοῦντα ἔνδον ὡς ἐν οἰκίᾳ πράττεσθαι καὶ τοῖς ἔξω ἐδείκνυε, λέγω δὴ τοῖς θεαταῖς ; \textit{idem} Suid.~ε~132 (ἐγκυκλήθητι ; ἐγκύκλημα) ; Poll.~4.128 καὶ τὸ μὲν ἐκκύκλημα, ἐπὶ ξύλων ὑψηλὸν βάθρον, ᾧ ἐπίκειται θρόνος. δείκνυσι δὲ καὶ τὰ ὑπὸ τὴν σκηνὴν ἐν ταῖς οἰκίαις ἀπόρρητα πραχθέντα. καὶ τὸ ῥῆμα τοῦ ἔργου καλεῖται ἐγκυκλεῖν. ἐφ᾽ οὗ δὲ εἰσάγεται τὸ ἐγκύκλημα, εἰσκύκλημα ὀνομάζεται· καὶ χρὴ τοῦτο νοεῖσθαι καθ᾽ ἑκάστην θύραν, οἱονεὶ καθ᾽ ἑκάστην οἰκίαν ; Ar.~\textit{Vesp.}~698–699 σκέψαι τοίνυν ὡς ἐξόν σοι πλουτεῖν καὶ τοῖσιν ἅπασιν / ὑπὸ τῶν ἀεὶ δημιζόντων οὐκ οἶδ᾽ ὅπῃ ἐγκεκύκλησαι ; Luc.~\textit{Lex.}~8 εἶτ᾽ ἐχριόμεθα βακχάριδι καὶ εἰσεκύκλησέ τις ἡμῖν τὴν ποδοκτύπην καὶ τριγωνίστριαν ; Clem.Al.~\textit{Protr.}~7.76.5–6 ἤδη δὲ ἐν Ἴωνι τῷ δράματι γυμνῇ τῇ κεφαλῇ ἐκκυκλεῖ τῷ θεάτρῳ τοὺς θεούς ; Σ Ar.~\textit{Thesm.}~277 παρεπιγραφή. ἐκκυκλεῖται ἐπὶ τὸ ἔξω τὸ θεσμοφόριον ; Σ Ar.~\textit{Vesp.}~1475 εἰσκεκύκληκεν: Εἰσήνεγκεν, εἰσήγαγεν ; \textit{idem} Suid.~ει~281 (add. Ἀριστοφάνης· νὴ τὸν Διόνυσον, ἄπορά γ᾽ ἡμῖν πράγματα δαίμων τις εἰσκεκύκληκεν ἐς τὴν οἰκίαν) ; Σ Soph.~\textit{Aj.}~346 ἐνταῦθα ἐκκύκλημά τι γίνεται ἵνα φανῇ ἐν μέσοις ὁ Αἴας ποιμνίοις ; Σ Eur.~\textit{Hipp.}~171 τοῦτο σεσημείωται τῷ Ἀριστοφάνει ὅτι καίτοι τῷ ἐκκυκλήματι χρώμενος τὸ ἐκκομίζουσα προσέθηκεν περισσῶς ; Hsch.~ε~1159 εἰσκυκλεῖσθαι· περιφέρεσθαι ; Hsch.~ε~1465 ἐκκυκλεῖ· ἐκκαλύπτει ; \textit{idem} Synag.~ε~204~Cunningham, Suid.~ε~699 (add. καὶ Ἐκκυκλούμενος, ἐκκύκλημα ἔχων) ; Hsch.~ε~4014 ἐξώστρα· ἐπὶ τῆς σκηνῆς τὸ ἐκκύκλημα ; \textit{idem} Phot.~ε~1294 (ἔξωστρον ; add. Φρύνιχος) ; Poll.~4.129 τὴν δὲ ἐξώστραν ταὐτὸν τῷ ἐγκυκλήματι νομίζουσιν ; Phot.~ε~418 ἐκκύκλημα καὶ ἐκκύκληθρον· μηχάνημα ὑπότροχον, ἐφ᾽ οὗ ἐδείκνυτο τὰ ἐν τῇ σκηνῇ ; Eur.~\textit{HF}~1029–1030 ἴδεσθε, διάνδιχα κλῇθρα / κλίνεται ὑψιπύλων δόμων ; Aesch.~\textit{Cho.}~973 ἴδεσθε χώρας τὴν διπλῆν τυραννίδα ; Soph.~\textit{Aj.}~346–347 ἰδού, διοίγω· προσβλέπειν δ᾽ ἔξεστί σοι / τὰ τοῦδε πράγη, καὐτὸς ὡς ἔχων κυρεῖ ; Soph.~\textit{Ant.}~1293 ὁρᾶν πάρεστιν· οὐ γὰρ ἐν μυχοῖς ἔτι ; Epiced. ap. BKT~5(1) p.~84 (DCLP~64374) [καὶ νῦ]ν ἰά[μ]βων κωμικῶν πεπαυμέν[ος] / [ἡρῷ᾽ ἔπη τ]ὸ λοιπὸν εἰσκυκλήσομ[αι.] ; Luc.~\textit{Deor.Conc.}~9 ἀλλ᾽ ὁ Ἄττης γε, ὦ Ζεῦ, καὶ ὁ Κορύβας καὶ ὁ Σαβάζιος, πόθεν ἡμῖν ἐπεισεκυκλήθησαν οὗτοι ; Luc.~\textit{Hist.Conscr.}~13 ὧν ἀμελήσαντες οἱ πολλοὶ τὰ μηδὲν προσήκοντα ἐπεισκυκλοῦσιν ; Luc.~\textit{Philops.}~29 καὶ τὸ τοῦ λόγου, θεὸν ἀπὸ μηχανῆς ἐπεισκυκληθῆναί μοι τοῦτον ᾤμην ὑπὸ τῆς Τύχης ; Max.Tyr.~17.5e εἰσκυκλεῖς ποιητῶν χορὸν ἀκόλαστον καὶ βακχεύοντα ; Ael.~fr.~145 Hercher ἐπεὶ δὲ ἐς τὸ θέατρον ἐξεκύκλησαν αὐτόν ; Philostr.~\textit{VA}~6.10 τρίποδάς τε ἐσκυκλήσει πίνοντι καὶ χρυσοῦς θρόνους ; Philostr.~\textit{VS}~2.21.3 ἡ μελέτη δὲ τῆς προτεραίας προεωραμένη ἐσεκυκλεῖτο ; Ath.~6.270e θεασάμενος πλῆθος ἰχθύων καὶ ἄλλων παντοδαπῶν ὄψων παρασκευὴν εἰσκυκλουμένην ; Clem.Al.~\textit{Strom.}~7.14.88.4 οὐ γὰρ ἐκκυκλεῖν χρὴ τὸ μυστήριον ; Plb.~11.5.8 τῆς τύχης ὥσπερ ἐπίτηδες ἐπὶ τὴν ἐξώστραν ἀναβιβαζούσης τὴν ὑμετέραν ἄγνοιαν ; Ap.Dysc.~\textit{Synt.}~3.98 ὡς μάτην εἰσκυκλεῖται ἡ τῶν παρῳχημένων χρόνων φωνὴ κατὰ τὴν ἔγκλισιν ; LXX~2 Ma.~2.24 τὴν οὖσαν δυσχέρειαν τοῖς θέλουσιν εἰσκυκλεῖσθαι τοῖς τῆς ἱστορίας διηγήμασιν ; Σ Ar.~\textit{Nub.}~184 ὁρᾷ δὲ ὡς φιλοσόφους κομῶντας, στραφέντος τοῦ ἐγκυκλήματος ; Σ Ar.~\textit{Nub.}~218 παρεγκύκλημα. δεῖ γὰρ κρεμᾶσθαι τὸν Σωκράτην ἐπὶ κρεμάθρας καθήμενον ; Σ Il.~18.476 δαιμονίως τὸν πλάστην αὐτὸς διέπλασεν, ὥσπερ ἐπὶ σκηνῆς ἐκκυκλήσας καὶ δείξας ἡμῖν ἐν φανερῷ τὸ ἐργαστήριον ; Σ Aesch.~\textit{Cho.}~973 ἀνοίγεται ἡ σκηνὴ καὶ ἐπὶ ἐκκυκλήματος ὁρᾶται τὰ σώματα ἃ λέγει διπλῆν τυραννίδα ; Σ Aesch.~\textit{Eum.}~64b καὶ δευτέρα δὲ γίνεται φαντασία· στραφέντα γὰρ μηχανήματα ἔνδηλα ποιεῖ τὰ κατὰ τὸ μαντεῖον ὡς ἔχει ; Σ Eur.~\textit{Med.}~96 τάδε λέγει ἡ Μήδεια ἔσω οὖσα, οὐδέπω ἐκκεκυκλημένη ; Hsch.~β~563 σημαίνει δὲ καὶ τὸ κύκλημα (γράφεται δὲ καὶ ἐκκύκλημα) ; Hsch.~ε~1160 εἰσκυκλήσας· περιελθών ; Hsch.~ε~6254 †ἐσκύκληνται· διεσκεδασμένοι εἰσίν ; Hsch.~κ~4464 κυκλήσομεν· ἐφ᾽ ἁμαξῶν ἄξομεν ; Synag.~ε~154~Cunningham εἰσκυκλεῖ· εἰσφέρει, καὶ ἐπισυνάπτει. καὶ ἐπεισκυκλεῖ ; \textit{idem} Suid.~ει~289 (add. καὶ Εἰσκυκλεῖσθαι, κυκλοειδῶς τρίβεσθαι) ; Suid.~ε~88 Ἐγκεκύκλησαι: ἀνέστρεψαι. Ἀριστοφάνης· οὐκ οἶδ᾽ ὅπη ἐγκεκύκλησαι}
\lp{266 ἁνὴρ μὲν ἡμῖν οὑτοσὶ καὶ δὴ γυνή]}{Ar.~\textit{Thesm.}~141–143 πότερον ὡς ἀνὴρ τρέφει; … ἀλλ᾽ ὡς γυνὴ δῆτ᾽ ; Ar.~\textit{Thesm.}~862–863 αὖθις αὖ γίγνει γυνή ; Eur.~\textit{Ba.}~822 τί δὴ τόδ᾽; ἐς γυναῖκας ἐξ ἀνδρὸς τελῶ ; Eur.~\textit{Ba.}~915–917 σκευὴν γυναικὸς μαινάδος βάκχης ἔχων, / … / πρέπεις δὲ Κάδμου θυγατέρων μορφὴν μιᾷ ; Σ Ar.~\textit{Thesm.}~266 ὁ Μνησίλοχος ; Ar.~\textit{Thesm.}~136 ποδαπὸς ὁ γύννις; τίς πάτρα; τίς ἡ στολή ; Ar.~\textit{Thesm.}~574 φίλαι γυναῖκες ξυγγενεῖς τοὐμοῦ τρόπου ; \textit{Il.}~8.163 γυναικὸς ἄρ᾽ ἀντὶ τέτυξο ; \textit{idem} Hsch.~γ~1016 (add. ἀντὶ γυναικὸς ἄρα γεγέννησο) ; Pl.~\textit{Lg.}~944d–e Καινέα τὸν Θετταλὸν ἐκ γυναικὸς μεταβαλόντα εἰς ἀνδρὸς φύσιν· … εἰς γυναῖκα ἐξ ἀνδρὸς μεταβαλοῦσα, τιμωρία τούτῳ γενομένη ; Apollod.~3.6.7 θεασάμενος περὶ Κυλλήνην ὄφεις συνουσιάζοντας καὶ τούτους τρώσας ἐγένετο ἐξ ἀνδρὸς γυνή ; Apollod.~3.13.8 Θέτις προειδυῖα ὅτι δεῖ στρατευόμενον αὐτὸν ἀπολέσθαι, κρύψασα ἐσθῆτι γυναικείᾳ ὡς παρθένον Λυκομήδει παρέθετο ; Paus.~8.20.3 ταύτην οἷα δὴ παρθένος πλεξάμενος τὴν κόμην καὶ ἐσθῆτα ἐνδὺς γυναικείαν ἀφίκετο ὡς τὴν Δάφνην ; Hsch.~γ~1015 γύννις· δειλός, ἄνανδρος, γυναικώδης· Ἀριστοφάνης. γύννιδας· μαλακούς ; Hsch.~θ~491 θηλυδριῶδες· τὸ κατακεκλασμένον ; Ar.~\textit{Av.}~829–831 ὅπου θεὸς γυνὴ γεγονυῖα πανοπλίαν / ἕστηκ᾽ ἔχουσα, Κλεισθένης δὲ κερκίδα ; Hsch.~α~4745 ἀνδρόγυνος· ὁ ἑρμαφρόδιτος ; Hsch.~θ~490 θηλυδρίας· ὁ τεθηλυμμένος ; App. Prov.~1.88 Γυνὴ εἰς Ἡρακλέους οὐ φοιτᾷ: πρὸς τοὺς ἀναξίους τινῶν πράξεων, παρόσον Ἡρακλῆς ἐδούλευσεν Ὀμφάλῃ ; Suid.~αι~230 Αἶνος: λόγος παροιμιώδης. … αἶνός τίς ἐστιν· ὡς ἀνήρ τε καὶ οὐκ ἀνήρ, … εὐνοῦχος}
\lp{266 καὶ δή]}{Ar.~\textit{Vesp.}~1324 ὁδὶ δὲ καὶ δὴ σφαλλόμενος προσέρχεται ; Ar.~\textit{Thesm.}~214 ἀπόδυθι τουτὶ θοἰμάτιον. — καὶ δὴ χαμαί ; Hsch.~κ~213 καὶ δή· ἐπὶ τοῦ ἤδη ; \textit{idem} Suid.~κ~1162 (ἤδη pro ἐπὶ τοῦ ἤδη)}
\lp{266–267 καὶ δὴ γυνὴ / τό γ᾽ εἶδος]}{Tosi n°~2183 Barba non facit philosophum — ἀπὸ πώγωνος φιλόσοφοι / σοφοί / σοφισταί}
\lp{267–268 ἢν λαλῇς δ᾽, ὅπως τῷ φθέγματι γυναικιεῖς]}{Ar.~\textit{Thesm.}~191–192 σὺ δ᾽ εὐπρόσωπος λευκὸς ἐξυρημένος / γυναικόφωνος ἁπαλὸς εὐπρεπὴς ἰδεῖν ; Ar.~\textit{Eccl.}~149 ἄγε νυν ὅπως ἀνδριστὶ καὶ καλῶς ἐρεῖς ; Stat.~\textit{Ach.}~1.330–331 et picturato cohibens vestigia limbo / incessum motumque docet fandique pudorem ; Plu.~\textit{Caes.}~10.2–3 τοῦ δὲ Κλωδίου φήσαντος Ἅβραν περιμένειν Πομπηΐας, αὐτὸ τοῦτο καλουμένην, καὶ τῇ φωνῇ γενομένου καταφανοῦς ; Plu.~\textit{Cic.}~28.2 συνεῖσα τὴν φωνὴν οὐ γυναικείαν οὖσαν ἀνέκραγε καὶ συνεκάλει τὰς γυναῖκας ; Luc.~\textit{Eun.}~10 οὑτοσὶ ὁ τὰς γνάθους λεῖος καὶ τὸ φώνημα γυναικεῖος καὶ τὰ ἄλλα εὐνούχῳ ἐοικὼς ; Luc.~\textit{Rh.Pr.}~11 γυναικεῖον τὸ βλέμμα, μελιχρὸν τὸ φώνημα, … πάναβρόν τινα Σαρδανάπαλλον ἢ Κινύραν ἢ αὐτὸν Ἀγάθωνα, τὸν τῆς τραγῳδίας ἐπέραστον ἐκεῖνον ποιητήν ; Philostr.~\textit{VS}~1.8.1 διφυὴς δὲ ἐτέχθη καὶ ἀνδρόθηλυς, καὶ τοῦτο ἐδηλοῦτο μὲν καὶ παρὰ τοῦ εἴδους, … ἐδηλοῦτο δὲ καὶ τῷ φθέγματι, ὀξυηχὲς γὰρ ἠκούετο καὶ λεπτὸν καὶ ἐπίτονον ; Σ Ar.~\textit{Thesm.}~268 ὡς γυνὴ λαλήσεις ; Ar.~\textit{Nub.}~979–980 οὐδ᾽ ἂν μαλακὴν φυρασάμενος τὴν φωνὴν πρὸς τὸν ἐραστὴν / αὐτὸς ἑαυτὸν προαγωγεύων τοῖς ὀφθαλμοῖς ἐβάδιζεν ; Ar.~\textit{Thesm.}~184–187 ἐὰν γὰρ ἐγκαθεζόμενος λάθρᾳ / ἐν ταῖς γυναιξίν, ὡς δοκῶν εἶναι γυνή, / ὑπεραποκρίνῃ μου, σαφῶς σώσεις ἐμέ. / μόνος γὰρ ἂν λέξειας ἀξίως ἐμοῦ ; Crates fr.~24~K.-A. ἀνδριστὶ μιμεῖσθαι φωνήν· ὡς ἀνήρ. Κράτης Λαμίαι ; \textit{Od.}~4.277–279 τρὶς δὲ περίστειξας κοῖλον λόχον ἀμφαφόωσα, / … / πάντων Ἀργείων φωνὴν ἴσκουσ᾽ ἀλόχοισιν ; Arist.~\textit{HA}~538b13 Καὶ περὶ φωνῆς δέ, πάντα τὰ θήλεα λεπτοφωνότερα καὶ ὀξυφωνότερα ; Arist.~\textit{Pr.}~11.62,~906a Διὰ τί τὰ παιδία καὶ αἱ γυναῖκες καὶ οἱ εὐνοῦχοι καὶ οἱ γέροντες φθέγγονται ὀξύ ; D.Chr.~33.38 ὥστε γυναικῶν λαβεῖν φωνὴν ἅπαντας καὶ μηδένα δύνασθαι μήτε νέον μήτε πρεσβύτερον ἀνδριστὶ μηδὲν εἰπεῖν ; Epict.~1.16.12 ἐπὶ τῶν γυναικῶν ὥσπερ ἐν φωνῇ τι ἐγκατέμιξεν ἁπαλώτερον, οὕτως καὶ τὰς τρίχας ἀφεῖλεν ; Quint.~\textit{Inst.}~11.3.91 cum tamen in expositione aut senis sermo, ut in Hydriae prologo, aut mulieris, ut in Georgo, incidit, tremula vel effeminata voce pronuntiant ; Aret.~\textit{SD}~2.5 οἷσι δὲ οὐκ ἔνεστι ζωοῦσα ἡ θορή, ῥικνοί, ἀσθενέες, ὀξύφωνοι, ἄτριχες, ἀγένειοι, γυναικώδεις· δηλοῦσι οἱ εὐνοῦχοι ; Σ Soph.~\textit{Aj.}~334 καὶ γὰρ δεῖ τὸν ὑποκρινόμενον τὸ τοῦ Αἴαντος πρόσωπον ἀπηνεστέρῳ χρήσασθαι τῷ φθέγματι καὶ κυνικώτερον βαύζειν ; Men.~\textit{Mon.}~27 Jäkel ἀνδρὸς χαρακτὴρ ἐκ λόγου γνωρίζεται}
\lp{267 ὅπως … γυναικιεῖς]}{Ar.~\textit{Thesm.}~1204–1205 σὺ δ᾽ ὅπως ἀνδρικῶς / ὅταν λυθῇς τάχιστα θεύξει καὶ τενεῖς ; Ar.~\textit{Nub.}~489 ἄγε νυν ὅπως, ὅταν τι προβάλλω σοι σοφὸν / περὶ τῶν μετεώρων, εὐθέως ὑφαρπάσει ; Arar.~fr.~17~Kock (Suid.~α~1875) ὅπως δὲ τὴν νύμφην, ἐπειδὰν καιρὸς ᾖ, / μετέωρον ἐπὶ τὸ ζεῦγος ἀναθήσεις φέρων}
\lp{267 ἢν λαλῇς δ᾽]}{Ar.~\textit{Thesm.}~90–92 ἐκκλησιάσοντ᾽ ἐν ταῖς γυναιξὶ κἂν δέῃ / λέξονθ᾽ ὑπὲρ ἐμοῦ. — πότερα φανερῶς ἢ λάθρᾳ; / λάθρᾳ, στολὴν γυναικὸς ἠμφιεσμένον ; Ar.~\textit{Thesm.}~392–393 τὰς μοιχοτρόπους, τὰς ἀνδρεραστίας καλῶν, / τὰς οἰνοπότιδας, τὰς προδότιδας, τὰς λάλους ; Ar.~\textit{Ran.}~954 ἔπειτα τουτουσὶ λαλεῖν ἐδίδαξα— — φημὶ κἀγώ ; Eup.~fr.~95~Kock (Plut. Alc.~13.2) λαλεῖν ἄριστος, ἀδυνατώτατος λέγειν ; Ammon.~\textit{Diff.}~293 Nickau λαλεῖν καὶ λέγειν καὶ διαλέγεσθαι διαφέρει. λαλεῖν μὲν γάρ ἐστι τὸ ἀτάκτως ἐκφέρειν τὰ ὑποπίπτοντα ῥήματα, λέγειν δὲ τὸ τεταγμένως προφέρεσθαι τὸν λόγον ; Ar.~\textit{Ach.}~21 οἱ δ᾽ ἐν ἀγορᾷ λαλοῦσι κἄνω καὶ κάτω ; Ar.~\textit{Thesm.}~211–212 πῶς οὖν ποιήσω δῆτα; — τοῦτον μὲν μακρὰ / κλάειν κέλευ᾽, ἐμοὶ δ᾽ ὅ τι βούλει χρῶ λαβών ; Ar.~\textit{Thesm.}~578 ὀλίγῳ τι πρότερον κατ᾽ ἀγορὰν λαλούμενον ; Ar.~\textit{Ran.}~1069 εἶτ᾽ αὖ λαλιὰν ἐπιτηδεῦσαι καὶ στωμυλίαν ἐδίδαξας ; Plaut.~\textit{Aul.}~124–126 nam multum loquaces merito omnes habemur, / nec mutam profecto repertam ullam esse / aut hodie dicunt mulierem aut ullo in saeclo ; Hsch.~λ~240 λαλίστερον· ἀδολεσχέστερον ; Hsch.~λ~245 λαλῶν· φθεγγόμενος ; Soph.~\textit{Aj.}~293 γύναι, γυναιξὶ κόσμον ἡ σιγὴ φέρει}
\lp{267 τό γ᾽ εἶδος]}{Ar.~\textit{Plut.}~317 ὑμεῖς ἐπ᾽ ἄλλ᾽ εἶδος τρέπεσθ᾽ ; Hsch.~ε~773 εἶδος· καῦμα. σῶμα. χρῶμα}
\lp{268 γυναικιεῖς]}{Ar.~\textit{Thesm.}~862–863 αὖθις αὖ γίγνει γυνή, / πρὶν τῆς ἑτέρας δοῦναι γυναικίσεως δίκην ; Diocl. Com. fr.~4~Kock (Antiatt. p.~87.2~Bekker) γυναικισμός· Διοκλῆς Βάκχαις γυναικίζειν φησὶ καὶ γυναικηρόν ; Hp.~\textit{Aer.}~22 ἐνδύονται στολὴν γυναικείην καταγνόντες ἑωυτῶν ἀνανδρείην. γυναικίζουσί τε καὶ ἐργάζονται μετὰ τῶν γυναικῶν ἃ καὶ ἐκεῖναι ; Pl.~\textit{R.}~3.395d–e αἱ μιμήσεις … εἰς ἔθη τε καὶ φύσιν καθίστανται καὶ κατὰ σῶμα καὶ φωνὰς καὶ κατὰ τὴν διάνοιαν; … οὐ δὴ ἐπιτρέψομεν … γυναῖκα μιμεῖσθαι ἄνδρας ὄντας ; Luc.~\textit{Gall.}~19 καὶ ἔρια ἔξαινες καὶ κρόκην κατῆγες καὶ ἐγυναικίζου ἐς τὸ ἑταιρικόν ; Hsch.~α~2422 ἀκκίζειν· γυναικίζεσθαι. προσποιεῖσθαι μωραίνειν ; Suid.~γ~496 Γυναίκισις: ἡ γυναικεία μίμησις. πρὶν τῆς ἑτέρας γυναικίσεως δίκην δοῦναι ; \textit{idem} Zonar. col.~459–460~Tittmann (add. … Γυναικιζόμενος. θηλυνόμενος) ; Lib.~\textit{Or.}~64.74 οἱ διχόθεν μιμοῦνται τὰς γυναῖκας, τῷ τε λόγῳ καὶ τοῖς σχήμασιν; … μεστοὶ γυναικίσεως οἱ παρόντες κατέβησαν ; Plb.~32.15.7 ὅπερ ὁ Προυσίας εἴθιστο ποιεῖν γονυπετῶν καὶ γυναικιζόμενος ; Joseph.~\textit{BJ}~4.563 γυναικιζόμενοι δὲ τὰς ὄψεις ἐφόνων ταῖς δεξιαῖς ; Phld. Mus.~4, col.~43 (DCLP~62477) εἰς γυναικισ[μόν, ὅπερ Ἀγ]άθωνος ἀκατ[απαύστως οἱ] κωμικοὶ κα[τηγοροῦσιν] ; Hsch.~α~2423 ἀκκίζεται· θηλύνεται ; Hsch.~α~8768 ἀφροδισιάζεσθαι· γυναικίζεσθαι ; Suid.~α~878 Ἀκκιζόμενος: προσποιούμενος, γυναικιζόμενος, ἢ μωραίνων, ἀπό τινος γυναικὸς μωρᾶς, ἥτις ἐκαλεῖτο Ἀκκώ ; Suid.~γ~495 Γυναικεία φύσις. Γυναικίας δὲ ἀνὴρ, ὁ θηλυμανής, διὰ τοῦ ι. καὶ Γυναικιζόμενος, ὁ θηλυνόμενος ; Hsch.~α~75 ἀβαριστάν· γυναικιζομένην, καθαιρομένην καταμηνίοις. Κύπριοι ; Phot.~α~2691 ἀπότριψις· μελέτησις, ἀπάγγελσις, ἀγώνισις, … διάρπασις, γυναίκισις, λάλησις ; Diogenian.~2.4 Ἀκκίζεσθαι: Ἀκκὼ γυνή τις ἐπὶ μωρίᾳ διαβαλλομένη, ἥτις ἐσοπτριζομένη τῇ οἰκείᾳ σκιᾷ ὡς ἑτέρᾳ διελέγετο ; Diogenian.~4.10 Γραῦς ἀναθυᾷ: ἀντὶ τοῦ, καπρᾷ. Ἐπὶ τῶν παρ᾽ ἡλικίαν γυναικιζομένων}
\lp{268 εὖ καὶ πιθανῶς]}{Ar.~\textit{Eq.}~629 πιθανώταθ᾽· ἡ βουλὴ δ᾽ ἅπασ᾽ ἀκροωμένη ; Ar.~\textit{Vesp.}~153 εὖ κἀνδρικῶς· κἀγὼ γὰρ ἐνταῦθ᾽ ἔρχομαι ; Ar.~\textit{Thesm.}~464 ἀσύνετ᾽ ἀλλὰ πιθανὰ πάντα ; Ar.~\textit{Thesm.}~656 ξυζωσαμένας εὖ κἀνδρείως τῶν θ᾽ ἱματίων ἀποδύσας ; Eur.~\textit{Or.}~903–906 ἀνήρ τις ἀθυρόγλωσσος, ἰσχύων θράσει, / … / πιθανὸς ἔτ᾽ αὐτοὺς περιβαλεῖν κακῷ τινι ; Th.~3.36.6 Κλέων ὁ Κλεαινέτου, … ὢν καὶ ἐς τὰ ἄλλα βιαιότατος τῶν πολιτῶν τῷ τε δήμῳ παρὰ πολὺ ἐν τῷ τότε πιθανώτατος ; Th.~6.35.2 παρελθὼν δ᾽ αὐτοῖς Ἀθηναγόρας, ὃς δήμου τε προστάτης ἦν καὶ ἐν τῷ παρόντι πιθανώτατος τοῖς πολλοῖς, ἔλεγε τοιάδε ; Pl.~\textit{Grg.}~521a καὶ νῦν εὖ καὶ γενναίως εἰπέ ; Plu.~\textit{Mor.}~783b ὡς εἰρημένον εὖ καὶ πιθανῶς ; Phot.~π~877 πιθανώτατος τοῖς πολλοῖς· ὁ πιστικώτατος καὶ πείθων τοὺς πολλοὺς παρὰ Θουκυδίδῃ (6, 35, 2) ; Ar.~\textit{Eq.}~379 σκεψόμεσθ᾽ εὖ κἀνδρικῶς ; Ar.~\textit{Vesp.}~450 προσαγαγὼν πρὸς τὴν ἐλάαν ἐξέδειρ᾽ εὖ κἀνδρικῶς ; Ar.~\textit{Eccl.}~253 κακῶς κεραμεύειν, τὴν δὲ πόλιν εὖ καὶ καλῶς ; Ar.~\textit{Eccl.}~638 οὐκοῦν ἄγξουσ᾽ εὖ καὶ χρηστῶς ἑξῆς τὸν πάντα γέροντα ; Cratin.~fr.~6~Kock (Ath.~4.164e) εἶδες τὴν Θασίαν ἅλμην οἷ᾽ ἄττα βαΰζει, / ὡς εὖ καὶ ταχέως ἀπετείσατο ; Alex.~fr.~133.5~Kock (Ath.~7.322c) τῷ σιλφίῳ μάστιξον εὖ τε καὶ καλῶς ; Aesch.~\textit{Ag.}~485–487 πιθανὸς ἄγαν ὁ θῆλυς ὅρος ἐπινέμεται / ταχύπορος· ἀλλὰ ταχύμορον / γυναικογήρυτον ὄλλυται κλέος ; Th.~4.21.3 Κλέων ὁ Κλεαινέτου, ἀνὴρ δημαγωγὸς κατ᾽ ἐκεῖνον τὸν χρόνον ὢν καὶ τῷ πλήθει πιθανώτατος ; Pl.~\textit{Lg.}~1.648c καὶ γυμνασάμενον μὲν εὖ καὶ ἀνδρείως ἀζήμιον ἀπαλλάττοις ἄν ; Hsch.~π~2264 πιθανάς· εὐπίστους ; Hsch.~π~2266 πιθανόν· ἀληθινόν. πιστόν. εὐπ⟨ε⟩ιθές}
\lp{268 πειράσομαι]}{Eur.~\textit{Hel.}~1004 ἔχουσα σῴζειν, Μενέλεως, πειράσομαι ; Eur.~\textit{Med.}~724 πειράσομαί σου προξενεῖν δίκαιος ὤν ; Soph.~\textit{Phil.}~1249 αἰσχρὰν ἁμαρτὼν ἀναλαβεῖν πειράσομαι ; Theoc.~15.62 καλλίστα παίδων· πείρᾳ θην πάντα τελεῖται}
\lp{269 βάδιζε τοίνυν]}{Ar.~\textit{Thesm.}~277 ἔκσπευδε ταχέως· ὡς τὸ τῆς ἐκκλησίας / σημεῖον ἐν τῷ Θεσμοφορείῳ φαίνεται ; Ar.~\textit{Thesm.}~485 ἐς τὸν κοπρῶν᾽ οὖν ἔρχομαι. βάδιζέ νυν ; Ar.~\textit{Eccl.}~867 βάδιζε τοίνυν ὕστερος ; Ar.~\textit{Eq.}~1217 βάδιζέ νυν καὶ δεῦρο πρὸς τὴν Παφλαγόνος ; Ar.~\textit{Thesm.}~25 βάδιζε δευρὶ καὶ πρόσεχε τὸν νοῦν. — ἰδού ; Ar.~\textit{Thesm.}~157–158 ὅταν σατύρους τοίνυν ποιῇς, καλεῖν ἐμέ, / ἵνα συμποιῶ σοὔπισθεν ἐστυκὼς ἐγώ ; Hsch.~β~34 βαδίζει· περιπατεῖ ; Hsch.~τ~1065 τοίνυν· διὰ τοῦτο}
\lp{269 μὰ τὸν Ἀπόλλω οὔκ, ἤν γε μή]}{Ar.~\textit{Ach.}~59–60 μὰ τὸν Ἀπόλλω ᾽γὼ μὲν οὔ, / ἢν μὴ περὶ εἰρήνης γε πρυτανεύσητέ μοι ; Ar.~\textit{Av.}~439–440 μὰ τὸν Ἀπόλλω ᾽γὼ μὲν οὔ ; Ar.~\textit{Lys.}~208 μὰ τὴν Ἀφροδίτην οὔκ, ἐάν γε μὴ λάχῃς ; Ar.~\textit{Ran.}~508 μὰ τὸν Ἀπόλλω οὐ μή σ᾽ ἐγὼ ; Ar.~\textit{Eccl.}~155–158 ἐμοὶ μὲν οὐ δοκεῖ μὰ τὼ θεώ. / μὰ τὼ θεώ; τάλαινα ποῦ τὸν νοῦν ἔχεις; / … / μὰ Δί᾽ ἀλλ᾽ ἀνὴρ ὢν τὼ θεὼ κατώμοσας ; Ar.~\textit{Eccl.}~160 ὢ νὴ τὸν Ἀπόλλω. — παῦε τοίνυν ; Σ Ar.~\textit{Eccl.}~155 μὰ τὼ θεώ: (Ἐπειδὴ γυναικῶν ὁ ὅρκος.) παίζει. οὐ γὰρ δεῖ σε, φησίν, ὑποκρινομένην ἄνδρα γυναικεῖον ὅρκον ὀμνύναι τὰς θεάς ; Σ Ar.~\textit{Lys.}~917 μὰ τὸν Ἀπόλλω: Οὐκ ἔστιν ὅρκος γυναικῶν ; Ar.~\textit{Eq.}~14 ἵνα μὴ μάχωμαι. — μὰ τὸν Ἀπόλλω ᾽γὼ μὲν οὔ ; Ar.~\textit{Av.}~263 ὁρᾷς τιν᾽ ὄρνιν; — μὰ τὸν Ἀπόλλω ᾽γὼ μὲν οὔ ; Ar.~\textit{Lys.}~917–918 μὰ τὸν Ἀπόλλω μή σ᾽ ἐγὼ / καίπερ τοιοῦτον ὄντα κατακλινῶ χαμαί ; Hsch.~μ~413 μὰ τὼ θεώ· οὐ μόνον γυναῖκες, ἀλλὰ καὶ ἄνδρες ὤμνυον ; Phot. s.v.~μὰ τὼ θεώ (Naber I, p.~409) μὰ τὼ θεώ· γυναικεῖος ὅρκος· δυϊκῶς δὲ ὀμνύουσι τὴν κόρην καὶ τὴν Δήμητραν· ἀνδράσι δὲ οὐ πρέπει τοῦτον ὀμνύναι ; Ar.~\textit{Vesp.}~1366 οὔτοι καταπροίξει μὰ τὸν Ἀπόλλω τοῦτο δρῶν ; Ar.~\textit{Thesm.}~536 εἰ δὲ μή, ἡμεῖς ; Ar.~\textit{Thesm.}~717 μάτην λαλεῖτε· τὴν δ᾽ ἐγὼ οὐκ ἀφήσω}
\lp{270 ὀμόσῃς ἐμοί]}{Eur.~\textit{IT}~743–746 ὄμνυ· σὺ δ᾽ ἔξαρχ᾽ ὅρκον ὅστις εὐσεβής. / δώσω, λέγειν χρή, τήνδε τοῖσι σοῖς φίλοις. / τοῖς σοῖς φίλοισι γράμματ᾽ ἀποδώσω τάδε. / κἀγὼ σὲ σώσω κυανέας ἔξω πέτρας ; Ar.~\textit{Lys.}~187–190 τίν᾽ ὅρκον ὁρκώσεις ποθ᾽ ἡμᾶς; — … / εἰς ἀσπίδ᾽ ὀμόσῃς μηδὲν εἰρήνης πέρι ; Apostol.~12.98 Ὅρκοις ἄνδρας, ἀστραγάλοις δὲ παῖδας ἐξαπατῶσιν}
\lp{270 τί χρῆμα;]}{Ar.~\textit{Nub.}~816 ὦ δαιμόνιε, τί χρῆμα πάσχεις ὦ πάτερ ; Eur.~\textit{Hel.}~826 τί χρῆμα δρᾶσαι; τίν᾽ ὑπάγεις μ᾽ ἐς ἐλπίδα ; Eur.~\textit{IT}~737–738 ἦ κἀντιδώσεις τῷδε τοὺς αὐτοὺς λόγους; — τί χρῆμα δράσειν ἢ τί μὴ δράσειν; λέγε ; Eur.~\textit{Med.}~746–749 ὄμνυ πέδον Γῆς, πατέρα θ᾽ Ἥλιον πατρὸς / τοὐμοῦ, θεῶν τε συντιθεὶς ἅπαν γένος. / τί χρῆμα δράσειν ἢ τί μὴ δράσειν; λέγε ; Eur.~\textit{Hipp.}~905 ἔα, τί χρῆμα ; Eur.~\textit{Ion}~1002 τί χρῆμα; μέλλον γάρ τι προσφέρεις ἔπος ; Eur.~\textit{Hipp.}~876 τί χρῆμα; λέξον, εἴ τί μοι λόγου μέτα ; Eur.~\textit{Hipp.}~909 τί χρῆμα πάσχει; τῷ τρόπῳ διόλλυται ; Eur.~\textit{IA}~726 τί χρῆμα; πείθεσθαι γὰρ εἴθισμαι σέθεν ; Aesch.~\textit{Ag.}~1306 τί δ᾽ ἐστὶ χρῆμα; τίς σ᾽ ἀποστρέφει φόβος ; Hsch.~χ~709 χρῆμα· πρᾶγμα. πλοῦτος, οὐσία. λῆμμα}
\lp{270 συσσώσειν ἐμέ]}{Eur.~\textit{Hel.}~1388–1389 ἢν δυνώμεθα / σωθέντες αὐτοὶ καὶ σὲ συνσῶσαί ποτε ; Eur.~\textit{Cyc.}~202 ἢ ζῶντες αἶνον τὸν πάρος συσσώσομεν ; Eur.~\textit{IA}~1209 πιθοῦ. τὸ γάρ τοι τέκνα συνσῴζειν καλόν ; Pl.~\textit{Lg.}~4.707c ἵν᾽ οὕτως λέγωμεν περὶ τῶν τότε συσσωσασῶν ἡμᾶς μαχῶν}
\lp{270–271 συσσώσειν ἐμὲ / πάσαις τέχναις, ἤν μοί τι περιπίπτῃ κακόν]}{Ar.~\textit{Thesm.}~926–927 οὐ γὰρ προδώσω σ᾽ οὐδέποτ᾽, ἤνπερ ἐμπνέω, / ἢν μὴ προλίπωσ᾽ αἱ μυρίαι με μηχαναί ; Ar.~\textit{Thesm.}~765 ἄγε δὴ τίς ἔσται μηχανὴ σωτηρίας ; Eur.~\textit{Or.}~452–455 ἀλλ᾽ ἀντιλάζου καὶ πόνων ἐν τῷ μέρει, / … / ὄνομα γάρ, ἔργον δ᾽ οὐκ ἔχουσιν οἱ φίλοι / οἱ μὴ ᾽πὶ ταῖσι συμφοραῖς ὄντες φίλοι ; Eur.~\textit{Hec.}~1226–1227 ἐν τοῖς κακοῖς γὰρ ἁγαθοὶ σαφέστατοι / φίλοι ; Tosi n°~437 Amicus certus in re incerta cernitur}
\lp{271 πάσαις τέχναις]}{Ar.~\textit{Nub.}~1322–1323 ὦ γείτονες καὶ ξυγγενεῖς καὶ δημόται, / ἀμυνάθετέ μοι τυπτομένῳ πάσῃ τέχνῃ ; Ar.~\textit{Thesm.}~65 Ἀγάθωνά μοι δεῦρ᾽ ἐκκάλεσον πάσῃ τέχνῃ ; Ar.~\textit{Thesm.}~93–94 τὸ πρᾶγμα κομψὸν καὶ σφόδρ᾽ ἐκ τοῦ σοῦ τρόπου· / τοῦ γὰρ τεχνάζειν ἡμέτερος ὁ πυραμοῦς ; Ar.~\textit{Eq.}~592–593 γὰρ τοῖς ἀνδράσι τοῖσδε πάσῃ / τέχνῃ πορίσαι σε νίκην ; Ar.~\textit{Eq.}~756 νῦν δή σε πάντα δεῖ κάλων ἐξιέναι σεαυτοῦ ; Ar.~\textit{Lys.}~412–413 σὺ δ᾽ ἢν σχολάσῃς, πάσῃ τέχνῃ πρὸς ἑσπέραν / ἐλθὼν ἐκείνῃ τὴν βάλανον ἐνάρμοσον ; \textit{idem} Suid.~π~753 (praem. Πάσῃ τέχνῃ· ἐκ παντὸς τρόπου. Ἀριστοφάνης Λυσιστράτῃ ; οὐδ᾽ pro σὺ δ᾽ ; ἐλθών, ἐκείνῃ) ; Ar.~\textit{Thesm.}~87 ἀτὰρ τίν᾽ ἐκ τούτων σὺ μηχανὴν ἔχεις ; Ar.~\textit{Thesm.}~198–199 τὰς συμφορὰς γὰρ οὐχὶ τοῖς τεχνάσμασιν / φέρειν δίκαιον ἀλλὰ τοῖς παθήμασιν ; Ar.~\textit{Ran.}~957 νοεῖν ὁρᾶν ξυνιέναι στρέφειν ἐρᾶν τεχνάζειν ; Ar.~\textit{Ran.}~1235 ἀλλ᾽ ὦγάθ᾽ ἔτι καὶ νῦν ἀπόδος πάσῃ τέχνῃ ; Eur.~\textit{Heracl.}~1002–1003 πάντα κινῆσαι πέτρον, / κτείνοντα κἀκβάλλοντα καὶ τεχνώμενον ; \textit{idem} Hsch.~π~389 (πάντα κινήσας πέτρον· παροιμία) ; Eur.~\textit{IT}~1031–1032 ταῖς σαῖς ἀνίαις χρήσομαι σοφίσμασι. / δειναὶ γὰρ αἱ γυναῖκες εὑρίσκειν τέχνας ; Eur.~\textit{Med.}~278 ἐχθροὶ γὰρ ἐξιᾶσι πάντα δὴ κάλων ; Hdt.~3.152 καίτοι πάντα σοφίσματα καὶ πάσας μηχανὰς ἐπεποιήκεε ἐς αὐτοὺς Δαρεῖος ; X.~\textit{An.}~4.5.16 ἐδεῖτο αὐτῶν πάσῃ τέχνῃ καὶ μηχανῇ μὴ ἀπολείπεσθαι ; Lys.~19.11 δέομαι δ᾽ ὑμῶν πάσῃ τέχνῃ καὶ μηχανῇ μετ᾽ εὐνοίας ἀκροασαμένους ἡμῶν διὰ τέλους ; Hsch.~μ~1301 μηχαναί· τέχναι. ἐπίνοιαι. βουλεύματα. ἅμαξαι. συμπεράσματα. καὶ ὄργανα τί⟨να⟩ μηχανικά ; Hsch.~π~390 πάντα λίθον κινεῖν· παροιμία, ἧς μέμνηται Ἀρίσταρχος. ὅτε γὰρ Ξέρξης ἐπὶ τοὺς Ἕλληνας ἐστράτευσεν, ἡττηθεὶς παρὰ Σαλαμῖνι ⟨…⟩ ; Hsch.~τ~722 τέχνη· ἐπιστήμη (Γ 61 . .). ἢ δόλος (δ 455 . .). καὶ ἀτεχνῶς· ἀδόλως ; Zen.~5.62 Πάντα κάλων σεῖε: παροιμία ἐπὶ τῶν πάσῃ προθυμίᾳ χρωμένων· παρῆκται δὲ ἀπὸ τῶν τὰ ἄρμενα χαλώντων ; \textit{idem} Hsch.~π~388 (σείει ; om. παρῆκται … χαλώντων) ; Zen.~5.63 Πάντα λίθον κίνει: παροιμία. … Τὸν δὲ Ἀπόλλωνα ἀποκρίνασθαί φασι, πάντα λίθον κίνει· ὅθεν ἡ παροιμία ἐκράτησεν ; Ar.~\textit{Nub.}~885 ἐὰν δὲ μή, τὸν γοῦν ἄδικον πάσῃ τέχνῃ ; Ar.~\textit{Thesm.}~430 ἢ φαρμάκοισιν ἢ μιᾷ γέ τῳ τέχνῃ ; Ar.~\textit{Thesm.}~1130–1132 σκαιοῖσι γάρ τοι καινὰ προσφέρων σοφὰ / μάτην ἀναλίσκοις ἄν, ἀλλ᾽ ἄλλην τινὰ / τούτῳ πρέπουσαν μηχανὴν προσοιστέον ; Ar.~\textit{Eccl.}~366 Ἀντισθένη τις καλεσάτω πάσῃ τέχνῃ ; Ar.~\textit{Eccl.}~534 ἥπερ μεθῆκέ μ᾽ ἐξιέναι πάσῃ τέχνῃ ; Pl.~\textit{Prt.}~338a μήτ᾽ αὖ Πρωταγόραν πάντα κάλων ἐκτείναντα, οὐρίᾳ ἐφέντα, φεύγειν εἰς τὸ πέλαγος τῶν λόγων ἀποκρύψαντα γῆν ; Suid.~δ~1083 Δικαίως ἀδίκως: ἀντὶ τοῦ παντὶ τρόπῳ. … τουτέστι πάσῃ τέχνῃ}
\lp{271 ἤν μοί τι περιπίπτῃ κακόν]}{Ar.~\textit{Ran.}~309 οἴμοι, πόθεν μοι τὰ κακὰ ταυτὶ προσέπεσεν ; Ar.~\textit{Ran.}~969 ὃς ἢν κακοῖς που περιπέσῃ καὶ πλησίον παραστῇ ; Ar.~\textit{Thesm.}~580 σκοπῆτε καὶ τηρῆτε μὴ καὶ προσπέσῃ ; Ar.~\textit{Thesm.}~650–651 κακοδαίμων ἐγώ, / εἰς οἷ᾽ ἐμαυτὸν εἰσεκύλισα πράγματα ; Crates fr.~36~Kock (Suid.~ο~388) Ὄνος ἐν μελίτταις· Κράτης Τόλμαις. καί, ὄνος ἐν μύρῳ. παροιμία ; \textit{idem} Diogenian.~7.32 (Ὄνος ἐν μελίσσαις ἐπὶ τῶν ἐν κακοῖς περιπεσόντων) ; Eur.~\textit{Hec.}~497–498 φεῦ φεῦ· γέρων μέν εἰμ᾽, ὅμως δέ μοι θανεῖν / εἴη πρὶν αἰσχρᾷ περιπεσεῖν τύχῃ τινί ; Aesch.~\textit{Sept.}~833–834 ὦ μέλαινα καὶ τελεία / γένεος Οἰδίπου τ᾽ ἀρά, / κακόν με καρδίαν τι περιπίτνει κρύος ; Hdt.~6.106.2 μὴ περιιδεῖν πόλιν ἀρχαιοτάτην ἐν τοῖσι Ἕλλησι δουλοσύνῃ περιπεσοῦσαν πρὸς ἀνδρῶν βαρβάρων ; Philostr.~\textit{VA}~1.33 δεῖ γάρ τι περιπεσεῖν, ὃ τὸν σὸν ἐλέγξει λόγον ; Σ Pi.~\textit{P.}~4.106c μελίσσας δὲ τὰς ἱερείας, κυρίως μὲν τὰς τῆς Δήμητρος, καταχρηστικῶς δὲ καὶ τὰς πάσας, διὰ τὸ τοῦ ζῴου καθαρόν}
\lp{272 ὄμνυμι τοίνυν αἰθέρ᾽ οἴκησιν Διός]}{\textit{Test.} Suid.~π~356 Παρακεκινδυνευμένον: ἀντὶ τοῦ ὑψηλόν, ἢ βλάσφημον. Ἀριστοφάνης· … αἰθέρα, Διὸς δωμάτιον, ἢ χρόνου πόδα, … Εὐριπίδης· ὄμνυμι δ᾽ ἱερὸν αἰθέρ᾽ οἴκησιν Διός ; Ar.~\textit{Thesm.}~274 ὄμνυμι τοίνυν πάντας ἄρδην τοὺς θεούς ; Ar.~\textit{Ran.}~100–102 αἰθέρα Διὸς δωμάτιον, ἢ χρόνου πόδα, / ἢ φρένα μὲν οὐκ ἐθέλουσαν ὀμόσαι καθ᾽ ἱερῶν, / γλῶτταν δ᾽ ἐπιορκήσασαν ἰδίᾳ τῆς φρενός ; \textit{idem} Ar.~\textit{Ran.}~311 (v. 100 tantum) ; Eur.~\textit{Med.}~752 ὄμνυμι Γαῖαν ⟨Ἡλίου θ᾽ ἁγνὸν σέβας⟩ ; Eur.~fr.~487~N.² ὄμνυμι δ᾽ ἱερὸν αἰθέρ᾽, οἴκησιν Διός ; \textit{idem} Σ Ar.~\textit{Ran.}~100 (ἱρὸν ; praem. αἰθέρα Διός· Εὐριπίδου ἐκ Μελανίππης) ; Σ Ar.~\textit{Ran.}~311 αἰθέρα Διός: Ἀντὶ τοῦ τὸν Εὐριπίδην· αὐτοῦ γὰρ ὁ ἴαμβος. ἢ τὴν ἐπιθυμίαν Εὐριπίδου τοῦ ταῦτα λέγοντος ; Σ Ar.~\textit{Thesm.}~272 ἐκ σοφῆς Μελανίππης Εὐριπίδου ; \textit{idem} Σ Ar.~\textit{Lys.}~1124 (praem. ὁ στίχος) ; Ar.~\textit{Thesm.}~450–451 νῦν δ᾽ οὗτος ἐν ταῖσιν τραγῳδίαις ποιῶν / τοὺς ἄνδρας ἀναπέπεικεν οὐκ εἶναι θεούς ; Ar.~\textit{Ran.}~889–891 ἕτεροι γάρ εἰσιν οἷσιν εὔχομαι θεοῖς. — ἴδιοί τινές σοι, κόμμα καινόν; — καὶ μάλα. / ἴθι δὴ προσεύχου τοῖσιν ἰδιώταις θεοῖς ; Eur.~\textit{Hipp.}~1025–1026 νῦν δ᾽ ὅρκιόν σοι Ζῆνα καὶ πέδον χθονὸς / ὄμνυμι ; Eur.~\textit{IT}~749 ἐγὼ δ᾽ ἄνακτά γ᾽ οὐρανοῦ, σεμνὸν Δία ; Eur.~\textit{Tro.}~884–889 ὦ γῆς ὄχημα κἀπὶ γῆς ἔχων ἕδραν, / ὅστις ποτ᾽ εἶ σύ, δυστόπαστος εἰδέναι, / Ζεύς, εἴτ᾽ ἀνάγκη φύσεος εἴτε νοῦς βροτῶν, / προσηυξάμην σε· … — τί δ᾽ ἔστιν; εὐχὰς ὡς ἐκαίνισας θεῶν ; Eur.~fr.~480~N.² Ζεὺς ὅστις ὁ Ζεύς, οὐ γὰρ οἶδα πλὴν λόγῳ ; \textit{idem} Plu.~\textit{Amat.}~13,~756b–c (praem. ἀκούεις δὲ δήπου τὸν Εὐριπίδην, ὡς ἐθορυβήθη ποιησάμενος ἀρχὴν τῆς Μελανίππης ἐκείνης ; add. … ἤλλαξε τὸν στίχον ὡς νῦν γέγραπται Ζεύς, ὡς λέλεκται τῆς ἀληθείας ὕπο) ; Eur.~fr.~484.1–3~N.² κοὐκ ἐμὸς ὁ μῦθος, ἀλλ᾽ ἐμῆς μητρὸς πάρα, / ὡς οὐρανός τε γαῖά τ᾽ ἦν μορφὴ μία· / ἐπεὶ δ᾽ ἐχωρίσθησαν ἀλλήλων δίχα ; Suid.~χ~535 Χρόνου πόδα: ἀπὸ Εὐριπίδου. Ἀριστοφάνης· αἰθέρα Διὸς δωμάτιον ἢ χρόνου πόδα. ἀντὶ τοῦ Εὐριπίδου τοῦ ταῦτα λέγοντος}
\lp{272 αἰθέρ᾽]}{Ar.~\textit{Nub.}~264–265 ὦ δέσποτ᾽ ἄναξ ἀμέτρητ᾽ Ἀήρ, ὃς ἔχεις τὴν γῆν μετέωρον, / λαμπρός τ᾽ Αἰθὴρ σεμναί τε θεαὶ Νεφέλαι βροντησικέραυνοι ; Ar.~\textit{Nub.}~627 μὰ τὴν Ἀναπνοὴν μὰ τὸ Χάος μὰ τὸν Ἀέρα ; Ar.~\textit{Thesm.}~14 αἰθὴρ γὰρ ὅτε τὰ πρῶτα διεχωρίζετο ; Ar.~\textit{Ran.}~892 αἰθὴρ ἐμὸν βόσκημα καὶ γλώσσης στρόφιγξ ; Eur.~fr.~877~N.² ἀλλ᾽ αἰθὴρ τίκτει σε, κόρα, / Ζεὺς ὃς ἀνθρώποις ὀνομάζεται ; Eur.~fr.~941~N.² ὁρᾷς τὸν ὑψοῦ τόνδ᾽ ἄπειρον αἰθέρα / καὶ γῆν πέριξ ἔχονθ᾽ ὑγραῖς ἐν ἀγκάλαις; / τοῦτον νόμιζε Ζῆνα ; Eur.~fr.~985~N.² ὁ πετόμενος ἱερὸν ἀνὰ Διὸς αἰθέρα γοργοφόνος ; Ar.~\textit{Thesm.}~1065–1068a ὦ νὺξ ἱερά, / ὡς μακρὸν ἵππευμα διώκεις / ἀστεροειδέα νῶτα διφρεύουσ᾽ / αἰθέρος ἱερᾶς / τοῦ σεμνοτάτου δι᾽ Ὀλύμπου ; Eur.~\textit{Hel.}~584 αἰθήρ, ὅθεν σὺ θεοπόνητ᾽ ἔχεις λέχη ; Eur.~\textit{Ion}~1078–1079 καὶ Διὸς ἀστερωπὸς / ἀνεχόρευσεν αἰθήρ ; Eur.~fr.~839.1–3~N.² Γαῖα μεγίστη καὶ Διὸς Αἰθήρ, / ὃ μὲν ἀνθρώπων καὶ θεῶν γενέτωρ ; Aesch.~fr.~70~N.² (Clem. Al. Strom.~5.14.114) Ζεύς ἐστιν αἰθήρ, Ζεὺς δὲ γῆ, Ζεὺς δ᾽ οὐρανός ; Hsch.~α~1860 αἰθήρ· ὁ ὑπὲρ τὰ νέφη τόπος ; Ar.~\textit{Nub.}~570 Αἰθέρα σεμνότατον βιοθρέμμονα πάντων ; Ar.~\textit{Thesm.}~43,~51 ἐχέτω δὲ πνοὰς νήνεμος αἰθήρ, … — τίς ὁ φωνήσας; — νήνεμος αἰθήρ ; Eur.~\textit{Ba.}~291–293 Ζεὺς δ᾽ ἀντεμηχανήσαθ᾽ οἷα δὴ θεός. / ῥήξας μέρος τι τοῦ χθόν᾽ ἐγκυκλουμένου / αἰθέρος, ἔθηκε τόνδ᾽ ὅμηρον ἐκδιδούς ; Eur.~fr.~1023~N.² Αἰθέρα καὶ Γαῖαν πάντων γενέτειραν ἀείδω ; Hes.~\textit{Th.}~124 Νυκτὸς δ᾽ αὖτ᾽ Αἰθήρ τε καὶ Ἡμέρη ἐξεγένοντο ; Hsch.~α~1854 αἰθέριος ῥύμβος· οὐρανός (Critias VS 88 B 19, 1–2) ; Hsch.~α~1855 αἰθέρα· ἐμπυρισμόν, ἀπὸ τοῦ αἴθεσθαι. Εὐριπίδης Τρῳάσιν (1079) ; Suid.~αι~121 Αἰθήρ: ὁ ἐν ὕψει ἀήρ, ὁ ἐπάνω τοῦ ἀέρος καιόμενος ἐκ τοῦ ἡλίου ; Zen.~1.39 Αἰθέρα νήνεμον ἐρέσσειν: ἐπὶ τῶν μάτην πονούντων}
\lp{272 οἴκησιν Διός]}{\textit{Il.}~2.412 Ζεῦ κύδιστε μέγιστε κελαινεφὲς αἰθέρι ναίων ; Eur.~\textit{Hec.}~1100–1101 ἀμπτάμενος οὐράνιον / ὑψιπετὲς ἐς μέλαθρον ; Hes.~\textit{Op.}~18 θῆκε δέ μιν Κρονίδης ὑψίζυγος, αἰθέρι ναίων ; Hsch.~α~1853 αἰθέρος· ἀέρος (Β 458). ἢ ὑπὲρ τὸν ἀέρα ὄντος (Β 412) ; Hsch.~κ~956 κάσα· οἰκία, καλύβη, οἴκησις παρὰ Ῥωμαίοις}
\lp{273 τί μᾶλλον ἤ]}{Plaut.~\textit{Merc.}~486–487 Qui potius quam voles? … Qui potius quam auro expendas?}
\lp{273 Ὑποκράτους]}{Ar.~\textit{Nub.}~1001 τοῖς Ἱπποκράτους υἱέσιν εἴξεις καὶ σε καλοῦσι βλιτομάμμαν ; Eup.~fr.~112~K.-A. (=~103~Kock) εἴσ᾽ Ἱπποκράτους τε παῖδες ἐκβόλιμοί τινες, / βληχητὰ τέκνα καὶ οὐδαμῶς τοῦ ⟨–⟩ τρόπου ; \textit{idem} Suid.~ε~951 (praem. Ἐμβόλιμος· ἀτελής ; om. εἴσ᾽ ; ἐμβόλιμοι ; cetera om.) ; Ath.~3.96e μήτηρ τῶν Ἱπποκράτους υἱῶν, οὓς εἰς ὑωδίαν κωμῳδουμένους οἶδα ; Σ Ar.~\textit{Nub.}~1001 τοῖς Ἱπποκράτους : Ὑώδεις τινὲς καὶ ἀπαίδευτοι κωμῳδοῦνται, καὶ τάχα ἂν εἴησαν προκέφαλοί τινες, ὡς ἐν Γεωργοῖς καὶ ἐν Τριφάλητι. (καὶ Εὔπολις ἐν Δήμοις Ἱπποκράτους τε παῖδες ἐκβόλιμοί τινες, βληχητὰ τέκνα κοὐδαμῶς τοῦ νῦν τρόπου.) τὰ δὲ ὀνόματα αὐτῶν Τελέσιππος, Δημοφῶν, Περικλῆς. (βλιτομάμαν : Μωρόν. καὶ γὰρ τὸ βλίτον μωρὸν εἶναι δοκεῖ λάχανον.) ; \textit{idem} Suid.~τ~1135 (praem. Τοῖς Ἱπποκράτους υἱέσιν εἴξεις, καί σε καλοῦσι βλιτομάμμαν ; οὗτοι ὑώδεις … ἐκωμῳδοῦντο ; ὡς ἐν Γεωργοῖς φησι ; καὶ Εὔπολις Δήμοις· Ἱπποκράτους τε παῖδες ἐμβόλιμοί τινες, βληχητὰ τέκνα, καὶ οὐδαμῶς τοῦ τρόπου ; om. βλιτομάμαν … λάχανον), Suid.~υ~125 (Ὑώδεις· οἱ Ἱπποκράτους υἱεῖς ὑώδεις … ἐκωμῳδοῦντο· καὶ τάχ᾽ ἂν εἴησαν προκέφαλοι. καὶ Εὔπολις· εἷς. Ἱπποκράτους παῖδες ἐμβόλιμοί τινες, βληχητὰ τέκνα ; om. ὡς ἐν Γεωργοῖς … Τριφάλητι et βλιτομάμαν … λάχανον) ; Σ Ar.~\textit{Thesm.}~273 Ἱπποκράτης ἐγένετο Ἀθηναῖος, ὃς εἶχεν υἱοὺς * * διεβάλλετο δὲ σὺν αὐτοῖς ὡς ἀπαίδευτος ; Phot.~σ~845 σῦς· ὗς τοὺς Ἱπποκράτους υἱοὺς ἔλεγον καὶ τοὺς Παναιτίου καὶ Μέμνονος εἰς ὑηνίαν κωμῳδοῦντες ; Th.~4.66.3 Ἱπποκράτη τε τὸν Ἀρίφρονος ; Th.~4.101.2 καὶ Ἱπποκράτης ὁ στρατηγός ; Gal.~\textit{QAM}~4 οἱ δὲ Ἱπποκράτους υἱεῖς, οὓς ἐπὶ μωρίᾳ σκώπτουσιν οἱ κωμικοὶ}
\lp{273 Ὑποκράτους (ὗς, ὑηνία)]}{Ar.~\textit{Eq.}~986 τῆς ὑομουσίας ; Pi.~fr.~83 S.-M. (Σ \textit{Pind.\,Ol.}~6.152x ; Gal. Protr.~7) ἦν ὅτε σύας τὸ Βοιώτιον ἔθνος ἔλεγον ; Plu.~\textit{De\,esu\,carn.}~1.6,~995e τοὺς γὰρ Βοιωτοὺς ἡμᾶς οἱ Ἀττικοὶ καὶ παχεῖς καὶ ἀναισθήτους καὶ ἠλιθίους, μάλιστα διὰ τὰς ἀδηφαγίας προσαγορεύουσιν ; Σ Pi.~\textit{O.}~6.152x εἰ φεύγομεν, Βοιωτίαν ὗν: ὅτι διὰ τὴν ἀγροικίαν καὶ τὴν ἀναγωγίαν τὸ παλαιὸν οἱ Βοιωτοὶ ὗες ἐκαλοῦντο ; Σ Luc.~\textit{JTr.}~32 εἰς ἀμαθίαν ἐσκώπτοντο οἱ Βοιωτοί· ἀφ᾽ οὗ καὶ ἡ παροιμία Βοιωτία ὗς ; Hsch.~σ~2214 συηνία· ταραχή, ἀηδία, ἀπὸ τῶν συῶν (Pherecr. fr. 271 K.–A.) ; Hsch.~υ~958 ὑωδία· ἀηδία. φυρμὸς ἐν βορβόρῳ ; Hdn.~\textit{Orth.}~3.20.4.1 Ὕαντες ἔθνος περὶ Ἀλαλκομένιον. λέγονται καὶ Ὑάντειοι· καὶ Ὑάντιος καὶ Ὑαντία καὶ Ὑαντῖνοι ; Phot.~υ~36 ὑηνῶν· σκαιῶν καὶ ἀμαθῶν· καὶ ὑηνία· σκαιότης καὶ ἀμαθία. οὕτως Καλλίας (fr. 38 K.-A.) ; Suid.~β~583 Βοιωτία ὗς· ἐπὶ τῶν ἀναισθητούντων καὶ ἀπαιδεύτων ; \textit{idem} Macar.~2.79 (ἀναισθήτων), Tosi n°~1866 (cetera om.) ; Pl.~\textit{La.}~196d κατὰ τὴν παροιμίαν ἄρα τῷ ὄντι οὐκ ἂν πᾶσα ὗς γνοίη ; Σ Pi.~\textit{O.}~6.148b ἡ δὲ παροιμία ἐντεῦθεν· Ὕαντες, ἔθνος βάρβαρον ἀρχαῖον, κατῴκησαν τὴν Βοιωτίαν· οὕστινας ἐκ παραφθορᾶς τοῦ ὀνόματος, ἀντὶ τοῦ Ὕαντας λέγειν, Ὗς ἐκάλουν· … τούς τε βουλομένους εἰς ἀπαιδευσίαν ἐπισκώπτειν τινὰς, τοῦτο λέγειν· Βοιωτία ὗς ; \textit{idem} Σ Pi.~\textit{O.}~6.153c (Ὕαντες ἔθνος Βοιωτίας pro ἡ δὲ παροιμία … τὴν Βοιωτίαν ; ἐκάλουν Ὗς) ; Ar.~\textit{Pax}~928 ἵνα μὴ γένηται Θεογένους ὑηνία ; Cratin.~fr.~77~K.-A. (Σ \textit{Pind.\,Ol.}~6.152x) οὗτοι δ᾽ εἰσὶν Συοβοιωτοὶ, κρουπεζοφόρον γένος ἀνδρῶν ; Pi.~\textit{O.}~6.89–90 ἀρχαῖον ὄνειδος ἀλαθέσιν / λόγοις εἰ φεύγομεν, Βοιωτίαν ὗν ; Pl.~\textit{Tht.}~166c ὗς δὲ δὴ καὶ κυνοκεφάλους λέγων οὐ μόνον αὐτὸς ὑηνεῖς, ἀλλὰ καὶ τοὺς ἀκούοντας τοῦτο δρᾶν εἰς τὰ συγγράμματά μου ἀναπείθεις ; Clem.Al.~\textit{Strom.}~1.12.55 ἀκροατῶν τοῖς ὑώδεσί τε καὶ ἀπαιδεύτοις ; Σ Ar.~\textit{Eq.}~984 τῆς ὑομουσίας δέ, τῆς χοιρωδίας, τῆς ἀπαιδευσίας ; Hsch.~σ~2754 συοβοιωτοί· οἱ Βοιωτοὶ σύες (Cratin. fr. 77 K.–A.) ; Hdn.~\textit{Κλίσ.}~2.2 Ὕαντες δὲ λέγονται οἱ κατοικοῦντες τὴν Βοιωτίαν ; Suid.~α~2954 Ἀπαιδευτότερος Φιλωνίδου τοῦ Μελιτέως. οὗτος ὁ Φιλωνίδης οὐ μόνον μέγας ἦν, ἀλλὰ καὶ ἀμαθὴς καὶ ὑώδης. κωμῳδεῖ δὲ αὐτὸν Ἀριστοφάνης ; EM~311.36 Ἐγκτῆνες, οἱ Βοιωτοὶ οὕτως ὠνομάζοντο, ὥσπερ καὶ Ὕαντες· διὰ τὸ κτηνώδεις εἶναι καὶ χοιρώδεις ; Diogenian.~3.46 Βοιώτιον οὖς: ἐπὶ τῶν ἀναισθήτων ; Greg. Cypr.~1.91 Βοιώτιος νοῦς: ἐπὶ τῶν ἀπαιδεύτων· εἰς τοῦτο ἀεὶ ὠνειδίζοντο οἱ Βοιωτοί ; Erasm.~\textit{Adag.}~1.10.6 (n°~906) Boeotica sus ; Hor.~\textit{Ep.}~2.1.244 Boeotum in crasso iurares aere natum ; Crates fr.~6~K.-A. (Phot.~υ~296) ὗς διὰ ῥόδων· ἐπὶ τῶν σκαιῶν καὶ ἀναγώγων. Κράτης Γείτοσιν ; Plaut.~\textit{Mil.}~587 occisam saepe sapere plus multo suem ; \textit{idem} Tosi n°~1861 (praem. Sat edepol scio) ; Theoc.~5.23 ὗς ποτ᾽ Ἀθαναίαν ἔριν ἤρισεν ; \textit{idem} Apostol.~17.73 (Ἀθηναίαν ; add. Θεόκριτος. ἐπὶ τῶν τοῖς κρείττοσι φιλονεικούντων) ; Pl.~\textit{Lg.}~819d ἀλλὰ ὑηνῶν τινων εἶναι μᾶλλον θρεμμάτων ; Plu.~\textit{De\,vit.\,pud.}~18,~536a καὶ μὴ πάθος πάσχειν ὑῶδες ; Plu.~\textit{Dem.}~11.5 ἐμὲ Δημοσθένης, ἡ ὗς τὴν Ἀθηνᾶν ; Clem.Al.~\textit{Protr.}~10 (= Democr. fr.~147~D.-K.) ὑώδεις τινὲς ἄνθρωποι ; Teles fr.~1 (Περὶ αὐταρκείας) εἰ δὲ γυνὴ ὑώδης ; Σ Ar.~\textit{Plut.}~308 ἕπεσθε· παροιμία ἐπὶ τῶν ἀπαιδεύτων λεγομένη ; Hsch.~β~728 βλήχημα· μωρός, προβατώδης ; Hsch.~υ~23 ὑανία· τύρβη, μάχη, ὕβρις, ἀγερωχία (Epicharm. fr. 146, 3–4 K.–A. [θυανια codd.]) ; Hsch.~υ~959 ὑωδίνων· φύρων. βορβορῶν. μολύνων ; Phot.~υ~35 ὑηνεῖς· ὑϊκόν τι καὶ ζῳῶδες ποιεῖς ; Diogenian.~8.60 Ὗς ἐκώμασεν: ἐπὶ τῶν ἀσέμνως τι ποιούντων ; Cic.~\textit{Acad.}~1.18 nam etsi non sus Minervam, ut aiunt, tamen inepte quisquis Minervam docet ; Cic.~\textit{De\,or.}~2.233 docebo sus, ut aiunt, oratorem}
\lp{273 τὴν Ὑποκράτους ξυνοικίαν]}{Ar.~\textit{Eq.}~1001 ἐμοὶ δ᾽ ὑπερῷον καὶ ξυνοικία δύο ; Aeschin.~1.124 ὅπου μὲν γὰρ πολλοὶ μισθωσάμενοι μίαν οἴκησιν διελόμενοι ἔχουσι, συνοικίαν καλοῦμεν, ὅπου δ᾽ εἷς ἐνοικεῖ, οἰκίαν ; Aeschin.~1.125 τὴν ἐν Κολωνῷ συνοικίαν τὴν Δήμωνος καλουμένην ; Σ Ar.~\textit{Eq.}~1001 ξυνοικία δύο : Δυϊκῶς ἐπήγαγε τὰ ξυνοικία. λέγονται δὲ ξυνοικίαι καὶ αἱ μικραὶ οἰκίαι καὶ ἀποστάσεις, ἢ οὓς νῦν φανόπτας φαμέν ; Ar.~\textit{Plut.}~180 ὁ Τιμοθέου δὲ πύργος— — ἐμπέσοι γέ σοι ; Ar.~fr.~137~K.-A. (=~133~Kock) (Poll.~9.39) ἐπὶ τοῦ περιδρόμου στᾶσα τῆς συνοικίας ; Herod.~6.52–53 ὁ δ᾽ ἕτερος ἐγγὺς τῆς συνοικίης οἰκέων / τῆς Ἑρμοδώρου, τὴν πλατεῖαν ἐκβάντι ; Th.~3.74.2 ἐμπιπρᾶσι τὰς οἰκίας τὰς ἐν κύκλῳ τῆς ἀγορᾶς καὶ τὰς ξυνοικίας ; Is.~5.27 τὴν συνοικίαν ἣν εἶχεν ἀντὶ τῆς προικός ; D.~36.6 ὅσα Πασίων ἐπὶ γῇ καὶ συνοικίαις δεδανεικὼς ἦν ; D.~45.28 συνοικίαν ἑκατὸν μνῶν ; Ael.~\textit{VH}~13.12 ἐν δὲ τοῖς καὶ τὴν συνοικίαν τὴν αὑτοῦ κατέπρησεν· ἐγειτνία δὲ αὕτη τῇ Ποικίλῃ ; Hsch.~ν~120 ναύκληρος· ὁ συνοικίας προεστώς. ἢ μεμισθωμένος ὅλην καὶ ἀπομισθῶν κατὰ μέρος καλούμενος σταθμοῦχος. καὶ ὁ δεσπότης τοῦ πλοίου ; Hsch.~σ~2667 συνοίκια· τὰ πλείονα ἐφέστια ; Poll.~9.39 καὶ τὸ μὲν πλῆθος τῶν οἰκοδομημάτων οἰκίαι καὶ συνοικίαι, καὶ οἰκίας περίδρομος ; Ammon.~\textit{Diff.}~350 Nickau οἶκος μὲν γὰρ λέγεται ἡ πᾶσα οὐσία, οἰκία δὲ ἡ ὑπὸ ἑνὸς ἢ δευτέρου κατοικουμένη, συνοικία δὲ ἡ πολλοῖς καὶ μισθουμένη ; Philem.~fr.~68.3–5~K.-A. (=~65~Kock) (Stob.~4.31a.13) ἀργύριόν ἐστι· τοῦτ᾽ ἐὰν ἔχῃς, λέγε / † πρὸς τοῦτ᾽ εἰ βούλει, πάντα σοι γενήσεται, / φίλοι, βοηθοί, μάρτυρες, συνοικίαι ; Aesch.~\textit{Eum.}~916 δέξομαι Παλλάδος ξυνοικίαν ; Aesch.~\textit{Supp.}~267 δρακονθόμιλον δυσμενῆ ξυνοικίαν ; Is.~2.27 ἐπιδειξάτω ὑμῖν ὁποῖον χωρίον ἢ συνοικίαν ἢ οἰκίαν κατέλιπεν ἐκεῖνος ; Is.~6.19–21 ἣ ἐναυκλήρει συνοικίαν ἐν Πειραιεῖ αὐτοῦ καὶ παιδίσκας ἔτρεφε ; Aeschin.~1.43 εὑρίσκουσι τοῦτον ἐν συνοικίᾳ μετὰ ξένων τινῶν συναριστῶντα ; Aeschin.~1.105 ἀλλ᾽ οὐκ ἔστι τούτῳ λοιπὸν οὐκ οἰκία, οὐ συνοικία ; D.~29.3 παραδοὺς τὴν μὲν συνοικίαν Αἰσίῳ, τὸν δ᾽ ἀγρὸν Ὀνήτορι ; D.~36.34 ἐκ τίνος τρόπου πρεσβεῖα λαβὼν τὴν συνοικίαν κατὰ τὴν διαθήκην ἔχει ; D.~38.7 καὶ τὰ χωρία καὶ τὰς συνοικίας ἐπρίανθ᾽ οἱ ἐπίτροποι ; D.~53.13 τίθημι οὖν τὴν συνοικίαν ἑκκαίδεκα μνῶν ; Plu.~\textit{Alc.}~17.5 καταπρῆσαι δὲ τὴν οἰκίαν νύκτωρ ; Plu.~\textit{Nic.}~13.5 νύκτωρ ἐμπρήσαντα τὴν οἰκίαν ; Harp. s.v.~ναύκληρος ἐπὶ τῷ τὰ ἐνοίκια ἐκλέγειν ἢ οἰκίας ἢ συνοικίας ; \textit{idem} Phot.~ν~36 (praem. ναύκληρος· Ὑπερίδης οὐ μόνον ὡς ἡ συνήθεια κέχρηται τῷ ὀνόματι, ἀλλὰ καὶ ἐπὶ τοῦ μεμισθωμένου) ; Suid.~σ~1680 Συσκηνία: συνοικία}
\lp{273 τί μᾶλλον ἢ τὴν Ὑποκράτους ξυνοικίαν (serment par un objet)]}{Ar.~\textit{Nub.}~627 μὰ τὴν Ἀναπνοὴν μὰ τὸ Χάος μὰ τὸν Ἀέρα ; Ar.~\textit{Av.}~520–521 ὤμνυ τ᾽ οὐδεὶς τότ᾽ ⟨ἂν⟩ ἀνθρώπων θεόν, ἀλλ᾽ ὄρνιθας ἅπαντες· / Λάμπων δ᾽ ὄμνυσ᾽ ἔτι καὶ νυνὶ τὸν χῆν᾽, ὅταν ἐξαπατᾷ τι ; Ar.~\textit{Ran.}~100 αἰθέρα Διὸς δωμάτιον, ἢ χρόνου πόδα ; Cratin.~fr.~249~K.-A. (Phot.~ρ~17) οἷς ἦν μέγιστος ὅρκος ἅπαντι λόγῳ κύων· ἔπειτα χήν· θεοὺς δ᾽ ἐσίγων ; Pl.~\textit{Phdr.}~236e ὄμνυμι γάρ σοι—τίνα μέντοι, τίνα θεῶν; ἢ βούλει τὴν πλάτανον ταυτηνί ; Ath.~9.370a–b καὶ Ἀνάνιος δέ φησι· καὶ σὲ πολλὸν ἀνθρώπων / ἐγὼ φιλέω μάλιστα, ναὶ μὰ τὴν κράμβην. καὶ Τηλεκλείδης Πρυτάνεσιν «ναὶ μὰ τὰς κράμβας» ἔφη. καὶ Ἐπίχαρμος ἐν Γᾷ καὶ Θαλάσσᾳ «ναὶ μὰ τὰν κράμβαν». Εὔπολις Βάπταις «ναὶ μὰ τὴν κράμβην» ; Σ Ar.~\textit{Av.}~521 Λάμπων δ᾽ ὄμνυσι : (Τῶν εἰκῇ δαιμόνων. ὅτι πρῶτοι οἱ Σωκρατικοὶ ἐπετήδευσαν οὕτως ὀμνύναι. Σωσικράτης γὰρ ἐν τῷ β΄ τῶν Κρητικῶν οὕτως φησὶ « Ῥαδάμανθυς δὲ δοκεῖ διαδεξάμενος τὴν βασίλειαν δικαιότατος γεγενῆσθαι πάντων ἀνθρώπων. λέγεται δὲ αὐτὸν πρῶτον οὐδένα ἐᾶν ὅρκους ποιεῖσθαι κατὰ τῶν θεῶν, ἀλλ᾽ ὀμνύναι κελεῦσαι χῆνα καὶ κύνα καὶ κριόν, καὶ τὰ ὅμοια. » ; \textit{idem} Suid.~λ~93 (Λάμπων ὄμνυσι τὸν χῆν᾽, ὅταν ἐξαπατᾷ τινά· πρῶτοι Σωκρατικοὶ … ; Ῥαδάμανθυς δὲ πρῶτος οὐκ εἴα ὅρκους ποιεῖσθαι κατὰ θεῶν, ἀλλ᾽ ὀμνύναι χῆνα καὶ κύνα καὶ κριὸν καὶ τὰ ὅμοια ; add. … εἴρηται οὖν ἐπὶ τῶν μετὰ ὅρκου ἐξαπατώντων τινά) ; Hsch.~χ~398 χῆνα ὀμνύειν· ἔθος ἦν τὸ κατὰ τῶν τοιούτων ὀμνύειν τοῖς παλαιοῖς (Cratin. fr. 249, 2 K.–A. Ar. Av. 521) ; Phot.~ρ~17 Ῥαδαμάνθυος ὅρκος· ὁ κατὰ χηνὸς ἢ κυνὸς ἢ πλατάνου ἢ κριοῦ ἤ τινος ἄλλου τοιούτου· … τοιοῦτοι δὲ καὶ οἱ Σωκράτους ὅρκοι ; \textit{idem} Suid.~ρ~13 (add. σημεῖον τοῦτο περὶ ὅρκου, ὃ καὶ νῦν ἐπιχωριάζει· πολλοὶ γὰρ πρὸς λάχανα ὀμνύουσιν, εὐορκίαν ἐπιδεικνύμενοι), Apostol.~15.17 (praem. ἐπὶ τῶν ἐπὶ δικαιοσύνῃ μαρτυρουμένων. ὁ δὲ ὅρκος ἦν ; om. θεοὺς δ᾽ ἐσίγων. Κρατῖνος Χείρωσι) ; Zen.~5.81 Ῥαδαμάνθυος ὅρκος: Κρατῖνός φησιν, ἐπὶ τῷ χηνὶ καὶ τῷ κυνὶ καὶ τοῖς τοιούτοις ὅρκον Ῥαδαμάνθυϊ ἀνατιθέασιν· ὡς καὶ Σωσικράτης ἐν δευτέρῳ Κρητικᾶν· ἵνα μὴ θεοὺς ὀμνύωσιν ; \textit{idem} Hsch.~ρ~14 (τὸν ἐπὶ χηνὶ καὶ κυνὶ καὶ τοῖς τοιούτοις ὅρκον ; cetera om.) ; Macar.~7.49 Ῥαδαμάνθυος ὅρκος: ἐπὶ τῶν κατὰ μικροῦ τινος ὀμνύντων, ἤγουν κριοῦ, κύκνου, λαχάνων καὶ τῶν ὁμοίων ; App. Prov.~2.91 Εὐριπίδειος ὅρκος: ἴσως ὁ κατὰ κυνὸς ἢ χηνός ; \textit{idem} Macar.~4.19 (τὰ κατὰ χηνὸς καὶ κυνός) ; Ar.~\textit{Nub.}~814 οὔτοι μὰ τὴν Ὁμίχλην ἔτ᾽ ἐνταυθοῖ μενεῖς ; Ar.~\textit{Nub.}~1240–1241 θαυμασίως ἥσθην θεοῖς, / καὶ Ζεὺς γέλοιος ὀμνύμενος τοῖς εἰδόσιν ; Pl.~\textit{Ap.}~22a καὶ νὴ τὸν κύνα, ὦ ἄνδρες Ἀθηναῖοι ; Pl.~\textit{Grg.}~482b μὰ τὸν κύνα τὸν Αἰγυπτίων θεόν ; Luc.~\textit{Vit.Auct.}~16 καὶ μὴν ὀμνύω γέ σοι τὸν κύνα καὶ τὴν πλάτανον οὕτω ταῦτα ἔχειν. Ἡράκλεις τῆς ἀτοπίας τῶν θεῶν ; Philostr.~\textit{VA}~6.19 ἐγένετό τις ἔφη Σωκράτης Ἀθηναῖος ἀνόητος, ὥσπερ ἡμεῖς, γέρων, ὃς τὸν κύνα καὶ τὸν χῆνα καὶ τὴν πλάτανον θεούς τε ἡγεῖτο καὶ ὤμνυ ; App. Prov.~3.100 Ναὶ ναὶ μὰ μήκωνος χλόην: ἐπὶ χλευασμοῦ ὁ ὅρκος ; Porph.~\textit{Abst.}~3.16 Σωκράτης δὲ καὶ ὤμνυεν κατ᾽ αὐτῶν, καὶ ἔτι πρὸ αὐτοῦ Ῥαδάμανθυς}
\lp{274 ὄμνυμι τοίνυν … τοὺς θεούς]}{Ar.~\textit{Ran.}~889–892 ἕτεροι γάρ εἰσιν οἷσιν εὔχομαι θεοῖς. / ἴδιοί τινές σοι, κόμμα καινόν; … αἰθὴρ ἐμὸν βόσκημα καὶ γλώσσης στρόφιγξ ; Eur.~\textit{Hipp.}~1025–1026 νῦν δ᾽ ὅρκιόν σοι Ζῆνα καὶ πέδον χθονὸς / ὄμνυμι τῶν σῶν μήποθ᾽ ἅψασθαι γάμων ; Eur.~\textit{Med.}~746–747 ὄμνυ πέδον Γῆς, πατέρα θ᾽ Ἥλιον πατρὸς / τοὐμοῦ, θεῶν τε συντιθεὶς ἅπαν γένος ; Eur.~\textit{Med.}~752–753 ὄμνυμι Γαῖαν ⟨Ἡλίου θ᾽ ἁγνὸν σέβας⟩ / θεούς τε πάντας ἐμμενεῖν ἅ σου κλύω ; \textit{Il.}~14.278–279 ὄμνυε δ᾽ ὡς ἐκέλευε, θεοὺς δ᾽ ὀνόμηνεν ἅπαντας / τοὺς ὑποταρταρίους ; Hp.~\textit{Jusj.}~1 ὄμνυμι Ἀπόλλωνα ἰητρὸν καὶ Ἀσκληπιὸν καὶ Ὑγείαν καὶ Πανάκειαν καὶ θεοὺς πάντας τε καὶ πάσας ; D.~54.41 καὶ νῦν ὀμνύω τοὺς θεοὺς καὶ τὰς θεὰς ἅπαντας καὶ πάσας ; Hsch.~α~7098 Ἀρδήττους· τοὺς ῥᾳδίως ἐπὶ τοὺς ὅρκους ἰόντας Ἀρδήττους ἐκάλουν. τόπος περὶ τὸν Ἰλισσόν, ἐγγὺς τοῦ Παναθηναϊκοῦ σταδίου ; Ar.~\textit{Nub.}~247–248 ποίους θεοὺς ὀμεῖ σύ; πρῶτον γὰρ θεοὶ / ἡμῖν νόμισμ᾽ οὐκ ἔστι ; Ar.~\textit{Thesm.}~331–334 εὔχεσθε τοῖς θεοῖσι τοῖς Ὀλυμπίοις / … / καὶ ταῖσι Δηλίαισι τοῖς ⟨τ᾽⟩ ἄλλοις θεοῖς ; Aesch.~\textit{Ag.}~513–514 τούς τ᾽ ἀγωνίους θεοὺς / πάντας προσαυδῶ ; Harp. s.v.~Ἀρδηττός ἈΡΔΗΤΤΟΣ … τόπος Ἀθήνησιν ὑπὲρ τὸ στάδιον τὸ Παναθηναϊκὸν … ἐν τούτῳ, φασὶ, δημοσίᾳ πάντες ὤμνυον Ἀθηναῖοι τὸν ὅρκον ; Suid.~α~3807 καὶ Ἀρδηττὸς τόπος ἐστὶν Ἀθήνησιν, ἐν ᾧ πάντες Ἀθηναῖοι δημοσίᾳ ὤμνυον τὸν ὅρκον τὸν Ἡλιαστικόν ; Ar.~\textit{Nub.}~1239 οὔ τοι μὰ τὸν Δία τὸν μέγαν καὶ τοὺς θεοὺς ; Ar.~\textit{Vesp.}~1136 τουτὶ τὸ κακὸν τί ἐστι πρὸς πάντων θεῶν ; Ar.~\textit{Av.}~865 καὶ ὄρνισιν Ὀλυμπίοις καὶ Ὀλυμπίῃσι πᾶσι καὶ πάσῃσιν ; X.~\textit{An.}~6.6.17 ἐγώ, ὦ ἄνδρες, ὄμνυμι θεοὺς καὶ θεὰς ; X.~\textit{Cyr.}~5.4.31 ὄμνυμί σοι θεούς, οἳ καὶ ὁρῶσι πάντα καὶ ἀκούουσι πάντα}
\lp{274 πάντας ἄρδην]}{\textit{Test.} Hsch.~α~7097 ἄρδην· φοράδην. ἢ σφόδρα. ἢ παντελῶς. ἢ ἀθρόως. ἢ ἐξ ὁλοκλήρου (Eur. Hec. 887) ; Eur.~\textit{Phoen.}~1146 τί μέλλετ᾽ ἄρδην πάντες ἐμπίπτειν πύλαις ; D.~19.61 ἀλλ᾽ ἐκ τοῦ σπείσασθαι πάντας ἄρδην ἀπολέσθαι ; Σ Ar.~\textit{Thesm.}~274 ὅτι κατὰ δύναμιν τῆς λέξεως, ἀντὶ τοῦ ἀθρόους πάντας. τὸ γὰρ ἄρδην ὅμοιόν ἐστι τῷ φοράδην ἐνέγκαι· ἡ γὰρ ἄρσις τοῦτο δηλοῖ. εἰ μὴ νῦν τοῦτο βούλεται, πάντας ὁμοῦ αἴρω ; \textit{idem} Suid.~α~3806 (praem. Ἄρδην ; om. πάντας ; ἀντὶ τοῦ πάντας ὁμοῦ αἴρω pro εἰ μὴ νῦν τοῦτο βούλεται, πάντας ὁμοῦ αἴρω) ; Eust.~707.37–39 Τὸ μέντοι ἄρδην ἐπίῤῥημα οὐκ ἂν οὐδ᾽ αὐτὸ παραγάγοι τινὰ ὡς ἐκ τοῦ ἄρδω γίνεται, ἀλλὰ μάλιστα ἐκ τοῦ αἴρειν, ὥσπερ ἐκ τοῦ ἀείρειν γίνεται τὸ ἀέρδην. ἔστι γάρ, φασι, τὸ ἄρδην ἀντὶ τοῦ φοράδην ; Ael.Dion. α~169 Erbse (= fr.~66 Schwabe) ἄρδην· φοράδην. Εὐριπίδης Ἀλκήστιδι (607 s.) νέκυν … ἤδη … πρόσπολοι φέρουσιν ἄρδην ἐς τάφον τε καὶ πυράν ; Eur.~\textit{Hec.}~887 καὶ Λῆμνον ἄρδην ἀρσένων ἐξῴκισαν ; Aesch.~\textit{PV}~1051 Τάρταρον ἄρδην ῥίψειε δέμας ; Pl.~\textit{R.}~421a ὁρᾷς δὴ ὅτι πᾶσαν ἄρδην πόλιν ἀπολλύασιν ; Σ Eur.~\textit{Phoen.}~1146 ἄρδην : ὁμοῦ ; Hsch.~α~1383 ἀέρδην· ἄνω ἢ φοράδην (Aesch. Ag. 234) ; Hsch.~κ~1191 ⟨κατ᾽ ἄκρων·⟩ ἄρδην. φοράδην. παντελῶς. σφόδρα (Eur. Phoen. 1176) ; Hsch.~π~1080 πασσύρως· ἄρδην. πανοικί ; Zen.~2.31 Αὐτῷ κανῷ: αὕτη τάττεται κατὰ τῶν ἄρδην ὁτιοῦν λυμαινομένων ; Diph.~fr.~105.2~K.-A. (=~105~Kock) (Stob.~4.32b.27) ἄρδην ἀνεῖλε καὶ κατέστρεψεν βίον ; Eur.~\textit{Alc.}~608 φέρουσιν ἄρδην πρὸς τάφον τε καὶ πυράν ; Eur.~\textit{Ion}~1274 ἄρδην ἂν ἐξέπεμψας εἰς Ἅιδου δόμους ; Eur.~\textit{Phoen.}~1620 πόθεν; τί μ᾽ ἄρδην ὧδ᾽ ἀποκτείνεις, Κρέον ; Soph.~\textit{Aj.}~1279 πηδῶντος ἄρδην Ἕκτορος τάφρων ὕπερ ; Soph.~\textit{Ant.}~430 ἔκ τ᾽ εὐκροτήτου χαλκέας ἄρδην πρόχου ; Pl.~\textit{Lg.}~677c θῶμεν δὴ τὰς ἐν τοῖς πεδίοις πόλεις καὶ πρὸς θαλάττῃ κατοικούσας ἄρδην ἐν τῷ τότε χρόνῳ διαφθείρεσθαι ; Pl.~\textit{Lg.}~708b ἤδη δέ ποτε καὶ συνάπασα πόλις τινῶν ἔφυγεν, ἄρδην κρείττονι κρατηθεῖσα πολέμῳ ; Pl.~\textit{Lg.}~716b ἑαυτόν τε καὶ οἶκον καὶ πόλιν ἄρδην ἀνάστατον ἐποίησεν ; Isoc.~14.19 τοὺς δ᾽ ἄρδην ἀπολλύοντες ; Aeschin.~3.93 τὰς δὲ συνεδρίας καὶ τὰς συντάξεις, ἐξ ὧν ἰσχύσειν ὁ πόλεμος ἔμελλεν, ἄρδην ἀπέδοτο ; D.~27.26 ἄρδην ὅλον τὸ ἐργαστήριον ἀφανίζουσι ; Euseb.~\textit{HE}~5.24.9 ἀκοινωνήτους πάντας ἄρδην τοὺς ἐκεῖσε ἀνακηρύττων ἀδελφούς ; Σ Eur.~\textit{Rhes.}~58 σύρδην: ὁμοῦ, ἄρδην ; Hsch.~φ~916 φροῦδον· ἄρδην, παντελῶς ἀφανές, οὐ κεχωρηκός, ἔρημον ; A.D. Adv. p.~198 Schneider (GG II.1) καὶ παρὰ τὴν πέμπτην δὲ συζυγίαν, αἴρω ἄρδην, φύρω φύρδην ; Σ Eur.~\textit{Phoen.}~1377 ὅθεν παροιμία ἐπὶ τῶν ἄρδην ἀπολλυμένων οὐδὲ πυρφόρος ἐσώθη}
\lp{275–276 ἡ φρὴν ὤμοσεν, ἡ γλῶττα δ᾽ οὐκ ὀμώμοκ᾽]}{\textit{Test.} Eust.~443.35–37 Τὸ δέ, ἑκὼν ἀέκοντί γε θυμῷ, ἔπεσε μὲν εἰς κοινὴν παροιμίαν τήν, ἑκὼν ἀέκων τόδε τι πεποίηκα, δι᾽ ἧς βοηθοῖτ᾽ ἂν καὶ τό, ἡ γλῶσσ᾽ ὀμώμοχ᾽, ἡ δὲ φρὴν ἀνώμοτος, ὃ δὴ πικρῶς λογιστεύει ὁ κωμικός ; Ar.~\textit{Ach.}~396–400 οὐκ ἔνδον ἔνδον ἐστίν, εἰ γνώμην ἔχεις. / πῶς ἔνδον εἶτ᾽ οὐκ ἔνδον; … ὁ νοῦς μὲν ἔξω ξυλλέγων ἐπύλλια / οὐκ ἔνδον, αὐτὸς δ᾽ ἔνδον ἀναβάδην ποιεῖ / τραγῳδίαν ; Ar.~\textit{Ran.}~101–102 ἢ φρένα μὲν οὐκ ἐθέλουσαν ὀμόσαι καθ᾽ ἱερῶν, / γλῶτταν δ᾽ ἐπιορκήσασαν ἰδίᾳ τῆς φρενός ; Ar.~\textit{Ran.}~1471 ἡ γλῶττ᾽ ὀμώμοκ᾽, Αἰσχύλον δ᾽ αἱρήσομαι ; Eur.~\textit{Hipp.}~612 ἡ γλῶσσ᾽ ὀμώμοχ᾽, ἡ δὲ φρὴν ἀνώμοτος ; \textit{idem} Σ Ar.~\textit{Ran.}~102 (praem. Παρὰ τὰ ἐξ Ἱππολύτου Εὐριπίδου), Σ Ar.~\textit{Ran.}~1471 (praem. Παρὰ τὰ ἐξ Ἱππολύτου ; γλῶττ᾽), Σ Soph.~\textit{OC}~230 (praem. περὶ τὴν αὐτὴν ὑπόθεσιν καὶ τὸ ; add. καὶ ἐκεῖνος γὰρ ἀπατηθεὶς ὤμοσεν), Σ Il.~19.132 (praem. ἐπαινετέον δὲ τὸ ἀψευδὲς Διός· οὐ γὰρ εἶπεν ὡς Εὐριπίδης ; ὀμώμοκεν ; cetera om.), Just.~\textit{Apol.}~1.39 (praem. Δυνατὸν γὰρ ἦν τὸ λεγόμενον ; ὀμώμοκεν ; add. ποιεῖν ἡμᾶς εἰς τοῦτο), Basil.~\textit{Leg.\,lib.\,gent.}~6.5 (praem. ἀλλ᾽ ; γλῶττα μὲν ὀμώμοκεν ; add. κατ᾽ Εὐριπίδην ἐρεῖ), Max.Tyr.~40.6 (γλῶττα ἐπώμοσεν), Philop.~\textit{in\,de\,An.}~3 (CAG XV) (γλῶττ᾽), Suid.~α~2592 (praem. Ἀνώμοτος· χωρὶς ὅρκου. ὁ δὲ ἀνὴρ καὶ ἀνώμοτος πιστεύειν ἄξιος. καὶ αὖθις), Suid.~η~63 (γλῶττ᾽), Suid.~μ~1342 (praem. Μῶρος· παρὰ Ἀττικοῖς προπερισπᾶται. ὦ μῶροι, οὐκ ἴστε τὸν σοφώτατον τῆς τραγῳδίας εἰπόντα ; γλῶττ᾽ ὀμώμοκεν), Tosi n°~2252, Erasm.~\textit{Adag.}~2.5.41 (Lingua juravit. Illud è Tragœdia cessit in proverbium· ἡ γλῶσσ᾽ ὀμώμοχεν ; cetera om.) ; Eur.~\textit{Hipp.}~656–658 εὖ δ᾽ ἴσθι, τοὐμόν σ᾽ εὐσεβὲς σῴζει, γύναι· / εἰ μὴ γὰρ ὅρκοις θεῶν ἄφαρκτος ᾑρέθην, / οὐκ ἄν ποτ᾽ ἔσχον μὴ οὐ τάδ᾽ ἐξειπεῖν πατρί ; Eur.~\textit{Hipp.}~1062–1063 οὐ δῆτα· πάντως οὐ πίθοιμ᾽ ἂν οὕς με δεῖ, / μάτην δ᾽ ἂν ὅρκους συγχέαιμ᾽ οὓς ὤμοσα ; \textit{Il.}~9.312–313 ἐχθρὸς γάρ μοι κεῖνος ὁμῶς Ἀΐδαο πύλῃσιν / ὅς χ᾽ ἕτερον μὲν κεύθῃ ἐνὶ φρεσίν, ἄλλο δὲ εἴπῃ ; Pi.~\textit{I.}~6.72 γλῶσσα δ᾽ οὐκ ἔξω φρενῶν ; Ov.~\textit{Her.}~21.135–138 Quae iurat, mens est. sed nil iuravimus illa; / illa fidem dictis addere sola potest. / Consilium prudensque animi sententia iurat, / et nisi iudicii vincula nulla valent ; Pl.~\textit{Smp.}~199a ἡ γλῶσσα οὖν ὑπέσχετο, ἡ δὲ φρὴν οὔ· χαιρέτω δή ; Pl.~\textit{Tht.}~154d Εὐριπίδειόν τι συμβήσεται· ἡ μὲν γὰρ γλῶττα ἀνέλεγκτος ἡμῖν ἔσται, ἡ δὲ φρὴν οὐκ ἀνέλεγκτος ; Arist.~\textit{Rh.}~3.15.8 ὥσπερ Εὐριπίδης πρὸς Ὑγιαίνοντα ἐν τῇ ἀντιδόσει κατηγοροῦντα ὡς ἀσεβής, ὅς γ᾽ ἐποίησε κελεύων ἐπιορκεῖν, ἡ γλῶσσ᾽ ὀμώμοχ᾽, ἡ δὲ φρὴν ἀνώμοτος ; Luc.~\textit{Vit.Auct.}~9 τὸ Εὐριπίδειον ἐκεῖνο ζηλώσεις μικρὸν ἐναλλάξας. τὸ ποῖον; ἡ φρήν σοι ἀλγήσει, ἡ δὲ γλῶσσα ἔσται ἀνάλγητος ; Cic.~\textit{Off.}~3.108 Scite enim Euripides: Iurávi lingua, méntem iniuratám gero ; Σ Ar.~\textit{Ach.}~398 ὁ νοῦς μὲν ἔξω συλλέγων : (Σκώπτει πάλιν τὸν Εὐριπίδην διὰ τὸ ἐν τοῖς λόγοις εἶναι συλλογιστικόν, καὶ οὗ ἂν λέγῃ τὸ ἐναντίον πάλιν κατασκευάζοντα, οἷον [Hippol. 612] ἡ γλῶσσ᾽ ὀμώμοχ᾽, ἡ δὲ φρὴν ἀνώμοτος.) οἷον, αὐτὸς μὲν ἔσω ἐστίν, ὁ δὲ νοῦς αὐτοῦ συλλογίζεται τῶν ἔξω τι καὶ μετεωρίζεται ; \textit{idem} Suid.~ο~868 (σκώπτει τὸν Εὐριπίδην ὡς συλλογιστικόν ; ἐὰν ; ὀμώμοκεν ; om. οἷον, αὐτὸς μὲν … μετεωρίζεται) ; Σ Ar.~\textit{Thesm.}~276 ἐπειδὴ ἐν Ἱππολύτῳ τοῦτο εἶπεν Εὐριπίδης, νῦν εὐλαβούμενος λέγει ὅτι ἡ φρήν σου ὀμώμοκε καὶ οὐχ ἡ γλῶττα ; Σ Eur.~\textit{Hipp.}~612 οὐ καθολικῶς δὲ τὸ τοιοῦτον, ἀλλὰ νῦν … ὤμοσε, φησὶν, ἡ γλῶσσα προπετέστερον καὶ παρὰ γνώμην τῆς φρενός. Ἀριστοφάνης καθολικώτερον νοήσας φησὶ τὸν Εὐριπίδην ὑβρίζειν ⟨διὰ⟩ τῆς γλώσσης τὴν διάνοιαν ; Suid.~π~356 Παρακεκινδυνευμένον: ἀντὶ τοῦ ὑψηλόν, ἢ βλάσφημον. … ὅστις φθέγξεται τοιουτονί τι παρακεκινδυνευμένον, αἰθέρα, Διὸς δωμάτιον, ἢ χρόνου πόδα, ἢ φρένα μὲν οὐκ ἐθέλουσαν ὀμόσαι καθ᾽ ἱερῶν, γλῶτταν δ᾽ ἐπιορκήσασαν ἄνευ τῆς φρενός. Εὐριπίδης· ὄμνυμι δ᾽ ἱερὸν αἰθέρ᾽ οἴκησιν Διός. Ἀλεξάνδρου δέ· καὶ χρόνου πρόβαινε πούς. καὶ ἐξ Ἱππολύτου Εὐριπίδου· ἡ γλῶσσ᾽ ὀμώμοχ᾽, ἡ δὲ φρὴν ἀνώμοτος ; App. Prov.~2.100 Ἡ γλῶττ᾽ ἀνέγνωχ᾽, ἡ δὲ φρὴν οὐ μανθάνει: ἐπὶ τῶν ἀμαθῶν ; \textit{idem} Macar.~4.36 ; Eur.~\textit{Hipp.}~611 ὦ τέκνον, ὅρκους μηδαμῶς ἀτιμάσῃς ; Eur.~\textit{Med.}~737 λόγοις δὲ συμβὰς καὶ θεῶν ἀνώμοτος ; Liv.~1.58.9 mentem peccare, non corpus, et unde consilium afuerit, culpam abesse ; Them.~\textit{in\,APo.} (CAG V~1) ἀλλ᾽ ἡ γλῶττα τούτου κατ᾽ Εὐριπίδην φιλονεικεῖ ; Anon. in Rh. (fr. comm.)~15 (CAG XXI~2) κατηγόρησε τοῦ Εὐριπίδου οὗτος ὁ Ὑγιαίνων ὡς ἀσεβοῦς γράψαντος ἐν τῷ δράματι Ἱππολύτου τὸ ἡ γλῶσσ᾽ ὀμώμοχ᾽, ἡ δὲ φρήν ; Hsch.~φ~866 φρενί· διανοίᾳ (ζ 65)}
\lp{275 μέμνησο … ὅτι]}{Ar.~\textit{Lys.}~931–932 μέμνησό νυν· / μή μ᾽ ἐξαπατήσῃς τὰ περὶ τῶν διαλλαγῶν ; Ar.~\textit{Eq.}~495 καὶ σπεῦδε ταχέως. — ταῦτα δρῶ. — μέμνησό νυν ; Ar.~\textit{Thesm.}~1134 μέμνησο Περσεῦ μ᾽ ὡς καταλείπεις ἀθλίαν ; Epich.~fr.~218~K.-A. (Polyb.~18.40.4) νᾶφε καὶ μέμνασ᾽ ἀπιστεῖν· ἄρθρα ταῦτα τᾶν φρενῶν ; \textit{idem} Cic.~\textit{Att.}~1.19.8 (praem. ut crebro mihi vafer ille Siculus insusurret Epicharmus cantilenam illam suam) ; Ar.~\textit{Nub.}~887–888 τοῦτό νυν μέμνησ᾽, ὅπως / πρὸς πάντα τὰ δίκαι᾽ ἀντιλέγειν δυνήσεται ; Ar.~\textit{Nub.}~1107 δίδασκε καὶ κόλαζε καὶ μέμνησ᾽ ὅπως ; Ar.~\textit{Ran.}~1520 μέμνησο δ᾽ ὅπως ὁ πανοῦργος ἀνὴρ ; Hsch.~μ~808 μέμνησο· μνημόνευε (Eur. Or. 125 . .)}
\lp{276 οὐδ᾽ ὥρκωσ᾽ ἐγώ]}{Cratin.~fr.~401~K.-A. (Phryn.~\textit{Ecl.}~337 F.) ὥρκωσε καὶ ὁρκωτὴς δ᾽ ἐγώ ; \textit{idem} Phryn.~\textit{Ecl.}~337 F. (338 Lobeck) (add. οὕτω Κρατῖνος φησιν. μᾶλλον δὲ διὰ τοῦ ω λέγε ἢ διὰ τοῦ ι ὥρκισεν) ; Hsch.~ο~1237 ὁρκωτής· ὁ ὁρκίζων ; Phot.~ο~485 ὁρκίζειν καὶ ὁρκοῦν· ἑκατέρως ; Zen.~3.80 Ἐπακτὸς ὅρκος: οὐκ ἐγχώριος, ἀλλ᾽ ἀπὸ ξένης ἐπηγμένος· ἢ ὃν ἕτερος ἐπάγει, οὐκ αὐθαίρετος ; \textit{idem} Diogenian.~4.60 (om. οὐκ ἐγχώριος … ἢ ; add. οἷον βίαιος) ; Ar.~\textit{Lys.}~187 τίν᾽ ὅρκον ὁρκώσεις ποθ᾽ ἡμᾶς ; Hsch.~α~8725 ἄφορκος· θρασύς. ἄπιστος ; Hsch.~ο~1231 ὁρκίδδει· ὅρκον ποιεῖ ὀμνύειν ; Hsch.~ο~1233 ὅρκοι· δεσμοί, σφραγίδες ; Poll.~1.38–39 ὁρκῶσαι, ὁρκῶ … ὁρκοῦν, ἐξορκοῦν, ἐνόρκως τι εἰπεῖν … ἔνορκον, ἐνώμοτον. ἄνορκον, ἀνώμοτον ; Harp. s.v.~ἐπακτὸς ὅρκος ὃν αὐτός τις ἑκὼν αὑτῷ ἐπάγεται, τουτέστιν αἱρεῖται· ἄλλοι δὲ τοὐναντίον ὁ ἀλλαχόθεν ἐπιφερόμενος, ἀλλ᾽ οὐκ αὐθαίρετος ; Th.~4.74.2 ὁρκώσαντες πίστεσι μεγάλαις ; Th.~8.75.2 ὥρκωσαν πάντας τοὺς στρατιώτας τοὺς μεγίστους ὅρκους ; Isoc.~1.23 ὅρκον ἐπακτὸν προσδέχου διὰ δύο προφάσεις, ἢ σεαυτὸν αἰτίας αἰσχρᾶς ἀπολύων, ἢ φίλους ἐκ μεγάλων κινδύνων διασώζων ; D.~19.151 ἢ γὰρ παρόντων καὶ κατὰ τὸ ψήφισμ᾽ αὐτὸν ἐξορκωσάντων ; D.~19.278 οὐ τὸ μὲν ψήφισμα τοὺς ἄρχοντας ὁρκοῦν τοὺς ἐν ταῖς πόλεσιν, οὗτοι δ᾽ … τούτους ὥρκισαν ; Lys.~20.26 ἀφικομένου γὰρ ἐκεῖσε Συρακουσίου ὅρκιον ἔχοντος καὶ ἑτοίμου ὄντος ὁρκοῦν ; Hsch.~ο~1229 ὅρκια· ⟨σκῆπτρον⟩ καθ᾽ οὗ ὤμνυον. ἤτοι ὁρκωμόσια}
\lp{276 ὀμώμοκ᾽]}{Phryn.~\textit{Ecl.}~25 Fischer Ὦμοκε τελέως ἀηθές· χρὴ γὰρ ὀμώμοκε λέγειν ; Archipp.~fr.~36~K.-A. (Phot.~ο~330) ὀμώμοκα· Ἄρχιππος Ὄνῳ}
\lp{275–276 ὤμοσεν … ὀμώμοκ᾽ (serment sans valeur)]}{Hsch.~α~8771 ἀφροδίσιος ὅρκος· παροιμία, ἣν καὶ ἀναγράφουσιν· Ἀφροδίσιος ὅρκος οὐ δάκνει· πρῶτος δὲ Ἡσίοδος (fr. 124 M.–W.) ἔπλασε τὰ περὶ τὸν Δία καὶ τὴν Ἰώ ; Diogenian.~3.37 Ἀφροδίσιος ὅρκος οὐκ ἐμποίνιμος: ἐπὶ τῶν ἐρωτικῶν ; \textit{idem} Greg. Cypr.~1.2 (om. ἐπὶ τῶν ἐρωτικῶν), Suid.~α~4652 (praem. καὶ παροιμία ; ἐπὶ τῶν δι᾽ ἔρωτα ὀμνυόντων πολλάκις καὶ ἐπιορκούντων), Tosi n°~332 (praem. Iuppiter ex alto periuria ridet amantum … ; om. ἐπὶ τῶν ἐρωτικῶν) ; Soph.~fr.~811~R. (Hellad. ap. Phot. Bibl.~530a15) ὅρκους ἐγὼ γυναικὸς εἰς ὕδωρ γράφω ; \textit{idem} Macar.~6.48 (χρὴ γράφειν pro ἐγὼ … γράφω ; add. Σοφοκλέους αὕτη ἐστίν), Tosi n°~337 (add. … ὅρκον δ᾽ ἐγὼ γυναικὸς εἰς οἶνον γράφω) ; Pl.~\textit{Smp.}~183b ὃ δὲ δεινότατον, ὥς γε λέγουσιν οἱ πολλοί, ὅτι καὶ ὀμνύντι μόνῳ συγγνώμη παρὰ θεῶν ἐκβάντι τῶν ὅρκων—ἀφροδίσιον γὰρ ὅρκον οὔ φασιν εἶναι ; App. Prov.~4.33 Ὅρκος Ἀφροδίσιος συγγινώσκεται: παρόσον οἱ ἐρῶντες εὐχερῶς ὀμνύουσι ; \textit{idem} Macar.~6.47 (Ἀφροδίτης ; ἐπιορκοῦσιν) ; Suid.~τ~201 Ταχυβάμονας ὅρκους: Ἀρίσταρχος κατ᾽ ἀντίφρασιν ἀκούει, ἀντὶ τοῦ βραδεῖς. βέλτιον δὲ τοὺς ἐν ἑτέρῳ τινὶ γενομένους νοεῖν· ἀβλαβεῖς γὰρ οἱ ἐπιορκήσαντες ; App. Prov.~2.88 Εὐορκότερος Τενεάτου, ἢ Γενεάτου: ἐπὶ τῶν εὐορκούντων}
\lp{post 276 ὀλολύζουσι]}{Aesch.~\textit{Sept.}~268 ὀλολυγμὸν ἱερὸν εὐμενῆ παιώνισον ; \textit{Il.}~6.301 αἳ δ᾽ ὀλολυγῇ πᾶσαι Ἀθήνῃ χεῖρας ἀνέσχον ; Hdt.~4.189.3 δοκέει δ᾽ ἔμοιγε καὶ ὀλολυγὴ ἐν ἱροῖσι ἐνθαῦτα πρῶτον γενέσθαι· κάρτα γὰρ ταύτῃ χρέωνται καλῶς αἱ Λίβυσσαι ; Σ Ar.~\textit{Thesm.}~129 ὀλολύζει. παρεπιγραφὴ τοῦτο· εἰπόντος γὰρ τοῦ Ἀγάθωνος ὁ κηδεστὴς τοῦ Εὐριπίδου ὀλολύζει εἰς τὴν θηλύτητα αὐτοῦ καὶ μόνον οὐχὶ γυναικὸς εὐχήν ; Hsch.~ο~613 ὀλολυγή· φωνὴ γυναικῶν, ἣν ποιοῦνται ἐν τοῖς ἱεροῖς εὐχόμεναι (Ζ 301). καὶ ἄνθος τι παρὰ λίμναις γινόμενον ; Ar.~\textit{Eq.}~1327 ἀλλ᾽ ὀλολύξατε φαινομέναισιν ταῖς ἀρχαίαισιν Ἀθήναις ; Eur.~\textit{Med.}~1176 εἶτ᾽ ἀντίμολπον ἧκεν ὀλολυγῆς μέγαν ; Aesch.~\textit{Ag.}~595 ὀλολυγμὸν ἄλλος ἄλλοθεν κατὰ πτόλιν ; Σ Ar.~\textit{Pax}~97 ὀλολύζειν : Τὸ μετὰ κραυγῆς εὔχεσθαι. — Ὅμηρος [Il. Ζ, 301] « αἱ δ᾽ ὀλολυγῇ πᾶσαι. » ; \textit{idem} Suid.~ο~193 (αἱ δ᾽ ὀλόλυξαν) ; Hsch.~ο~612 ὀλολυγή· ποιὰ φωνὴ λυπηρά, ὀδύνην καρδίας ἀσήμῳ τινὶ φθόγγῳ παριστῶσα ; Ar.~\textit{Pax}~97 μηδὲν γρύζειν ἀλλ᾽ ὀλολύζειν}
\lp{post 276 τὸ ἱερὸν ὠθεῖται]}{Σ Ar.~\textit{Thesm.}~277 παρεπιγραφή. ἐκκυκλεῖται ἐπὶ τὸ ἔξω τὸ θεσμοφόριον ; Hsch.~ε~4014 ἐξώστρα· ἐπὶ τῆς σκηνῆς τὸ ἐκκύκλημα (Phryn. trag. 3 F 21a, com. fr. 90 K.–A.⟩ ; Poll.~4.128 καὶ τὸ μὲν ἐκκύκλημα, ἐπὶ ξύλων, ὑψηλὸν βάθρον, ᾧ ἐπίκειται θρόνος. δείκνυσι δὲ καὶ τὰ ὑπὸ τὴν σκηνὴν ἐν ταῖς οἰκίαις ἀπόρρητα πραχθέντα. καὶ τὸ ῥῆμα τοῦ ἔργου καλεῖται ἐγκυκλεῖν ; Poll.~4.129 τὴν δὲ ἐξώστραν ταὐτὸν τῷ ἐγκυκλήματι νομίζουσιν ; Cic. Prov. cons.~14 quibuscum iam in exostra helluatur, antea post siparium solebat ; Plb.~11.5.8 τῆς τύχης ὥσπερ ἐπίτηδες ἐπὶ τὴν ἐξώστραν ἀναβιβαζούσης τὴν ὑμετέραν ἄγνοιαν}
\lp{post 276 παρεπιγραφή (Σ)]}{Σ Eur.~\textit{Or.}~1384 τινὲς τοῦτο παρεπιγραφὴν εἶναι ὡς εἰς τὰ κωμικὰ δράματα ; Suid.~δ~804 Διαύλειον προσαυλεῖ τις: τοῦτο παρεπιγραφή. φασὶ δὲ διαύλειον λέγεσθαι, ὅταν ἡσυχίας πάντων γινομένης, ἔνδον ὁ αὐλητὴς ᾄσῃ. Ἀριστοφάνης Βατράχοις}
\lp{277–278 ἔκσπευδε ταχέως· ὡς τὸ τῆς ἐκκλησίας σημεῖον ἐν τῷ Θεσμοφορίῳ φαίνεται]}{\textit{Test.} Suid.~σ~325 καὶ ὅτε δὲ ἔμελλε γίνεσθαι ἐκκλησία, σημεῖον ἐτίθετο. οὕτως οὖν καὶ τῶν γυναικῶν μελλουσῶν ἐν τοῖς Θεσμοφορίοις ἐκκλησιάζειν, σημεῖον ἐτέθη. Ἀριστοφάνης· ὡς τὸ τῆς ἐκκλησίας σημεῖον ἐν τῷ Θεσμοφόρῳ φαίνεται}
\lp{277 ἔκσπευδε ταχέως]}{Ar.~\textit{Eq.}~495 καὶ σπεῦδε ταχέως ; Tosi n°~696 Festina lente ; Erasm.~\textit{Adag.}~2.1.1 Σπεῦδε βραδέως, i. e. Festina lente. … Nec absurde mihi conjectasse videatur, si quis effictum existimet ex eo, quod est apud Aristophanem in Equitibus, σπεῦδε ταχέως, i. e. Propera propere, ut allusor, quisquis is fuit, τὴν ἀναδίπλωσιν in ἐναντίωσιν commutarit ; Ar.~\textit{Ran.}~188 ταχέως ἔμβαινε ; Aesch.~\textit{Pers.}~742 ἀλλ᾽ ὅταν σπεύδῃ τις αὐτός, χὠ θεὸς συνάπτεται ; Gell.~10.11.5 Nam et dicere in sermonibus et scribere in epistulis solitum esse aiunt σπεῦδε βραδέως ; Polyaen.~8.24 Σεβαστὸς τοῖς ἑαυτοῦ στρατηγοῖς παρήγγελλεν ἀσφαλείας πεφροντικέναι μάλιστα καὶ συνεχῶς ἐπέλεγε τὸ ῾σπεῦδε βραδέως ; Ar.~\textit{Ach.}~777 φώνει δὴ τὺ ταχέως χοιρίον ; Ar.~\textit{Ach.}~1094 ἀλλ᾽ ὡς τάχιστα σπεῦδε ; Ar.~\textit{Eq.}~95 ἀλλ᾽ ἐξένεγκέ μοι ταχέως οἴνου χοᾶ ; Ar.~\textit{Av.}~991 οὔκουν ἑτέρωσε χρησμολογήσεις ἐκτρέχων ; Ar.~\textit{Thesm.}~638 χάλα ταχέως τὸ στρόφιον ; Ar.~\textit{Thesm.}~731 ἀπόδυθι ταχέως ; Ar.~\textit{Thesm.}~1011 σημεῖον ὑπεδήλωσε Περσεὺς ἐκδραμών ; Ar.~\textit{Ran.}~605 ξυνδεῖτε ταχέως τουτονὶ τὸν κυνοκλόπον ; Ar.~\textit{Plut.}~414 καὶ δὴ βαδίζω. — σπεῦδέ νυν. — τοῦτ᾽ αὐτὸ δρῶ ; Soph.~\textit{Aj.}~1164 ἀλλ᾽ ὡς δύνασαι, Τεῦκρε, ταχύνας ; Hsch.~σ~1486 σπεύδων· σπουδάζων, ἐνεργῶν}
\lp{277–278 τὸ τῆς ἐκκλησίας σημεῖον … φαίνεται]}{Ar.~\textit{Ach.}~19–22 ὡς νῦν, ὁπότ᾽ οὔσης κυρίας ἐκκλησίας / ἑωθινῆς ἔρημος ἡ πνὺξ αὑτηί, / οἱ δ᾽ ἐν ἀγορᾷ λαλοῦσι κἄνω καὶ κάτω / τὸ σχοινίον φεύγουσι τὸ μεμιλτωμένον ; Ar.~\textit{Vesp.}~689–690 ὡς ὅστις ἂν ὑμῶν / ὕστερος ἔλθῃ τοῦ σημείου, τὸ τριώβολον οὐ κομιεῖται ; Ar.~\textit{Thesm.}~83–84 κἀν Θεσμοφόροιν μέλλουσι περί μου τήμερον / ἐκκλησιάζειν ; Ar.~\textit{Thesm.}~372–377 ἔδοξε τῇ βουλῇ τάδε / τῇ τῶν γυναικῶν· … / ἐκκλησίαν ποιεῖν ἕωθεν τῇ μέσῃ / τῶν Θεσμοφορίων, ᾗ μάλισθ᾽ ἡμῖν σχολή ; And.~1.36 ὥστ᾽ ἐπειδὴ τὴν βουλὴν εἰς τὸ βουλευτήριον ὁ κῆρυξ ἀνείποι ἰέναι καὶ τὸ σημεῖον καθέλοι, τῷ αὐτῷ σημείῳ ἡ μὲν βουλὴ εἰς τὸ βουλευτήριον ᾔει, οἱ δ᾽ ἐκ τῆς ἀγορᾶς ἔφευγον ; D.C.~37.28.2 καὶ αὐτό, ἕως μὲν ἡ ἐκκλησία ἦν, ἐφρούρουν, ὁπότε δὲ διαλυθήσεσθαι ἔμελλε, τό τε σημεῖον καθῃρεῖτο ; Σ Ar.~\textit{Ach.}~22 τὸ σχοινίον φεύγουσιν : … μεμιλτωμένῳ σχοινίῳ περιβάλλοντες αὐτοὺς συνήλαυνον εἰς τὴν ἐκκλησίαν. τοῦτο δὲ ἐποίουν ὑπὲρ τοῦ μὴ βραδῦναι. … σχοινίον ἐκτείνοντες διὰ τῆς ἀγορᾶς διώκειν τὸν ὄχλον εἰς τὴν ἐκκλησίαν, ὥς φησι Πλάτων ὁ κωμικός ; Σ Ar.~\textit{Thesm.}~278 ὅτε ἔμελλε γίνεσθαι ἐκκλησία, σημεῖον ἐτίθετο. οὕτως οὖν καὶ τῶν γυναικῶν μελλουσῶν ἐκκλησιάζειν σημεῖον τίθεται ; Σ Ar.~\textit{Vesp.}~690 τὸ τριώβολον οὐ κομιεῖται : Ὡς ἔθους ὄντος πρὸς σημεῖόν τι ἐκκείμενον ἀθροίζεσθαι τοὺς δικαστάς ; Ar.~\textit{Eq.}~952 οὑμός· τὸ γοῦν σημεῖον ἕτερον φαίνεται ; Ar.~\textit{Thesm.}~89–90 ἐς Θεσμοφόροιν ἐλθεῖν. — τί δράσοντ᾽; εἰπέ μοι. / ἐκκλησιάσοντ᾽ ἐν ταῖς γυναιξὶ ; Ar.~\textit{Eccl.}~282–284 ἀλλὰ σπεύσαθ᾽ ὡς εἴωθ᾽ ἐκεῖ / τοῖς μὴ παροῦσιν ὀρθρίοις ἐς τὴν πύκνα / ὑπαποτρέχειν ἔχουσι μηδὲ πάτταλον ; Ar.~\textit{Eccl.}~378–379 καὶ δῆτα πολὺν ἡ μίλτος ὦ Ζεῦ φίλτατε / γέλων παρέσχεν, ἣν προσέρραινον κύκλῳ ; Ar.~\textit{Eccl.}~381 ἀλλ᾽ ὕστερος νῦν ἦλθον, ὥστ᾽ αἰσχύνομαι ; Th.~1.49.1 ἐπειδὴ τὰ σημεῖα ἑκατέροις ἤρθη, ἐναυμάχουν ; Hsch.~ε~1442 ἐκκλησία· σύνοδος, συναγωγή, πανήγυρις ; Hsch.~σ~502 σημεῖον· τέρας (e. g. Xen. Cyr. 2, 4, 19). ἢ σφραγίς (e. g. Ar. Equ. 952–3) ; Poll.~8.104 ληξίαρχοι … καὶ σχοινίον μιλτώσαντες διὰ τῶν τοξοτῶν συνήλαυνον τοὺς ἐκ τῆς ἀγορᾶς εἰς τὴν ἐκκλησίαν ; Suid.~σ~323 Σημεῖα ἀρθῆναι ἑκατέροις: παρὰ Θουκυδίδῃ. σύμβολά τινα εἶχον, ἃ κατὰ τὸν καιρὸν τῆς μάχης δεικνύμενα. καὶ ὅτε μὲν ἀνετείνετο τὰ σύμβολα, ἤρχοντο τῆς μάχης· ὅτε δὲ κατεσπῶντο, ἐπαύοντο ; Suid.~σ~1810 Σχοινίον μεμιλτωμένον: εἰ βραδύνοιεν ἐπὶ τὴν ἐκκλησίαν, οἱ τοξόται σχοινίον μιλτοῦντες συνήλαυνον καὶ τὰ πρατήρια διέκλειον ; Eur.~\textit{Phoen.}~1377–1378 ἐπεὶ δ᾽ ἀφείθη πυρσὸς ὣς Τυρσηνικῆς / σάλπιγγος ἠχὴ σῆμα φοινίου μάχης ; Th.~8.95.4 σημεῖον δὲ αὐτοῖς ἐς τὸν Ὠρωπὸν ἐκ τῆς Ἐρετρίας, ὁπότε χρὴ ἀνάγεσθαι, ἤρθη ; X.~\textit{Cyr.}~7.1.23 ἦρε τοῖς κέρασι σημεῖον μηκέτι ἄνω πορεύεσθαι ; Hsch.~φ~59 φαίνεται· δοκεῖ. καὶ ὅμοια (λ 336 . .) ; Suid.~μ~564 Μεμιλτωμένον· Ἀριστοφάνης· τὸ σχοινίον φεύγουσι τὸ μεμιλτωμένον. ἐπεὶ ὀκνηρῶς εἶχον οἱ Ἀθηναῖοι πρὸς τὰς συνόδους}
\lp{278 ἐν τῷ Θεσμοφορίῳ]}{Ar.~\textit{Thesm.}~657–658 ζητεῖν, εἴ που κἄλλος τις ἀνὴρ ἀνελήλυθε, καὶ περιθρέξαι / τὴν πύκνα πᾶσαν καὶ τὰς σκηνὰς καὶ τὰς διόδους διαθρῆσαι ; Ar.~\textit{Thesm.}~880 ληροῦντι λῆρον; Θεσμοφορεῖον τουτογί ; Is.~8.19 αἵ τε γυναῖκες αἱ τῶν δημοτῶν μετὰ ταῦτα προὔκριναν αὐτὴν μετὰ τῆς Διοκλέους γυναικὸς τοῦ Πιθέως ἄρχειν εἰς τὰ Θεσμοφόρια ; Plu.~\textit{Dion}~56.5 καταβὰς εἰς τὸ τῶν Θεσμοφόρων τέμενος ὁ διδοὺς τὴν πίστιν, ἱερῶν τινων γενομένων, περιβάλλεται τὴν πορφυρίδα τῆς θεοῦ ; Teles \textit{Περὶ\,φυγῆς}~1 (First1K) Οὐδὲ γὰρ νῦν εἰς τὸ Θεσμοφόριον ἐξουσίαν ἔχω ; Hsch.~θ~383 θεσμοφόρια· τὰ μυστήρια τῆς Δήμητρος. ἢ ὄνομα ἑορτῆς ; Hdn.~\textit{Pros.\,cath.}~13 ὀσχοφόριον, προσφόριον, θεσμοφόριον, ὠμοφόριον, ἐμπόριον ; Hdt.~2.171.2 καὶ τῆς Δήμητρος τελετῆς πέρι, τὴν οἱ Ἕλληνες θεσμοφόρια καλέουσι, καὶ ταύτης μοι πέρι εὔστομα κείσθω ; Is.~3.80 ἠναγκάζετο ἂν ὑπὲρ τῆς γαμετῆς γυναικὸς καὶ θεσμοφόρια ἑστιᾶν τὰς γυναῖκας ; Plu.~\textit{Dem.}~30.5 κατέστρεψε δὲ ἕκτῃ ἐπὶ δέκα τοῦ Πυανεψιῶνος μηνός, ἐν ᾗ τὴν σκυθρωποτάτην τῶν Θεσμοφορίων ἡμέραν ἄγουσαι ; Luc.~\textit{DMeretr.}~2.1 εἶδον γὰρ αὐτὴν ἔναγχος ἐν τοῖς Θεσμοφορίοις μετὰ τῆς μητρός ; Clem.Al.~\textit{Strom.}~4.19.121 ἐρωτηθεῖσα γάρ, ποσταία γυνὴ ἀπὸ ἀνδρὸς εἰς τὸ θεσμοφόριον κάτεισιν ; Plb.~15.29.8 ἡ δ᾽ Οἰνάνθη περικακοῦσα παρῆν εἰς τὸ Θεσμοφορεῖον ; Plb.~15.33.8 ἐκ τοῦ Θεσμοφορείου τὴν Οἰνάνθην ἀποσπάσαντες ; Theon \textit{Prog.}~5 Θεανὼ ἡ Πυθαγορικὴ φιλόσοφος ἐρωτηθεῖσα ὑπό τινος, ποσταία γυνὴ ἀπ᾽ ἀνδρὸς καθαρὰ εἰς τὸ θεσμοφορεῖον κάτεισιν ; Hsch.~π~4128 πρυτανεῖον· τρία Ἀθήνησι συσσίτια· πρυτανεῖα, θεσμοφορεῖον, πρυτανεῖον ; Phot. p.~87 Porson (Naber I, p.~278) Θεσμοφόρια· ὄνομα ἑορτῆς ; Plu.~\textit{Pel.}~5.3 καὶ μὴ προσδοκῶσι τοῖς Θηβαίοις ἐπιθεμένου Θεσμοφορίων ὄντων}
\lp{279 ἐγὼ δ᾽ ἄπειμι]}{Ar.~\textit{Nub.}~887 ἐγὼ δ᾽ ἀπέσομαι. — τοῦτό νυν μέμνησ᾽, ὅπως ; Soph.~\textit{Phil.}~123–124 σὺ μὲν μένων νυν κεῖνον ἐνθάδ᾽ ἐκδέχου, / ἐγὼ δ᾽ ἄπειμι, μὴ κατοπτευθῶ παρών ; Ar.~\textit{Thesm.}~457 ἀλλ᾽ εἰς ἀγορὰν ἄπειμι· δεῖ γὰρ ἀνδράσιν ; Aesch.~\textit{Pers.}~839 ἐγὼ δ᾽ ἄπειμι γῆς ὑπὸ ζόφον κάτω ; Hsch.~α~5917 ἄπεισιν· ἀπέρχεται. ἀπέθανεν}
\lp{279 δεῦρό νυν ὦ Θρᾷτθ᾽ ἕπου]}{Σ Ar.~\textit{Thesm.}~279 μιμεῖται τὴν φωνὴν γυναικός}
\lp{279 δεῦρό νυν … ἕπου]}{Ar.~\textit{Plut.}~823–824 ἕπου μετ᾽ ἐμοῦ παιδάριον, ἵνα πρὸς τὸν θεὸν / ἴωμεν ; Hsch.~δ~731 δεῦρο· ἐνθάδε. ἐλθέ εἰς τοῦτον τὸν τόπον. ἢ ἐπὶ τοῦ παρόντος, ἢ νῦν. ἢ παρακελευστικῶς· «δεῦρό νυν, ἢ τρίποδος περιδώμεθα» (Ψ 485) ; Ar.~\textit{Lys.}~930 δεῦρό νυν ὦ χρύσιον ; Ar.~\textit{Eccl.}~1005 ἀλλ᾽ ἕπου δεῦρ᾽ ὡς ἐμέ ; Ar.~\textit{Eccl.}~1058 ἕπου μαλακίων δεῦρ᾽ ἀνύσας ; Eur.~\textit{Ba.}~363 ἀλλ᾽ ἕπου μοι κισσίνου βάκτρου μέτα ; Eur.~\textit{HF}~724 δεῦρ᾽ ἕπεσθε, πρόσπολοι ; Eur.~\textit{Ion}~741 ἕπου νυν· ἴχνος δ᾽ ἐκφύλασσ᾽ ὅπου τίθης ; \textit{Il.}~23.485 δεῦρό νυν ἢ τρίποδος περιδώμεθον ἠὲ λέβητος ; Σ Ar.~\textit{Nub.}~91 τὸ δὲ δεῦρο οὐκ ἔστι τοπικόν, ἀλλ᾽ ἀντὶ τοῦ ἐλθέ. R. V ; Ar.~\textit{Ach.}~204 τῇδε πᾶς ἕπου ; Ar.~\textit{Nub.}~91 δεῦρό νυν ἀπόβλεπε ; Ar.~\textit{Thesm.}~27 σίγα νυν ; Ar.~\textit{Ran.}~1368 ἴτε δεῦρό νυν, εἴπερ γε δεῖ καὶ τοῦτό με ; Ar.~\textit{Eccl.}~1074 μὴ σκῶπτέ μ᾽ ἀλλὰ δεῦρ᾽ ἕπου ; Plaut.~\textit{Curc.}~215 Sequere me ; Eur.~\textit{Phoen.}~1274 οὐ φίλ᾽, ἀλλ᾽ ἕπου ; Σ Eur.~\textit{Phoen.}~1274 οὐ φίλ᾽, ἀλλ᾽ ἕπου : οὐ προσφιλῆ, ἀλλ᾽ ἀκολούθει ; Hsch.~ε~5515 ἕπου· ἀκολούθει (Eur. Phoen. 1274 . .)}
\lp{279 Θρᾷττα]}{Ar.~\textit{Ach.}~273 τὴν Στρυμοδώρου Θρᾷτταν ἐκ τοῦ Φελλέως ; Ar.~\textit{Vesp.}~828 ἡ Θρᾷττα προσκαύσασα πρώην τὴν χύτραν — ; Ar.~\textit{Pax}~1138 χἄμα τὴν Θρᾷτταν κυνῶν ; Pl.~\textit{Tht.}~174a Θρᾷττά τις ἐμμελὴς καὶ χαρίεσσα θεραπαινὶς ἀποσκῶψαι λέγεται ; D.~59.35 καὶ θεραπαίνας δύο, Θρᾷτταν καὶ Κοκκαλίνην ; Str.~7.3.12 ἐξ ὧν γὰρ ἐκομίζετο, ἢ τοῖς ἔθνεσιν ἐκείνοις ὁμωνύμους ἐκάλουν τοὺς οἰκέτας, ὡς Λυδὸν καὶ Σύρον ; Suid.~θ~464 Θρᾷττα: Θρᾳκικὴ δούλη, ἐκ Θράκης. Ἀριστοφάνης· Στρυμνοδώρου Θρᾷτταν. τουτέστι Θρᾳκικὴν δούλην. λέγεται καὶ κοινῶς δούλη ; App. Prov.~3.5 Ἡ Θρᾷττα ἐπὶ τῷ ταρίχει: ἐπὶ τῶν κρυπτόντων τὰ ἴδια πάθη, ἢ ἐπὶ τῶν προφασιζομένων ἵνα τύχοιεν τινός. Καὶ γὰρ ἡ Θρᾷττα φίλοινος οὖσα, διψῆν ἔφασκε διὰ τὸ φαγεῖν τάριχος ; Ar.~\textit{Lys.}~184 καλῶς λέγεις. ποῦ ᾽σθ᾽ ἡ Σκύθαινα; ποῖ βλέπεις ; Ar.~\textit{Thesm.}~728 ἴωμεν ἐπὶ τὰς κληματίδας ὦ Μανία ; Ar.~\textit{Thesm.}~1112 ληρεῖς ὦ Σκύθα ; Ar.~\textit{Ran.}~1345 ὦ Μανία ξύλλαβε ; Eup.~fr.~262~K.-A. (=~243~Kock) (Ath.~7.326a) μήτηρ τις αὐτῷ Θρᾷττα ταινιόπωλις ἦν ; Herod.~1.1 Θ[ρέϊσ]σα, ἀράσσει τὴν θύρην τις ; Zen.~2.12 Ἁλώνητον ἀνδράποδον: ἶσον τῷ, βάρβαρον καὶ εὐτελές. … Ὅθεν καὶ ὁ κωμικὸς φησί· Θρᾷξ εὐγενὴς εἶ πρὸς ἅλας ἠγορασμένος ; \textit{idem} Diogenian.~1.100 (τὸ εὐτελὲς, τὸ πρὸς ἅλας ὠνηθέν pro ἶσον τῷ, βάρβαρον καὶ εὐτελές ; Μένανδρος pro ὁ κωμικὸς φησί ; ὠνημένος), Erasm.~\textit{Adag.}~1.7.79 (Sale emptum mancipium. Ἁλώνητον ἀνδράποδον … ; ἠγορασμένος) ; Ar.~\textit{Thesm.}~739 παράβαλλε πολλὰς κληματίδας ὦ Μανία ; Ar.~\textit{Thesm.}~754 δός μοι σφαγεῖον Μανία ; Hor.~\textit{Carm.}~3.9.9 me nunc Thressa Chloe regit ; Arist.~\textit{Ath.}~14.4 ὡς δ᾽ ἔνιοι λέγουσιν ἐκ τοῦ Κολλυτοῦ στεφανόπωλιν Θρᾷτταν, ᾗ ὄνομα Φύη ; Ael.~\textit{VH}~12.43 Θρᾴττης υἱὸς ἦν ; Hsch.~ζ~155 Ζιβυθίδες· αἱ Θρᾷσσαι, ἢ Θρᾷκες γνήσιοι ; Zen.~4.32 Θρᾷκες ὅρκια οὐκ ἐπίστανται}
\lp{280 ὦ Θρᾷττα θέασαι]}{Alciphr.~2.37 θέασαι πρὸ τῆς τελευταίας ἡμέρας τὰ κατ᾽ ἄστυ καλά. … καὶ ἡ νῦν ἑστῶσα σεμνοτάτη τῶν Θεσμοφορίων ἑορτή ; Ar.~\textit{Ach.}~770 οὐ δεινά; θᾶσθε τῶδε τὰς ἀπιστίας ; Ar.~\textit{Pax}~907 θέασ᾽ ὡς προθύμως ὁ πρύτανις παρεδέξατο ; Ar.~\textit{Eq.}~997 ἰδοὺ θέασαι, κοὐχ ἅπαντας ἐκφέρω ; Ar.~\textit{Vesp.}~1215 ὀροφὴν θέασαι, κρεκάδι᾽ αὐλῆς θαύμασον ; Eur.~\textit{HF}~1131 ἰδού, θέασαι τάδε τέκνων πεσήματα}
\lp{280 καομένων τῶν λαμπάδων]}{Ar.~\textit{Thesm.}~101–102 ἱερὰν χθονίαις δεξάμεναι / λαμπάδα κοῦραι ξὺν ἐλευθέρᾳ ; Ar.~\textit{Thesm.}~655 ἡμᾶς τοίνυν μετὰ τοῦτ᾽ ἤδη τὰς λαμπάδας ἁψαμένας χρὴ ; Ar.~\textit{Thesm.}~1152 ὄργια σεμνὰ θεοῖν, ἵνα λαμπάσι ; Ar.~\textit{Ran.}~340 ἔγειρε φλογέας λαμπάδας ἐν χερσὶ γὰρ ἥκει τινάσσων ; Eur.~\textit{Or.}~1573 ἔα, τί χρῆμα; λαμπάδων ὁρῶ σέλας ; Soph.~\textit{OC}~1048–1051 ἢ λαμπάσιν ἀκταῖς, / οὗ πότνιαι σεμνὰ τιθηνοῦνται τέλη / θνατοῖσιν ; \textit{h.Cer.}~47–48 ἐννῆμαρ μὲν ἔπειτα κατὰ χθόνα πότνια Δηὼ / στρωφᾶτ᾽ αἰθομένας δαΐδας μετὰ χερσὶν ἔχουσα ; D.~43.14 λαβόντες τὴν ψῆφον καομένων τῶν ἱερείων ; Apollod.~1.5.1 Δημήτηρ δὲ μετὰ λαμπάδων νυκτός τε καὶ ἡμέρας κατὰ πᾶσαν τὴν γῆν ζητοῦσα περιῄει ; Paus.~2.22.3 ἀφιᾶσι δὲ καὶ νῦν ἔτι ἐς τὸν βόθρον καιομένας λαμπάδας Κόρῃ τῇ Δήμητρος ; Ar.~\textit{Thesm.}~917 ὅστις σ᾽ ἀπάξει τυπτόμενος τῇ λαμπάδι ; Eur.~\textit{Ion}~1074–1077 εἰ παρὰ καλλιχόροισι παγαῖς / λαμπάδα θεωρὸν εἰκάδων / ὄψεται ἐννύχιος ; Ael.~\textit{VH}~13.2 ἔτι καομένων τῶν ἐμπύρων ; Σ Soph.~\textit{OC}~1048 τὰ περὶ Ἐλευσῖνα μέρη φησὶ καὶ εἴη λαμπάσιν ἀκταῖς ταῖς λαμπαδευομέναις καὶ καταλαμπομέναις καὶ δᾳδουχουμέναις ὑπὸ τῆς μυστικῆς φλογὸς καὶ τῶν ἱερῶν δᾴων ; Hsch.~δ~53 δαΐδες· λαμπάδες. λύχνοι}
\lp{280 καομένων]}{Moer. p.~201,~24~Bekker κλάειν καὶ κάειν σὺν τῷ ᾱ Ἀττικοί, μετὰ τοῦ ῑ Ἕλληνες ; Ar.~\textit{Vesp.}~1372–1375 ἐν ἀγορᾷ τοῖς θεοῖς δᾲς κάεται. / … / ἡ πίττα δήπου καομένης ἐξέρχεται ; Ar.~\textit{Lys.}~319–320 λιγνὺν δοκῶ μοι καθορᾶν καὶ καπνὸν ὦ γυναῖκες / ὥσπερ πυρὸς καομένου· σπευστέον ἐστὶ θᾶττον}
\lp{281 ὅσον τὸ χρῆμ᾽]}{Ar.~\textit{Ach.}~150 ὅσον τὸ χρῆμα παρνόπων προσέρχεται ; \textit{idem} Suid.~π~683 (praem. καὶ αὖθις ; add. τὸ πλῆθος σημαίνει τῆς στρατιᾶς) ; Ar.~\textit{Pax}~1192 ὅσον τὸ χρῆμ᾽ ἐπὶ δεῖπνον ἦλθ᾽ ἐς τοὺς γάμους ; Ar.~\textit{Eccl.}~394–395 ἀτὰρ τί τὸ πρᾶγμ᾽ ἦν, ὅτι τοσοῦτον χρῆμ᾽ ὄχλου / οὕτως ἐν ὥρᾳ ξυνελέγη ; Ar.~\textit{Eq.}~1219 ὅσον τὸ χρῆμα τοῦ πλακοῦντος ἀπέθετο ; Ar.~\textit{Nub.}~2 ὦ Ζεῦ βασιλεῦ τὸ χρῆμα τῶν νυκτῶν ὅσον ; Ar.~\textit{Lys.}~83 ὡς δὴ καλὸν τὸ χρῆμα τιτθίων ἔχεις ; Ar.~\textit{Ran.}~1278 ὦ Ζεῦ βασιλεῦ τὸ χρῆμα τῶν κόπων ὅσον ; Ar.~fr.~73~K.-A. (Phot.~ν~139) ὦ Ζεῦ, τὸ χρῆμα τῆς νεολαίας ὡς καλόν ; Σ Ar.~\textit{Ach.}~150 ὅσον τὸ χρῆμα : Ὡς πολλῶν ὄντων παρνόπων ἐν τῇ Ἀττικῇ. βούλεται δὲ πολὺ πλῆθος σημᾶναι ; Hdt.~1.36.1 ἐν τῷ Μυσίῳ Ὀλύμπῳ ὑὸς χρῆμα γίνεται μέγα ; Hsch.~χ~709 χρῆμα· πρᾶγμα. πλοῦτος, οὐσία. λῆμμα}
\lp{281 ἀνέρχεθ᾽]}{Ar.~\textit{Thesm.}~585 αὑτοῦ γέροντα δεῦρ᾽ ἀναπέμψαι τήμερον ; Ar.~\textit{Thesm.}~623 ἀνῆλθες ἤδη δεῦρο πρότερον ; Ar.~\textit{Thesm.}~893 οὗτος πανουργῶν δεῦρ᾽ ἀνῆλθεν ὦ ξένε ; Alciphr.~2.37 ἡ μὲν οὖν Ἄνοδος κατὰ τὴν πρώτην γέγονεν ἡμέραν, ἡ Νηστεία δὲ τὸ τήμερον εἶναι παρ᾽ Ἀθηναίοις ἑορτάζεται ; Σ Ar.~\textit{Thesm.}~585 ὅτι ἀναπέμψαι κυρίως. διὸ καὶ ἄνοδος ἡ πρώτη λέγεται, παρ᾽ ἐνίοις καὶ κάθοδος, διὰ τὴν θέσιν τῶν θεσμοφορίων. ἐπεὶ καὶ ἄνοδον τὴν εἰς τὸ θεσμοφόριον ἄφιξιν λέγουσιν· ἐπὶ ὑψηλοῦ γὰρ κεῖται τὸ θεσμοφόριον ; Σ Ar.~\textit{Thesm.}~623 καὶ τοῦτο πρὸς τὴν ἄνοδον, ἐπειδὴ ἄνοδος ἦν πρὸς τὸ ἱερόν ; Hsch.~α~5234 ἄνοδος· ἀνάβασις. ἡ ἑνδέκατη τοῦ Πυανεψιῶνος, ὅτε αἱ γυναῖκες ἀνέρχονται εἰς Θεσμοφόριον, οὕτω καλεῖται ; Ar.~\textit{Thesm.}~1045 ἐπὶ δὲ τοῖσδ᾽ ⟨ἐς⟩ τόδ᾽ ἀνέπεμψεν ; Σ Ar.~\textit{Thesm.}~80 ἑνδεκάτη Πυανεψιῶνος ἄνοδος. δωδεκάτη νηστεία ; Phot.~σ~550 Στήνια· ἑορτὴ Ἀθήνησιν, ἐν ᾗ ἐδόκει ἡ ἄνοδος γενέσθαι τῆς Δήμητρος· ἐλοιδοροῦντο δ᾽ ἐν αὐτῇ νυκτὸς αἱ γυναῖκες ἀλλήλαις. οὕτως Εὔβουλος (fr. 146 K.- A.) ; Paus.~1.22.3 ἔστι δὲ καὶ Γῆς Κουροτρόφου καὶ Δήμητρος ἱερὸν Χλόης ; Hsch.~σ~1825 στήνια· ἑορτὴ Ἀθήνησι· καὶ διασκώπτουσι καὶ λοιδοροῦσιν (Ar. Thesm. 834. Eubul. fr. 146 K.–A.)}
\lp{281 ὑπὸ τῆς λιγνύος]}{Ar.~\textit{Av.}~1241–1242 λιγνὺς δὲ σῶμα καὶ δόμων περιπτυχὰς / καταιθαλώσῃ σου Λικυμνίαις βολαῖς ; Ar.~\textit{Lys.}~319 λιγνὺν δοκῶ μοι καθορᾶν καὶ καπνὸν ὦ γυναῖκες ; Ar.~\textit{Ran.}~313–314 ἔγωγε, καὶ δᾴδων γέ με / αὔρα τις εἰσέπνευσε μυστικωτάτη ; Eur.~\textit{Ion}~1474–1475 οὐχ ὑπὸ λαμπάδων οὐδὲ χορευμάτων / ὑμέναιος ἐμός ; Aesch.~\textit{Sept.}~493–494 Τυφῶν᾽ ἱέντα πύρπνοον διὰ στόμα / λιγνὺν μέλαιναν, αἰόλην πυρὸς κάσιν ; Soph.~\textit{Ant.}~1126–1128 σὲ δ᾽ ὑπὲρ διλόφου πέτρας στέροψ ὄπωπε / λιγνύς, ἔνθα Κωρύκιαι / στείχουσι νύμφαι Βακχίδες ; \textit{Il.}~18.492–493 νύμφας δ᾽ ἐκ θαλάμων δαΐδων ὕπο λαμπομενάων / ἠγίνεον ἀνὰ ἄστυ ; \textit{Od.}~23.290 ἐσθῆτος μαλακῆς, δαΐδων ὕπο λαμπομενάων ; Archil.~fr.~123~Bergk (Σ Il.~18.492) ᾄδων ὑπ᾽ αὐλητῆρος ; \textit{idem} Σ Il.~18.492 (praem. ὅτι ἡ ὑπό ἀντὶ τῆς μετά, μετὰ δᾴδων· Ἀρχίλοχος) ; Arist.~\textit{Mete.}~387b6 (4.9) ἡ δὲ πίονος θυμίασις λιγνύς ; Σ Ar.~\textit{Thesm.}~281 Put. τῆς ἀναδόσεως τοῦ καπνοῦ ; Σ Pi.~\textit{O.}~4.3h τῇ δὲ ὑπό ἀντὶ τῆς μετά καὶ Ὅμηρος κέχρηται · δαΐδων ὑπὸ λαμπομενάων· ἀντὶ τοῦ μετὰ δαΐδων ; Hsch.~λ~972 λιγνύς· ἡ τοῦ λύχνου ἀτμίς ; Synag.~λ~100~Cunningham λιγνύς· ἡ ἀνάδοσις τοῦ καπνοῦ ; \textit{idem} Phot. s.v.~λιγνύς (Naber I, p.~387), Suid.~λ~502 (add. Ἀριστοφάνης· λιγνὺς δὲ σῶμα καὶ δόμων περιπτυχὰς καταιθαλώσει Λικυμνίαισι βολαῖς) ; Procl. Chrest. ap. Phot. Bibl.~239, p.~320b33–35~Bekker ὑπόρχημα δὲ τὸ μετ᾽ ὀρχήσεως ᾀδόμενον μέλος ἐλέγετο· καὶ γὰρ οἱ παλαιοὶ τὴν ὑπό ἀντὶ τῆς μετά πολλάκις ἐλάμβανον ; Eur.~\textit{Hel.}~638–640 ἃν ὑπὸ λαμπάδων κόροι λεύκιπποι / ξυνομαίμονες ὤλβισαν ; Soph.~\textit{Trach.}~794 τότ᾽ ἐκ προσέδρου λιγνύος διάστροφον ; Pi.~\textit{O.}~4.2 ὑπὸ ποικιλοφόρμιγγος ἀοιδᾶς ἑλισσόμεναί μ᾽ ἔπεμψαν ; A.R.~4.1188 θυέων δ᾽ ἀποτηλόθι κήκιε λιγνύς ; Th.~5.70.1 Λακεδαιμόνιοι δὲ βραδέως καὶ ὑπὸ αὐλητῶν πολλῶν ; Plb.~1.48.6 ἀποσκοτουμένους ὑπὸ τῆς εἰς αὐτοὺς φερομένης λιγνύος ; Dsc.~5.162.1 μέλαν, ᾧ γράφομεν, σκευάζεται ἐκ λιγνύος συναγομένης ἐκ δᾳδίων ; Σ Soph.~\textit{Ant.}~1008 κυρίως δὲ ἡ λιγνὺς τοῦ καπνοῦ [ἢ] ἡ ἀπὸ τῶν καιομένων γινομένη ; Hsch.~δ~54 δαΐδων· λαμπάδων (Σ 492) ; Hsch.~λ~1053 †λιγνύς· καπνός ; Arist.~\textit{Mete.}~374a24–26 (3.4) ἀπὸ λιγνύος τῆς ἀπὸ τῆς φλογὸς ἀπορρεούσης … καπνώδης γὰρ ἡ λιγνύς}
\lp{282 ἀλλ᾽, ὦ περικαλλεῖ Θεσμοφόρω]}{\textit{Test.} Phot.~α~1386 (Ael.~\textit{Dion.} α~113 Erbse) ἀμφορῆ· τὸ δυϊκὸν συναλείφοντες οἱ Ἀττικοὶ λέγουσιν ἀντὶ τοῦ ἀμφορέε καὶ τὰ ἄλλα ὁμοίως. Ἀριστοφάνης Θεσμοφοριαζούσαις· „ἀλλ᾽, ὦ περικαλλῆ Θεσμοφόρω, δέξασθέ με“ ; Ar.~\textit{Av.}~368 τῆς ἐμῆς γυναικὸς ὄντε ξυγγενεῖ καὶ φυλέτα ; Σ Ar.~\textit{Thesm.}~282 δυϊκόν. τὸ ἐντελὲς περικαλλέε}
\lp{282 ἀλλ᾽, ὦ]}{Ar.~\textit{Nub.}~1478 ἀλλ᾽ ὦ φίλ᾽ Ἑρμῆ μηδαμῶς θύμαινέ μοι ; Ar.~\textit{Vesp.}~323 ἀλλ᾽ ὦ Ζεῦ ⟨Ζεῦ⟩ μέγα βροντήσας ; Soph.~\textit{OC}~101 ἀλλά μοι, θεαί}
\lp{282 περικαλλεῖ]}{\textit{h.Cer.}~493 αὐτὴ καὶ κούρη περικαλλὴς Περσεφόνεια ; \textit{h.Hom.}~13.1–2 Δημήτρ᾽ ἠύκομον, σεμνὴν θεάν, ἄρχομ᾽ ἀείδειν, / αὐτὴν καὶ κούρην, περικαλλέα Περσεφόνειαν ; Ar.~\textit{Thesm.}~299 Πλούτῳ καὶ τῇ Καλλιγενείᾳ καὶ τῇ Κουροτρόφῳ ; \textit{Il.}~5.389 μητρυιὴ περικαλλὴς Ἠερίβοια ; \textit{h.Cer.}~405 τὴν δ᾽ αὖ Περσεφόνη περικαλλὴς ἀντίον ηὔδα ; Hsch.~κ~472 Καλλιγένειαν· οὐ τὴν γῆν, ἀλλὰ τὴν Δήμητραν. οὐδεὶς γὰρ οὕτως ἔφη τὴν γῆν καλλιγένειαν. οἱ μὲν τροφὸν αὐτῆς· οἱ δὲ ἱέρειαν· οἱ δὲ ἀκόλουθον (Ar. Thesm. 299) ; Thgn.~1277 τῆμος Ἔρως προλιπὼν Κύπρον, περικαλλέα νῆσον ; \textit{h.Merc.}~244 νύμφην τ᾽ οὐρείην περικαλλέα ; Hdt.~5.60.1 (épigramme) νικήσας ἀνέθηκε τεῒν περικαλλὲς ἄγαλμα ; \textit{idem} Hdt.~5.61.1 (épigramme) (μουναρχέων pro νικήσας) ; Hdt.~7.5.3 ὡς ἡ Εὐρώπη περικαλλὴς εἴη χώρη ; Hsch.~π~1723 περικαλλές· λίαν καλόν (Λ 632 . .) ; \textit{idem} Phot.~π~700}
\lp{282 Θεσμοφόρω]}{Ar.~\textit{Thesm.}~297–298 εὔχεσθε τοῖν / Θεσμοφόροιν τῇ Δήμητρι καὶ τῇ Κόρῃ ; Ar.~\textit{Thesm.}~1156 Θεσμοφόρω πολυποτνία ; Ar.~\textit{Eccl.}~155–158 ἐμοὶ μὲν οὐ δοκεῖ μὰ τὼ θεώ. / μὰ τὼ θεώ; τάλαινα ποῦ τὸν νοῦν ἔχεις; / … / μὰ Δί᾽ ἀλλ᾽ ἀνὴρ ὢν τὼ θεὼ κατώμοσας ; Σ Ar.~\textit{Eccl.}~155 μὰ τὼ θεώ : (Ἐπειδὴ γυναικῶν ὁ ὅρκος.) παίζει. οὐ γὰρ δεῖ σε, φησίν, ὑποκρινομένην ἄνδρα γυναικεῖον ὅρκον ὀμνύναι τὰς θεάς ; Hsch.~μ~413 μὰ τὼ θεώ· οὐ μόνον γυναῖκες, ἀλλὰ καὶ ἄνδρες ὤμνυον ; Phot. s.v.~μὰ τὼ θεώ (Naber I, p.~409) μὰ τὼ θεώ· γυναικεῖος ὅρκος· δυϊκῶς δὲ ὀμνύουσι τὴν κόρην καὶ τὴν Δήμητραν· ἀνδράσι δὲ οὐ πρέπει τοῦτον ὀμνύναι ; Ar.~\textit{Thesm.}~83 κἀν Θεσμοφόροιν μέλλουσι περί μου τήμερον ; Ar.~\textit{Thesm.}~383 φιλοτιμίᾳ μὲν οὐδεμιᾷ μὰ τὼ θεὼ ; Verg.~\textit{Aen.}~4.58 legiferae Cereri Phoeboque patrique Lyaeo ; Hsch.~θ~384 θεσμοφόρος· ὁ θεσμοὺς φέρων, ἤτοι νόμους ; Hsch.~ν~543 νὴ τὼ θεώ· ὅρκος κατὰ Δήμητρος καὶ Κόρης ; Suid.~θ~271 Θεσμοφόρος: ὁ τοὺς νόμους ἐπιφερόμενος. καὶ ἐν Ἐπιγράμμασι· ὡμολόγησ᾽ ἥξειν εἰς νύκτα μοι ἐπιβόητος, καὶ σεμνὴν ὤμοσε Θεσμοφόρον ; Suid.~θ~272 Θεσμοφόρος. ὅτι Βάττος ὁ Κυρήνην κτίσας τῆς Θεσμοφόρου τὰ μυστήρια ἐγλίχετο μαθεῖν καὶ προσῆγε βίαν λίχνοις ὀφθαλμοῖς χαριζόμενος ; Hdt.~6.91.2 Δήμητρος Θεσμοφόρου ; Plu.~\textit{De\,esu\,carn.}~1.2 (993c–994b) τί τὴν θεσμοφόρον ἀσεβεῖτε Δήμητρα}
\lp{282 δέξασθέ με]}{Aesch.~\textit{Supp.}~27 ὁσίων ἀνδρῶν, δέξασθ᾽ ἱκέτην ; Ar.~\textit{Plut.}~1147 ἀλλὰ ξύνοικον πρὸς θεῶν δέξασθέ με ; Soph.~\textit{OC}~44 ἀλλ᾽ ἵλεῳ μὲν τὸν ἱκέτην δεξαίατο ; Pi.~\textit{N.}~11.3 εὖ μὲν Ἀρισταγόραν δέξαι τεὸν ἐς θάλαμον ; X.~\textit{Cyr.}~2.1.1 προσηύχοντο αὖθις θεοῖς τοῖς Μηδίαν γῆν κατέχουσιν ἵλεως καὶ εὐμενεῖς δέχεσθαι αὐτούς ; Luc.~\textit{Peregr.}~36 δαίμονες μητρῷοι καὶ πατρῷοι, δέξασθέ με εὐμενεῖς ; Soph.~\textit{OC}~487 στέρνων δέχεσθαι τὸν ἱκέτην σωτήριον ; Alc. ap. Heph. Ench.~5.2 δέξαι με κωμάζοντα, δέξαι, λίσσομαί σε, λίσσομαι}
\lp{282–283 ἀλλ᾽, ὦ … δέξασθέ με]}{Soph.~\textit{El.}~67–68 ἀλλ᾽, ὦ πατρῴα γῆ θεοί τ᾽ ἐγχώριοι, / δέξασθέ μ᾽ εὐτυχοῦντα ταῖσδε ταῖς ὁδοῖς ; Σ Soph.~\textit{El.}~68 δέξασθέ μ᾽ εὐτυχοῦντα: ἐπ᾽ εὐτυχίᾳ με δέξασθε ὥστε ταύτας τὰς ὁδοὺς τελεσθῆναι}
\lp{283 ἀγαθῇ τύχῃ]}{Ar.~\textit{Vesp.}~869 ὦ Φοῖβ᾽ Ἄπολλον Πύθι᾽ ἐπ᾽ ἀγαθῇ τύχῃ ; Ar.~\textit{Av.}~435–436 ταύτην λαβόντε κρεμάσατον τύχἀγαθῇ ; Ar.~\textit{Eccl.}~131 περίθου δὴ τὸν στέφανον τύχἀγαθῇ ; Th.~4.118.11 τύχῃ ἀγαθῇ τῇ Ἀθηναίων ; Pl.~\textit{Cri.}~43d ἀλλ᾽, ὦ Κρίτων, τύχῃ ἀγαθῇ, εἰ ταύτῃ τοῖς θεοῖς φίλον, ταύτῃ ἔστω ; Tosi n°~1222 Quod bonum felix faustumque sit ; Ar.~\textit{Av.}~675 ἡγοῦ δὴ σὺ νῷν τύχἀγαθῇ ; Nicostr.~fr.~18.3~K.-A. (=~19~Kock) φέρε, τύχἀγαθῇ ; Pl.~\textit{Lg.}~625b ἀλλ᾽ ἴωμεν ἀγαθῇ τύχῃ ; Pl.~\textit{Smp.}~177e ἀλλὰ τύχῃ ἀγαθῇ καταρχέτω Φαῖδρος ; Pl.~\textit{Ti.}~26e ἀλλ᾽ ἀγαθῇ τύχῃ χρὴ λέγειν μὲν ὑμᾶς ; X.~\textit{Cyr.}~4.5.51 καὶ ἀγαθῇ τύχῃ ἡμεῖς τε ἱππεῖς γενοίμεθα ; Aeschin.~3.154 ἀφίησιν ἀγαθῇ τύχῃ τρέπεσθαι ἐπὶ τὰ ἑαυτῶν ; D.~3.18 ταῦτα ποιεῖτ᾽ ἀγαθῇ τύχῃ ; Thphr.~\textit{Char.}~14.7 σκυθρωπάσας καὶ δακρύσας εἰπεῖν· ἀγαθῇ τύχῃ ; Plu.~\textit{Alc.}~19.4 ἀλλὰ νῦν μὲν ἀγαθῇ τύχῃ πλεέτω ; Hsch.~α~246 Ἀγαθὴ Τύχη· ἡ Νέμεσις καὶ ἡ Θέμις}
\lp{283 καὶ δεῦρο καὶ πάλιν οἴκαδε]}{[Eur.]~\textit{Rh.}~216–217 ἀλλ᾽ εὖ σ᾽ ὁ Μαίας παῖς ἐκεῖσε καὶ πάλιν / πέμψειεν Ἑρμῆς ; Ar.~\textit{Av.}~449 ἀνελομένους θὤπλ᾽ ἀπιέναι πάλιν οἴκαδε ; Ar.~\textit{Thesm.}~1204–1206 ὡς τὴν γυναῖκα καὶ τὰ παιδί᾽ οἴκαδε ; \textit{Il.}~1.19 ἐκπέρσαι Πριάμοιο πόλιν, εὖ δ᾽ οἴκαδ᾽ ἱκέσθαι ; Theoc.~22.60 ἔλθοις, καὶ ξενίων γε τυχὼν πάλιν οἴκαδ᾽ ἱκάνοις ; Ar.~\textit{Ach.}~84 κᾆτ᾽ ἀπῆλθεν οἴκαδε ; Ar.~\textit{Nub.}~32 ἄπαγε τὸν ἵππον ἐξαλίσας οἴκαδε ; Ar.~\textit{Lys.}~792 κοὐκέτι κατῆλθε πάλιν οἴκαδ᾽ ὑπὸ μίσους ; Ar.~\textit{Ran.}~1167 οὐ φημὶ τὸν Ὀρέστην κατελθεῖν οἴκαδε ; Pi.~\textit{N.}~3.63 πάλιν οἴκαδ᾽ ἀνεψιὸς ζαμενὴς Ἑλένοιο Μέμνων μόλοι ; Hsch.~ο~227 οἴκαδε· οἶκόνδε, εἰς οἶκον. εἰς τὸν ἴδιον οἶκον, εἰς τὰ οἰκεῖα (Α 19)}
\lp{284 τὴν κίστην]}{Ar.~\textit{Ach.}~1086 βάδιζε τὴν κίστην λαβὼν καὶ τὸν χοᾶ ; Ar.~\textit{Lys.}~1184 ξενίσωμεν ὧν ἐν ταῖσι κίσταις εἴχομεν ; \textit{Od.}~6.76 μήτηρ δ᾽ ἐν κίστῃ ἐτίθει μενοεικέ᾽ ἐδωδὴν ; Hsch.~κ~2799 κίστη· ἀγγεῖον πλεκτόν, εἰς ὃ βρῶμα ἐνετίθετο καὶ ἱμάτια. κιβωτός. [ἢ] ἅπαξ ⟨εἰρημένον⟩ (ζ 76) ; Ar.~\textit{Ach.}~1098 παῖ παῖ φέρ᾽ ἔξω δεῦρο τὴν κίστην ἐμοί ; Ar.~\textit{Eq.}~1211 ἐγὼ φράσω σοι. τὴν ἐμὴν κίστην ἰὼν ; Σ Ar.~\textit{Thesm.}~284 Put. ἀγγεῖον ἐν ᾧ τὰ βρώματα ἐτίθουν ; Hsch.~κ~3274 κοίτη· κίστη, ἐν ᾗ τὰ βρώματα ἔφερον. αἱ δὲ μικραὶ κοιτίδες ; Hsch.~κ~3275 κοιτίς· ἡ μικροτέρα κίστη, ἐν ᾗ ⟨ὁ⟩ γυναικεῖος κόσμος ἀπετίθετο ; Poll.~6.13 ἵνα δὲ τὰ ὄψα, κίστη ὀψοφόρος, ἢ κίστη δειπνοφόρος ; App. Prov.~4.17 Ὁ κάνης τῆς τύχης ὑπερέχει: ἐπὶ τῶν τὰ μὴ ἀναγκαῖα μείζω καὶ πλείω τῶν ἀναγκαίων κεκτημένων. Κοίτη γὰρ ἡ κίστη, εἰς ἣν τὰ βρωτὰ ἐμβάλλουσιν, κάνης δέ, τὸ ἐπίθεμα τῆς κοίτης ; Ar.~\textit{Vesp.}~529 ἐνεγκάτω μοι δεῦρο τὴν κίστην τις ὡς τάχιστα ; Ar.~\textit{Pax}~666 σπονδῶν φέρουσα τῇ πόλει κίστην πλέαν ; Thphr.~\textit{HP}~3.10.4 ἐξ οὗ τὰς κίστας ποιοῦσιν ; Thphr.~\textit{HP}~3.13.1 δι᾽ ὃ καὶ τὰς κίστας ἐξ αὐτοῦ ποιοῦσιν ὥσπερ καὶ ἐκ τοῦ τῆς φιλύρας ; Thphr.~\textit{HP}~5.7.7 ἰτέᾳ δὲ πρός τε τὰς ἀσπίδας καὶ τὰς κίστας καὶ τὰ κανᾶ ; Hsch.~α~5457 ἀντίπηγα· κίστην. Εὐριπίδης Ἴωνι (1338). ἤγουν ἀγγεῖον ; Poll.~7.79 Εἰς ἃ δὲ κατετίθεντο τὰς ἐσθῆτας ταύτας, χηλοὶ μὲν καθ᾽ Ὅμηρον, κοῖταί τε καὶ κιβωτοί, καὶ κίσται, καὶ ζυγάστρια}
\lp{284 κάθελε, κᾆτ᾽ ἔξελε]}{Ar.~\textit{Ach.}~244 κατάθου τὸ κανοῦν ὦ θύγατερ, ἵν᾽ ἀπαρξώμεθα ; Eur.~\textit{El.}~140–141 θὲς τόδε τεῦχος ἐμῆς ἀπὸ κρατὸς ἑ- / λοῦσ᾽ ; Ar.~\textit{Vesp.}~936 αὐτὸς καθελοῦ· τοὺς μάρτυρας γὰρ ἐσκαλῶ ; Ar.~\textit{Pax}~1145 τῶν τε πυρῶν μεῖξον αὐτοῖς, τῶν τε σύκων ἔξελε ; Ar.~\textit{Eccl.}~222 ἐπὶ τῆς κεφαλῆς φέρουσιν ὥσπερ καὶ πρὸ τοῦ ; Ar.~\textit{Plut.}~1198 ἱδρυσόμεθα, λαβοῦσ᾽ ἐπὶ τῆς κεφαλῆς φέρε ; Ar.~fr.~708~K.-A. (=~686~Kock) τὸν σαπέρδην ἀποτῖλαι χρὴ / κᾆτ᾽ ἐκπλῦναι καὶ διαπλῦναι ; Ar.~\textit{Lys.}~1279 πρόσαγε χορόν, ἔπαγε ⟨δὲ⟩ Χάριτας}
\lp{284–285 τὴν κίστην … τὰ πόπαν᾽]}{Men.~\textit{Dysc.}~447–450 (Ath.~4.146e) κοίτας φέροντες σταμνία ⟨τ᾽⟩ οὐχὶ τῶν θεῶν / ἕνεκ᾽, ἀλλ᾽ ἑαυτῶν. ὁ λιβανωτὸς εὐσεβὲς / καὶ τὸ πόπανον· τοῦτ᾽ ἔλαβεν ὁ θεὸς ἐπὶ τὸ πῦρ ; Theoc.~26.7 ἱερὰ δ᾽ ἐκ κίστας ποπανεύματα χερσὶν ἑλοῖσαι}
\lp{285 τὰ πόπαν᾽]}{\textit{Test.} Eust.~437.4 καὶ ἐκ τοῦ πόπος καὶ ἄναξ, θείων λέξεων, ὡς προδεδήλωται, πόπανον τὸ θεῷ ἄνακτι καθοσιούμενον, ὡς ἐμφαίνει καὶ ὁ κωμικός ; Ar.~\textit{Plut.}~660 ἐπεὶ δὲ βωμῷ πόπανα καὶ προθύματα ; \textit{idem} Suid.~κ~936 (praem. Ἀριστοφάνης Πλούτῳ ; καθωσιώθη pro καὶ προθύματα ; add. κατεκλίνομεν τὸν Πλοῦτον) ; Hsch.~π~3004 πόπανα· πλακούντια ἀπὸ ἄρτου (Ar. Eccl. 843. Plut. 660) ; Phot.~π~1090 πόπανα· πλακούντια πλατέα καὶ λεπτὰ καὶ περιφερῆ ; \textit{idem} Suid.~π~2051 (add. ὅτι ἐπὶ τέσσαρσι ποπάνοις ἔθυον βοῦν καὶ ἐκάλουν αὐτὸν πέμπτον βοῦν. Πόπανα εἰδικόν, προθύματα γενικόν) ; Ar.~\textit{Ach.}~245–246 ὦ μῆτερ ἀνάδος δεῦρο τὴν ἐτνήρυσιν, / ἵν᾽ ἔτνος καταχέω τοὐλατῆρος τουτουί ; Ar.~\textit{Eccl.}~843 λαγῷ᾽ ἀναπηγνύασι, πόπανα πέττεται ; Ar.~\textit{Plut.}~680 εἴ που πόπανον εἴη τι καταλελειμμένον ; Pl.~\textit{R.}~455c τήν τε ὑφαντικὴν λέγοντες καὶ τὴν τῶν ποπάνων τε καὶ ἑψημάτων θεραπείαν ; Ael.~\textit{VH}~11.5 ἐς τὰ κανᾶ, ἔνθα ἦν αὐτοῖς ὅ τε λιβανωτὸς καὶ τὰ πόπανα ; D.L.~8.13 διὰ τὸ πυροὺς καὶ κριθὰς καὶ πόπανα μόνα τίθεσθαι ἐπ᾽ αὐτοῦ ἄνευ πυρός, ἱερεῖον δὲ μηδέν, ὥς φησιν Ἀριστοτέλης ἐν Δηλίων πολιτείᾳ ; Clem.Al.~\textit{Protr.}~2 οὐ σησαμαῖ ταῦτα καὶ πυραμίδες καὶ τολύπαι καὶ πόπανα πολυόμφαλα ; Porph.~\textit{Abst.}~2.19 οἱ δὲ τὰ περὶ τῶν ἱερουργιῶν γεγραφότες θυσιῶν τὴν περὶ τὰ πόπανα ἀκρίβειαν φυλάττειν ; Hsch.~α~2892 ἀλετρίδες· γίγνονταί τινες ⟨τῶν εὖ γεγονυιῶν⟩ παρθένοι, αἵτινες τὰ εἰς τὰς θυσίας πόπανα ἀλοῦσι, καὶ ἔστι τοῦτο ἔντιμον (Ar. Lys. 643) ; Hsch.~β~967 βοῦς· πόπανόν τι τῶν θυομένων οὕτως ἐν ταῖς ἁγιωτάταις Ἀθήνησι θυσίαις. ἦν δὲ βοῒ παραπλήσιον ; Hsch.~ε~1888 ἐλατήρ· τὸ πλατὺ πόπανον, ἀπὸ τοῦ ἐληλάσθαι εἰς μέγεθος, ἢ πέμμα. καὶ ἡνίοχος (Δ 145) ; Moer. p.~212,~1–2~Bekker φθοῖς Ἀττικοὶ … ἔστι δὲ πέμμα πλατὺ ἔχον ὀμφαλόν. πόπανον Ἕλληνες ; Hsch.~β~1004 βουφόνια· … φασὶ βοῦν καταφαγεῖν τὸ παρεσκευασμένον πόπανον τῇ θυσίᾳ, οἷον πλακούντιον ἐξ ἄρτου (Ar. Nub. 985) ; Hsch.~ε~2430 ἐμπέλανα· πόπανα ; Hsch.~ε~5054 ἐπιπελάνιαι· ὀλαί. καὶ πόπανα ; Hsch.~ε~5154 ἐπισέληνα· πόπανα μηνοειδῆ (Plat. com. fr. 188, 10 K.–A.) ; Hsch.~π~1281 πέλα[ι]να· πόπανα, μειλίγματα ; Hsch.~σ~379 σελήνας· πόπανα τῷ ἄστρῳ ὅμοια, πέμματα (Eur. fr. 350 N.) ; Hsch.~σ~719 σίρβηνον· πόπανόν τι, ὃ παρετίθετο τῇ Ἀφροδίτῃ ; Hsch.~χ~196 χάρις· … καὶ θῦμα ἐκ τριῶν ποπάνων συγκείμενον ; Suid.~α~2082 πέλανοι δὲ τὰ εἰς θεοὺς πέμματα· πόπανα δέ ; Suid.~ο~907 καὶ τὰς μὲν κριθὰς μέχρι νῦν ὅλας χέουσιν οἱ ἐπιθύοντες ταῖς σπονδαῖς, ἐπεὶ σύμβολον τῆς παλαιᾶς τροφῆς· τὰ δὲ πόπανα τῆς ἄρτι ὥρας}
\lp{285 ταῖν θεαῖν]}{Σ Ar.~\textit{Thesm.}~566 ὡς τὼ χεῖρε. οὐκέτι δὲ τοῖν θεοῖν, ἀλλὰ ταῖν θεαῖν ; And.~1.31 μεμύησθε καὶ ἑοράκατε τοῖν θεοῖν τὰ ἱερά ; Ar.~\textit{Thesm.}~948 ὅταν ὄργια σεμνὰ θεοῖν ἱεραῖς ὥραις ἀνέχωμεν}
\lp{286 δέσποινα πολυτίμητε Δήμητερ φίλη]}{Ar.~\textit{Thesm.}~594 ὦ πολυτιμήτω θεώ ; Ar.~\textit{Ran.}~337 ὦ πότνια πολυτίμητε Δήμητρος κόρη ; Ar.~\textit{Ran.}~386–387 Δήμητερ ἁγνῶν ὀργίων / ἄνασσα συμπαραστάτει ; Ar.~\textit{Ran.}~886–887 Δήμητερ ἡ θρέψασα τὴν ἐμὴν φρένα, / εἶναί με τῶν σῶν ἄξιον μυστηρίων ; Philem.~fr.~70~K.-A. (=~67~Kock) (Ath.~14.645a) Ἄρτεμι, φίλη δέσποινα, τοῦτόν σοι φέρω, / ὦ πότνι᾽, ἀμφιφῶντα καὶ σπονδήσιμα ; Ar.~\textit{Nub.}~269 ἔλθετε δῆτ᾽ ὦ πολυτίμητοι Νεφέλαι τῷδ᾽ εἰς ἐπίδειξιν ; Ar.~\textit{Pax}~271 εὖ γ᾽ ὦ πότνια δέσποιν᾽ Ἀθηναία ποιῶν ; Ar.~\textit{Pax}~1016 ταῦτ᾽ ὦ πολυτίμητ᾽ εὐχομένοις ἡμῖν δίδου ; Ar.~\textit{Lys.}~203 δέσποινα Πειθοῖ καὶ κύλιξ φιλοτησία ; Ar.~\textit{Lys.}~317 δέσποινα Νίκη ξυγγενοῦ τῶν τ᾽ ἐν πόλει γυναικῶν ; Ar.~\textit{Ran.}~324 Ἴακχ᾽ ὦ πολυτίμητ᾽ ἐν ἕδραις ἐνθάδε ναίων ; Ar.~\textit{Ran.}~399 Ἴακχε πολυτίμητε, μέλος ἑορτῆς ; Eur.~\textit{Hipp.}~88 ἄναξ — θεοὺς γὰρ δεσπότας καλεῖν χρεών — ; Eur.~\textit{Hipp.}~117 δέσποινα Κύπρι ; Eur.~\textit{Supp.}~1 Δήμητερ ἑστιοῦχ᾽ Ἐλευσῖνος χθονὸς ; Aesch.~fr.~388~Radt δέσποιν᾽ Ἑκάτη, τῶν βασιλείων / πρόδομος μελάθρων ; Trag. adesp. fr.~328d Kn.-Sn. (Phot.~α~1305) ὦ δέσποτ᾽ Ἀμφιάραε πολυτίμητ᾽ ἄναξ ; Anacr.~\textit{PMG}~348.1–3 (Heph. Poem.~4.8) γουνοῦμαί σ᾽, ἐλαφηβόλε, ξανθὴ παῖ Διός, ἀγρίων δέσποιν᾽ Ἄρτεμι θηρῶν ; Ar.~\textit{Ach.}~807 ὦ πολυτίμηθ᾽ Ἡράκλεις ; Ar.~\textit{Av.}~877 δέσποινα Κυβέλη, στροῦθε, μῆτερ Κλεοκρίτου ; Ar.~\textit{Thesm.}~987–988 κισσοφόρε Βάκχειε / δέσποτ᾽ ; Antiph.~fr.~143.2~K.-A. (=~145~Kock) πλακοῦντα χρηστόν, ὦ πολυτίμητοι θεοί ; Eub.~fr.~115.6~K.-A. (=~116–117~Kock) ὦ Ζεῦ πολυτίμητ᾽, εἶτ᾽ ἐγὼ κακῶς ποτε / ἐρῶ γυναῖκας ; Chariclid.~fr.~1~K.-A. (=~1~Kock) (Ath.~7.325d) δέσποιν᾽ Ἑκάτη τριοδῖτι, / τρίμορφε, τριπρόσωπε, / τρίγλαις κηλευμένα ; Hsch.~δ~707 δεσποίνας· γυναῖκας. Θεσσαλοί ; Hsch.~π~2930 πολυτιμητίζειν· πολλοὺς τιμᾶν. ἢ ὑπὸ πολλῶν τιμᾶσθαι ; Phot.~π~1063 πολυτιμητιζόμενος· θεραπευόμενος}
\lp{287 Φερρέφαττα]}{Ar.~\textit{Ran.}~671 χἠ Φερρέφατθ᾽, ἅτ᾽ ὄντε κἀκείνω θεώ ; Pl.~\textit{Cra.}~404c Φερρέφαττα δέ· πολλοὶ μὲν καὶ τοῦτο φοβοῦνται τὸ ὄνομα ; Moer. p.~212,~17~Bekker Φερέφαττα Ἀττικοί, Φερσεφόνη Ἕλληνες ; Epicr.~fr.~8.2~K.-A. (=~9~Kock) (Ael.~\textit{NA}~12.10) ἐπομνύουσα τὰν Κόραν τὰν Ἄρτεμιν τὰν Φερρέφατταν ; Aesch.~\textit{Cho.}~490 ὦ Περσέφασσα, δὸς δ᾽ ἔτ᾽ εὔμορφον κράτος ; D.~54.8 ἀναστρέφουσιν ἀπὸ τοῦ Φερρεφαττίου ; Hsch.~φ~316 Φερ⟨ρ⟩εφάττιον· τόπος ἐν ἀγορᾷ (Dem. 54, 8) ; Hsch.~φ~317 Φερσεφόνεια· ἡ τῆς Δήμητρος θυγάτηρ, ἡ φέρουσα τὸ ἄφενος, τουτέστι τὸν πλοῦτον, διὰ τὸν καρπόν, ⟨ἢ⟩ ἀπὸ τοῦ φέρειν ὄνησιν ; Eur.~\textit{Hel.}~175 Φερσέφασσα ; Soph.~\textit{Ant.}~894 πλεῖστον δέδεκται Φερσέφασσ᾽ ὀλωλότων ; \textit{h.Cer.}~56 ἥρπασε Περσεφόνην}
\lp{287 πολλὰ πολλάκις]}{Ar.~\textit{Eccl.}~1105 ὅμως δ᾽ ἐάν τι πολλὰ πολλάκις πάθω ; Eur.~\textit{Med.}~1165 δώροις ὑπερχαίρουσα, πολλὰ πολλάκις ; Eur.~\textit{Tro.}~1015 καίτοι σ᾽ ἐνουθέτουν γε πολλὰ πολλάκις ; Ar.~\textit{Eq.}~411 ἔγωγε νὴ τοὺς κονδύλους οὓς πολλὰ δὴ ᾽πὶ πολλοῖς ; Ar.~\textit{Ach.}~253 ὅπως τὸ κανοῦν καλὴ καλῶς ; D.~20.3 ἀλλὰ καὶ ψηφίσματα πολλὰ πολλάκις ἐξαπατηθέντες κεχειροτονήκατε ; [Luc.]~\textit{Dem.Enc.}~38 τῇ παραδόξῳ ῥοπῇ τῆς τύχης τῇ πολλὰ πολλάκις ἡμῖν συνειργασμένῃ}
\lp{286–288 δέσποινα … πολλάκις μέ σοι θύειν ἔχουσαν]}{Eur.~\textit{El.}~805–807 νύμφαι πετραῖαι, πολλάκις με βουθυτεῖν / καὶ τὴν κατ᾽ οἴκους Τυνδαρίδα δάμαρτ᾽ ἐμὴν / πράσσοντας ὡς νῦν, τοὺς δ᾽ ἐμοὺς ἐχθροὺς κακῶς ; Herod.~4.86–87 εἴη γάρ, ὦ μέγιστε, χὐγιίῃ πολλῇ / ἔλθοιμεν αὖτις μέζον᾽ ἱρ᾽ ἀγινεῦσαι ; Σ Ar.~\textit{Thesm.}~287 Put. ἐλλειπτικὸν τὸ σχῆμα ; Ar.~\textit{Thesm.}~950–952 πολλάκις αὐτοῖν ἐκ τῶν ὡρῶν / ἐς τὰς ὥρας ξυνεπευχόμενος}
\lp{288 θύειν ἔχουσαν]}{Herod.~4.14–17 οὐ γάρ τι πολλὴν οὐδ᾽ ἕτοιμον ἀντλεῦμεν, / ἐπεὶ τάχ᾽ ἂν βοῦν ἢ νενημένην χοῖρον / πολλῆς φορίνης, κοὐκ ἀλέκτορ᾽, ἴητρα / νούσων ἐποιεύμεσθα ; Ar.~\textit{Eq.}~1295 φασὶ ⟨μὲν⟩ γὰρ αὐτὸν ἐρεπτόμενον τὰ τῶν ἐχόντων ἀνέρων ; Ar.~\textit{Nub.}~509 χώρει· τί κυπτάζεις ἔχων περὶ τὴν θύραν ; Eur.~\textit{HF}~636 ἔχουσιν, οἳ δ᾽ οὔ ; Hes.~\textit{Op.}~336 κὰδ δύναμιν δ᾽ ἔρδειν ἱέρ᾽ ἀθανάτοισι θεοῖσιν}
\lp{288 εἰ δὲ μή, ἀλλὰ νῦν λαθεῖν]}{Ar.~\textit{Nub.}~1364 ἔπειτα δ᾽ ἐκέλευσ᾽ αὐτὸν ἀλλὰ μυρρίνην λαβόντα ; Ar.~\textit{Av.}~1598 ἐάν τι δίκαιον ἀλλὰ νῦν ἐθέλητε δρᾶν ; Pi.~\textit{O.}~1.64 εἰ δὲ θεὸν ἀνήρ τις ἔλπεταί τι λαθέμεν ἔρδων, ἁμαρτάνει ; Hor.~\textit{Ep.}~1.16.60–62 labra movet metuens audiri: “pulchra Laverna, / da mihi fallere, da iusto sanctoque videri, / noctem peccatis et fraudibus obice nubem.” ; X.~\textit{Mem.}~3.3.7 εἰ δὲ μή, ἀλλὰ νῦν γε πειράσομαι, ἔφη ; D.L.~1.36 ἠρώτησέ τις αὐτὸν εἰ λήθοι θεοὺς ἄνθρωπος ἀδικῶν· ἀλλ᾽ οὐδὲ διανοούμενος, ἔφη ; Tosi n°~681 Si non caste, saltim caute ; Ar.~\textit{Ach.}~250–252 ἀγαγεῖν τυχηρῶς τὰ κατ᾽ ἀγροὺς Διονύσια, … τὰς σπονδὰς δέ μοι / καλῶς ξυνενεγκεῖν ; Ar.~\textit{Thesm.}~419 αὐταῖς ταμιεῦσαι καὶ προαιρούσαις λαθεῖν ; Ar.~\textit{Ran.}~103 σὲ δὲ ταῦτ᾽ ἀρέσκει; — μἀλλὰ πλεῖν ἢ μαίνομαι ; Ar.~\textit{Ran.}~388–390 καὶ σῷζε τὸν σαυτῆς χορόν, / καί μ᾽ ἀσφαλῶς πανήμερον / παῖσαί τε καὶ χορεῦσαι ; Ar.~\textit{Eccl.}~26 ἢ θαἰμάτια τἀνδρεῖα κλεψάσαις λαθεῖν ; Eur.~\textit{Supp.}~3 εὐδαιμονεῖν με Θησέα τε παῖδ᾽ ἐμὸν ; Aesch.~\textit{Sept.}~253 θεοὶ πολῖται, μή με δουλείας τυχεῖν ; Ar.~\textit{Eq.}~74 ἀλλ᾽ οὐχ οἷόν τε τὸν Παφλαγόν᾽ οὐδὲν λαθεῖν ; Ar.~\textit{Thesm.}~408 οὐδὲ τοῦτ᾽ ἔστιν λαθεῖν ; Ar.~\textit{Eccl.}~23 δεῖ τὰς ἑτέρας πως κἀγκαθεζομένας λαθεῖν ; Pi.~\textit{P.}~1.67–68 Ζεῦ τέλει᾽, αἰεὶ δὲ τοιαύταν Ἀμένα παρ᾽ ὕδωρ / αἶσαν ἀστοῖς καὶ βασιλεῦσιν διακρίνειν ; Tosi n°~244 λάθε βιώσας}
\lp{289 τὴν θυγατέρα Χοιρίλ᾽]}{Vit. Eur.~1 γυναῖκα δὲ γῆμαι πρώτην Μελιτὼ, δευτέραν δὲ Χοιρίλην ; Vit. Eur.~2 λέγουσι δὲ αὐτὸν γήμαντα τὴν Μνησιλόχου θυγατέρα Χοιρίλην ; Vit. Eur.~5 ἔγημε δὲ ὅμως πρώτην μὲν Χοιρίλην, θυγατέρα Μνησιλόχου, ἐξ ἧς ἔσχε Μνησίλοχον καὶ Μνησαρχίδην καὶ Εὐριπίδην ; \textit{idem} Suid.~ε~3695 (Χοιρίνην ; add. ἀπωσάμενος δὲ ταύτην ἔσχε καὶ δευτέραν, καὶ ταύτης ὁμοίως ἀκολάστου πειραθείς) ; Σ Ar.~\textit{Thesm.}~289 ὡς τῆς θυγατρὸς αὐτοῦ κατ᾽ ἐπίκλησιν οὕτω καλουμένης. οἷον βοίδιον ἢ χρυσίον ἢ μύρτιον. ἴσως δὲ τὸ αἰδοῖον τὸ γυναικεῖον αἰνίττεται ; Σ Eur.~\textit{Hec.}~3 Φιλόχορος μὲν γὰρ ἐν τῷ περὶ τραγῳδιῶν συγγράμματι Χοιρίλην αὐτήν φησι καλεῖσθαι, ἴσως δὲ διὰ τὸ πολύπαιδα γεγενῆσθαι· ἡ γὰρ χοῖρος πολλὰ τίκτει ; Hsch.~ε~220 ἐγκεχοιριλωμένην· λεπράν (com. ad. fr. 322 K.–A.) ; Hsch.~ε~1439 ἐκκεχοιριλωμένη· οὐ Χοιρίλου οὖσα. Ἐκφαντίδῃ γὰρ τῷ κωμικῷ (Τ 5 K.–A.) Χοιρίλος θεράπων ἦν, ὃς συνεποιεῖτο κωμωδίας ; Ar.~\textit{Lys.}~930 δεῦρό νυν ὦ χρύσιον ; Ar.~\textit{Ran.}~1408 αὐτὸς τὰ παιδί᾽ ἡ γυνὴ Κηφισοφῶν ; Telecl.~fr.~41.1~K.-A. (=~39~Kock) (Vit. Eur.) Μνησίλοχός ἐστ᾽ ἐκεῖνος, ὃς φρύγει τι δρᾶμα καινὸν ; Luc.~\textit{DMeretr.}~2.2 ὦ Μύρτιον ; Gell.~15.20.6 sive quod duas simul uxores habuerat ; Σ Luc.~\textit{DMeretr.}~2 ⸢Μύρτιον καὶ Μου⸣σάριον ⸢καὶ Γλυκέριον καὶ εἴ⸣ τι τοι⸢οῦτον Ἀττικοί⸣ εἰσιν ⸢ὑποκορισμοί⸣ ; Suid.~χ~593 Χοιρίλη· ἡ Ἑκάβη}
\lp{289 Χοιρίλ᾽ (χοῖρον R)]}{Ar.~\textit{Ach.}~777 φώνει δὴ τὺ ταχέως χοιρίον ; Ar.~\textit{Vesp.}~573 εἰ δ᾽ αὖ τοῖς χοιριδίοις χαίρω, θυγατρὸς φωνῇ με πιθέσθαι ; Ar.~\textit{Vesp.}~1353 λυσάμενος ἕξω παλλακὴν ὦ χοιρίον ; Ar.~\textit{Thesm.}~538 ταύτης ἀποψιλώσομεν τὸν χοῖρον, ἵνα διδαχθῇ ; Suid.~χ~601 Χοῖρος: τὸ ζῷον. παρὰ Κορινθίοις δὲ τὸ γυναικεῖον αἰδοῖον. ἔνθεν καὶ παροιμία· Ἀκροκορινθία ἔοικας χοιροπωλήσειν. ἀντὶ τοῦ ἔοικας μισθαρνήσειν ἐν Κορίνθῳ· πολλαὶ γὰρ ἐκεῖ ἑταῖραι ; Ar.~\textit{Ach.}~781–782 αὕτα ᾽στὶ χοῖρος; — νῦν γε χοῖρος φαίνεται. / ἀτὰρ ἐκτραφείς γε κύσθος ἔσται ; Ar.~\textit{Thesm.}~540 μὴ δῆτα τόν γε χοῖρον ὦ γυναῖκες ; Ar.~\textit{Eccl.}~724 κατωνάκην τὸν χοῖρον ἀποτετιλμένας ; Hsch.~χ~601 χοιροτροφεῖον· περίζωμα γυναικεῖον. καὶ συφεόν (Eupol. fr. dub. 493. Phryn. com. fr. 45 K.–A.) ; Plu.~\textit{Prov.Alex.}~1.92 Ἀκροκορινθία ἔοικας χοιροπωλήσειν: ἐπὶ τῶν παρ᾽ ὥραν θρυπτομένων γυναικῶν· οἷον, ἔοικας μισθαρνήσειν ἐν Κορίνθῳ· τὸ γὰρ γυναικεῖον μόριον χοῖρος λέγεται ; Ar.~\textit{Vesp.}~1364 ὦ οὗτος οὗτος τυφεδανὲ καὶ χοιρόθλιψ ; Ar.~\textit{Plut.}~307–308 ὑμεῖς δὲ γρυλίζοντες ὑπὸ φιληδίας / ἕπεσθε μητρὶ χοῖροι}
\lp{289–290 ἀνδρός μοι τυχεῖν πλουτοῦντος]}{Eur.~\textit{Alc.}~165–166 τέκν᾽ ὀρφανεῦσαι τἀμά· καὶ τῷ μὲν φίλην / σύζευξον ἄλοχον, τῇ δὲ γενναῖον πόσιν ; Thgn.~187–190 οὐδὲ γυνὴ κακοῦ ἀνδρὸς ἀναίνεται εἶναι ἄκοιτις / πλουσίου, ἀλλ᾽ ἀφνεὸν βούλεται ἀντ᾽ ἀγαθοῦ. / χρήματα γὰρ τιμῶσι· καὶ ἐκ κακοῦ ἐσθλὸς ἔγημεν, / καὶ κακὸς ἐξ ἀγαθοῦ· πλοῦτος ἔμιξε γένος ; Phoenix fr.~2.10–11 Powell (Ath.~8.359f) θεοί, γένοιτο πάντ᾽ ἄμεμπτος ἡ κούρη / κἀφνειὸν ἄνδρα κὠνομαστὸν ἐξεύροι ; \textit{Od.}~6.180–181 σοὶ δὲ θεοὶ τόσα δοῖεν ὅσα φρεσὶ σῇσι μενοινᾷς, / ἄνδρα τε καὶ οἶκον}
\lp{290 πλουτοῦντος … ἠλιθίου κἀβελτέρου]}{Ar.~\textit{Eq.}~264–265 καὶ σκοπεῖς γε τῶν πολιτῶν ὅστις ἐστὶν ἀμνοκῶν, / πλούσιος καὶ μὴ πονηρὸς καὶ τρέμων τὰ πράγματα ; Ar.~\textit{Nub.}~842 γνώσει δὲ σαυτὸν ὡς ἀμαθὴς εἶ καὶ παχύς ; Eup.~fr.~172.7–8~K.-A. (=~159~Kock) (Ath.~6.236e–f) ἐκεῖ δ᾽ ἐπειδὰν κατίδω τιν᾽ ἄνδρα / ἠλίθιον, πλουτοῦντα δ᾽, εὐθὺς περὶ τοῦτον εἰμί ; Σ Ar.~\textit{Plut.}~179 κωμῳδεῖται δὲ ὡς εὔπορος καὶ ὡς μέγας τῷ σώματι καὶ ἠλίθιος. Νικοχάρης Γαλατείᾳ· τί δῆτ᾽ ἀπαιδευτότερος εἶ Φιλωνίδου τοῦ Μελιτέως ; Σ Ar.~\textit{Thesm.}~290 ἀναισθήτου καὶ μωροῦ, ἵνα κρατῇ αὐτοῦ ; Suid.~π~830 Παχεῖς: Ἀττικοὶ τοὺς πλουσίους καλοῦσι συνήθως. Ἀριστοφάνης· ὡς ἀμαθὴς εἶ καὶ παχύς. ἀντὶ τοῦ ἀναίσθητος καὶ παχὺς εἰς τὸ νοεῖν ; D.L.~6.47 τὸν ἀμαθῆ πλούσιον πρόβατον εἶπε χρυσόμαλλον ; Ar.~\textit{Thesm.}~413 δέσποινα γὰρ γέροντι νυμφίῳ γυνή ; Prop.~2.16.7–8 quare, si sapis, oblatas ne desere messis / et stolidum pleno vellere carpe pecus ; \textit{idem} Tosi n°~750 ; Σ Ar.~\textit{Plut.}~179 Φιλωνίδης ἀνήρ τις ἦν πλούσιος, οὗτινος ἀμόρφου ὄντος, διὰ τὰ χρήματα ἐρᾷ Λαῒς αὐτοῦ, ἡ ἐν Κορίνθῳ ἀδομένη πόρνη ; Hsch.~π~1165 πάχητες· πλούσιοι. παχεῖς ; Hsch.~π~1171 παχύνοοι· παχὺν νοῦν ἔχοντες, ἀνόητοι ; Apostol.~3.30 Ἀπαιδευτότερος Φιλωνίδου τοῦ Μελιτέως: οὗτος ὁ Φιλωνίδης οὐ μόνον μέγας ἦν, ἀλλὰ καὶ ἀμαθὴς καὶ ὑώδης}
\lp{290 ἄλλως τ᾽]}{Ar.~\textit{Av.}~1476–1477 χρήσιμον μὲν οὐδέν, ἄλλως / δὲ δειλὸν καὶ μέγα ; Ar.~\textit{Thesm.}~159 ἄλλως τ᾽ ἄμουσόν ἐστι ποιητὴν ἰδεῖν ; Ar.~\textit{Plut.}~975–976 ἦν μοί τι μειράκιον φίλον, / πενιχρὸν μέν, ἄλλως δ᾽ εὐπρόσωπον καὶ καλὸν}
\lp{290 ἠλιθίου]}{Hsch.~η~359 ἠλίθιον· μάταιον, ἀναίσθητον, μωρόν, ἀνόητον, ἄφρονα ; Ar.~\textit{Thesm.}~593 ἠλίθιος ὅστις τιλλόμενος ἠνείχετο ; Ar.~\textit{Ran.}~917 ἠλίθιος γὰρ ἦσθα ; Cratin.~fr.~45~K.-A. (=~43~Kock) (EM~196,~7) ὁ δ᾽ ἠλίθιος ὥσπερ πρόβατον βῆ βῆ λέγων βαδίζει ; Ar.~\textit{Ach.}~443 τοὺς δ᾽ αὖ χορευτὰς ἠλιθίους παρεστάναι ; Alc. Com. fr.~1~K.-A. (=~1~Kock) (Ath.~7.316b) ἠλίθιον εἶναι νοῦν τε πουλύποδος ἔχειν ; Eur.~\textit{Cyc.}~537 ἠλίθιος ὅστις μὴ πιὼν κῶμον φιλεῖ ; Aesch.~\textit{Ag.}~365–366 μήτε πρὸ καιροῦ μήθ᾽ ὑπὲρ ἄστρων / βέλος ἠλίθιον σκήψειεν ; Simon.~\textit{PMG}~542.37 (Pl.~\textit{Prt.}~346c) τῶν γὰρ ἠλιθίων ἀπείρων γενέθλα ; Pi.~\textit{P.}~3.11 χόλος δ᾽ οὐκ ἀλίθιος ; Hdt.~1.60.3 εὐηθείης ἠλιθίου ἀπηλλαγμένον ; Pl.~\textit{Chrm.}~162b ἤ τινος ἠλιθίου ἤκουσας τουτὶ λέγοντος ; X.~\textit{Mem.}~3.9.8 τοὺς ἠλιθίους δὲ ἀεὶ πάσχειν αὐτό ; Lys.~10.16 οὐκ ἂν ἠλίθιον αὐτὸν νομίζοιεν}
\lp{290 κἀβελτέρου]}{Ar.~\textit{Ran.}~989–991 τέως δ᾽ ἀβελτερώτατοι / κεχηνότες Μαμμάκυθοι / Μελιτίδαι καθῆντο ; Ar.~\textit{Eccl.}~768 μάλιστα πάντων. — τὸν μὲν οὖν ἀβέλτερον ; Pl. Com. fr.~65.2–3~K.-A. (=~64~Kock) (Ath.~2.68c) ὁ μὲν Λέαγρος, Γλαύκωνος ὢν μεγάλου γένους, / ⟨ἀβελτερο⟩κόκκυξ ἠλίθιος περιέρχεται ; Hsch.~α~108 ἀβέλτερος· ἀνόητος, μωρός, ὁ τὸ βέλτιον μὴ γινώσκων, ἠλίθιος ; Ar.~\textit{Nub.}~1201 τί κάθησθ᾽ ἀβέλτεροι ; Pl.~\textit{Smp.}~198d ἐγὼ μὲν γὰρ ὑπ᾽ ἀβελτερίας ᾤμην δεῖν τἀληθῆ λέγειν ; Suid.~α~32 Ἀβέλτερος. ἀνόητος, ἀσύνετος. βέλτερος γὰρ ὁ φρόνιμος ; Suid.~μ~1344 Μωροὶ καὶ ἀβέλτεροι καὶ κεχηνότες Μαμμάκυθος: Μελιτίδης, Βουταλίων, Κόροιβος ; Anaxandr.~fr.~22~K.-A. (=~21~Kock) (Stob. Ecl.~2.1.3) ἅπαντές ἐσμεν πρὸς τὰ θεῖ᾽ ἀβέλτεροι, / κοὐκ ἴσμεν οὐδέν ; Pl.~\textit{R.}~409c ἀβέλτερος αὖ φαίνεται, ἀπιστῶν παρὰ καιρὸν ; D.~9.14 ἀβελτερώτατος εἴη πάντων ἀνθρώπων ; Hsch.~α~106 ἀβελτερία· ἀφροσύνη, ἄνοια ; Suid.~α~31 Ἀβελτεροκόκκυξ: ὁ κενὸς καὶ ἀβέλτερος}
\lp{290 ἠλιθίου κἀβελτέρου (sots proverbiaux)]}{Zen.~4.58 Κοροίβου ἠλιθιώτερος: εὐήθης καὶ μωρός. Ἐπὶ γὰρ τοῦ μωραίνοντος ἔταττον τὸν Κόροιβον ; \textit{idem} Hsch.~κ~3649 (Κόροιβος· ἠλίθιος καὶ μωρός pro Κοροίβου ἠλιθιώτερος· εὐήθης καὶ μωρός ; add. ἀπὸ Κοροίβου τινὸς μωροῦ καὶ ἠλιθίου μετάγοντες, ὃν οἴονται τὸν Μυγδόνος εἶναι παῖδα τοῦ Φρυγός) ; Suid.~γ~118 Γελοιότερον Μελιτίδου, ἐπὶ τῶν ἐπὶ μωρίᾳ διαβεβλημένων. Μελιτίδης γὰρ ἀνὴρ κωμῳδούμενος ὑπὸ τῶν ποιητῶν ἐπὶ μωρίᾳ … γήμαντα δὲ τῆς νύμφης μὴ ἅψασθαι· φοβεῖσθαι γὰρ μὴ αὐτὸν ἡ παῖς τῇ μητρὶ διαβάλλῃ ; Σ Luc.~\textit{Philops.}~3 ὁ μὲν Κόροιβος μωρός τις ἀπομνημονεύεται ἄνθρωπος καὶ οὕτως ὥστε γυναῖκα ἀγαγόμενος μὴ ἂν ἑλέσθαι συγκαθευδῆσαι αὐτῇ διὰ τὸν πρὸς τὴν πενθερὰν αὐτοῦ φόβον ; Suid.~μ~187 Μαργίτης: ἀνὴρ ἐπὶ μωρίᾳ κωμῳδούμενος. ὅν φασιν ἀριθμῆσαι μὲν μὴ πλείω τῶν ε΄ δυνηθῆναι· νύμφην δὲ ἀγόμενον μὴ ἅψασθαι αὐτῆς, ἀλλὰ φοβεῖσθαι λέγοντα, μὴ τῇ μητρὶ αὐτὸν διαβάλῃ ; Diogenian.~5.12 Ἡλιθιώτερος τῆς Πραξίλλης … Ὁμοία τῇ, Ἀνοητότερος Ἰβύκου, καὶ Κοροίβου, καὶ Μελιτίδου ; Diogenian.~5.56 Κοροίβου ἠλιθιώτερος: οὗτος γὰρ τὰ κύματα τῆς θαλάσσης ἠρίθμει ; Ael.~\textit{VH}~13.15 καὶ Κόροιβον δὲ καὶ Μελιτίδην καὶ ἐκείνους ἀνοήτους φασίν ; Suid.~β~468 Βουταλίων καὶ Κόροιβος καὶ Μελιτίδης ἐπὶ μωρίᾳ διεβέβληντο ; App. Prov.~3.10 Ἡλιθιώτερος Λευκολόφου: οὗτος ἀποδόμενος τὴν οἰκίαν, ἀντεποιεῖτο τοῦ φρέατος}
\lp{291 πρὸς θυλάκιον]}{\textit{Test.} Poll.~10.151–152 (Bethe II, p.~235) καὶ θυλάκιον δὲ καὶ θυλακίσκον· Ἀριστοφάνης γοῦν ἐν Τριφάλητι (fg 545) τοῦτο ὑποδηλοῖ, ὅταν φῇ … ἐπὶ τοῦ αὐτοῦ δὲ καὶ σακίον, ὅταν φῇ ἐν Θεσμοφοριαζούσαις (fg 328)· σακίον, ἐν οἷσπερ τἀργύριον ταμιεύεται ; Ar.~\textit{Vesp.}~314 ἀνόνητον ἄρ᾽ ὦ θυλάκιόν σ᾽ εἶχον ἄγαλμα ; Ar.~\textit{Lys.}~493–495 ὑμεῖς ταμιεύσετε τἀργύριον; — τί ⟨δὲ⟩ δεινὸν τοῦτο νομίζεις; / οὐ καὶ τἄνδον χρήματα πάντως ἡμεῖς ταμιεύομεν ὑμῖν ; Ar.~\textit{Thesm.}~418–419 ἃ δ᾽ ἦν ἡμῖν πρὸ τοῦ / αὐταῖς ταμιεῦσαι καὶ προαιρούσαις λαθεῖν ; Ar.~\textit{Ran.}~1203 καὶ κῳδάριον καὶ ληκύθιον καὶ θύλακον ; Ar.~\textit{Eccl.}~211–212 καὶ γὰρ ἐν ταῖς οἰκίαις / ταύταις ἐπιτρόποις καὶ ταμίαισι χρώμεθα ; Ar.~\textit{Eccl.}~382 μὰ τὸν Δί᾽ οὐδέν᾽ ἄλλον ἢ τὸν θύλακον ; Hsch.~θ~48 θαλλικοποιοί· οἱ τὰ σκυτούμενα κιβώτια, καὶ τοὺς δερματίνους ῥίσκους ἐργαζόμενοι ; Hsch.~θ~50 θάλλικα· σάκκου εἶδος ; Hsch.~θ~51 θαλλίς· μάρσιππος μακρός ; Suid.~μ~226 Μάρσιππος: σάκκος, θυλάκιον, σακέλλιον. εἶδόν τινας ἐν πέτραις μαρσιππίους ἱματίων κατατιθεμένους. Ξενοφῶν ; Ar.~\textit{Av.}~503 ὀβολὸν κατεβρόχθισα· κᾆτα κενὸν τὸν θύλακον οἴκαδ᾽ ἀφεῖλκον ; Hdt.~3.46.2 θύλακον δὲ φέροντες ἔφασαν τὸν θύλακον ἀλφίτων δέεσθαι. οἳ δέ σφι ὑπεκρίναντο τῷ θυλάκῳ περιεργάσθαι ; Hsch.~θ~855 θυλλίς· θύλακος. γωρυτός. ἔλυτρον ; Hsch.~κ~4884 κώρυκος· θυλάκιον. ἔστι δὲ δερμάτινον ἀγγεῖον, ὅμοιον ἀσκῷ. οἱ δὲ πλοῖον· οἱ δὲ κόγχην ; Hsch.~σ~82 σάκτας· ὁ θύλακος (Ar. Plut. 681) ; Phot. s.v.~θύλακος (Naber I, p.~285) θύλακος· ὁ ῥαπτὸς φάσκωλος ; Suid.~κ~2301 Κώρυκος: θυλάκιον. τὸ παῤ ἡμῖν βουλγίδιον· ἢ πλέγμα δεκτικὸν ἄρτων ; Plu.~\textit{Glor.\,Ath.}~4,~348a (Corinn. test.) γελάσασα ἐκείνη τῇ χειρὶ δεῖν ἔφη σπείρειν, ἀλλὰ μὴ ὅλῳ τῷ θυλάκῳ ; Apostol.~15.6 Πτίσσε πτίσσε τὸν Ἀναξάρχου θύλακον· οὐ γὰρ πτίσσεις Ἀνάξαρχον ; App. Prov.~1.37 Ἁπλοῦν θύλακον πάτει: ἐπὶ τῶν εὐκαταφρονήτων ; Pl.~\textit{Tht.}~161a ὅτι με οἴει λόγων τινὰ εἶναι θύλακον ; Alex.~fr.~88.3–5~K.-A. (=~85~Kock) (Ath.~11.470e) κατά τε τὴν παροιμίαν / αἰεὶ ποτ᾽ εὖ μὲν ἀσκός, εὖ δὲ θύλακος / ἅνθρωπός ἐστιν}
\lp{291 πρὸς θυλάκιον (Σ τὸν παιδαρίσκον)]}{Ar.~\textit{Pax}~1300 εἰπέ μοι ὦ πόσθων, ἐς τὸν σαυτοῦ πατέρ᾽ ᾄδεις ; Σ Ar.~\textit{Thesm.}~291 τὸν παιδαρίσκον. ἴσως δὲ παρὰ τὴν πόσθην αὐτὸ συνέθηκεν. ἀντὶ τοῦ κατωφερῆ τὸν νοῦν μου ποίησον ; Hsch.~π~3097 πόσθων· πόσθην τὸ ἀνδρεῖον αἰδοῖον λέγουσι (Ar. Nub. 1014). πόσθωνας δὲ παρὰ τοῦτο τοὺς παῖδας, τινὲς δὲ τοὺς ψώλωνας, ἄλλοι μωροὺς ἢ παιδαριώδεις (Ar. Pac. 1300) ; Phot.~π~1114 πόσθων (Ar. Pac. 1300)· κυρίως λέγεται παιδάριον· οὕτως γὰρ ὑποκοριζόμενοι ἔλεγον ἀπὸ τοῦ αἰδοίου· πόσθιον γὰρ τοῦτο καλοῦσιν· δηλοῖ δὲ καὶ τὸν παιδαριώδη καὶ ἄφρονα κατὰ μεταφοράν ; Ar.~\textit{Lys.}~771 τοὺς ἔποπας φεύγουσαι, ἀπόσχωνταί τε φαλήτων ; Ar.~\textit{Thesm.}~254 νὴ τὴν Ἀφροδίτην ἡδύ γ᾽ ὄζει ποσθίου ; Ar.~\textit{Thesm.}~514–515 λέων λέων σοι γέγονεν, αὔτ᾽ ἔκμαγμα σόν, / τά τ᾽ ἄλλ᾽ ἁπαξάπαντα καὶ τὸ πόσθιον ; Epilyc. Κωραλίσκος (Ath.~4.140a) Ἐπίλυκός τε ἐν Κωραλίσκῳ λέγων οὕτως ; Hsch.~π~3096 πόσθιον· αἰδοῖον ἀνδρός (Ar. Thesm. 254) ; Hdn.~\textit{Pros.\,cath.}~6 (Κωραλίσκος) Τὰ εἰς σκος ι ἢ υ παραληγόμενα παροξύνεται, σατυρίσκος, νεανίσκος, παιδίσκος, ἀστερίσκος, ὀβελίσκος, πανίσκος, Κωραλίσκος, οὐρανίσκος ; Hsch.~π~3095 πόσθη· τοῦ αἰδοίου τὸ δέρμα ; Suid.~π~2113 Πόσθη: τὸ αἰδοῖον τοῦ ἀνδρός}
\lp{291 νοῦν ἔχειν … καὶ φρένας]}{Ar.~\textit{Lys.}~432 οὐ γὰρ μοχλῶν δεῖ μᾶλλον ἢ νοῦ καὶ φρενῶν ; Ar.~\textit{Ran.}~534a νοῦν ἔχοντος καὶ φρένας ; Men.~fr.~58~K.-A. (=~60~Kock) (Stob.~4.24e.37–38) εὐδαιμονία τοῦτ᾽ ἐστὶν υἱὸς νοῦν ἔχων· / ἀλλὰ θυγάτηρ κτῆμ᾽ ἐστὶν ἐργῶδες πατρί ; \textit{idem} Apostol.~8.8a (Εὐδαιμονίας ; add. Μενάνδρου ; om. v. 2) ; Hor.~\textit{Ep.}~1.4.8–11 quid voveat dulci nutricula maius alumno, / qui sapere et fari possit quae sentiat, et cui / gratia, fama, valetudo contingat abunde, / et mundus victus non deficiente crumina? ; D.~18.324 μὴ δῆτ᾽, ὦ πάντες θεοί, μηδεὶς ταῦθ᾽ ὑμῶν ἐπινεύσειεν, ἀλλὰ μάλιστα μὲν καὶ τούτοις βελτίω τινὰ νοῦν καὶ φρένας ἐνθείητε ; Ar.~\textit{Thesm.}~462 οὐκ ἄκαιρα, φρένας ἔχουσα ; D.~25.33 ἀλλὰ νοῦ καὶ φρενῶν ἀγαθῶν καὶ προνοίας πολλῆς ; Ar.~\textit{Nub.}~1122 προσεχέτω τὸν νοῦν, πρὸς ἡμῶν οἷα πείσεται κακά}
\lp{292 ποῦ ποῦ]}{Ar.~\textit{Eccl.}~320–321 ἀλλ᾽ ἐν καθαρῷ ποῦ, ποῦ τις ἂν χέσας τύχοι; / ἢ πανταχοῦ τοι νυκτός ἐστιν ἐν καλῷ ; Eur.~\textit{Or.}~470 ποῦ ποῦ θυγατρὸς τῆς ἐμῆς ἴδω πόσιν ; Ar.~\textit{Eq.}~1198 ποῦ ποῦ ; Ar.~\textit{Av.}~1122 ποῦ ποῦ ᾽στι, ποῦ ποῦ ποῦ ᾽στι, ποῦ ποῦ ποῦ ᾽στι, ποῦ ; Ar.~\textit{Thesm.}~1093 ποῖ ποῖ πεύγεις}
\lp{292 καθίζωμ᾽ ἐν καλῷ]}{Ar.~\textit{Eccl.}~86–87 νὴ τὸν Δί᾽ ὥστε δεῖ σε καταλαβεῖν ἕδρας / ὑπὸ τῷ λίθῳ τῶν πρυτάνεων καταντικρύ ; Pl. Com. fr.~199.1~K.-A. (=~183~Kock) (Plut. Them.~32.5) ὁ σὸς δὲ τύμβος ἐν καλῷ κεχωσμένος ; Theoc.~15.73–74 ὠθεῦνθ᾽ ὥσπερ ὕες. — θάρσει γύναι· ἐν καλῷ εἰμές. / κεἰς ὥρας κἤπειτα φίλ᾽ ἀνδρῶν ἐν καλῷ εἴης ; Ael.~\textit{VH}~2.13 καὶ δὴ καὶ ἐν καλῷ τοῦ θεάτρου ἐκάθητο ; Alciphr.~2.17 ἄγει μέ τις λαβὼν εἰς τὸ θέατρον, οὗ καθίσας ἐν καλῷ διαφόροις ἐψυχαγωγούμην θεωρίαις ; Ar.~\textit{Ach.}~28–29 ἐγὼ δ᾽ ἀεὶ πρώτιστος εἰς ἐκκλησίαν / νοστῶν κάθημαι ; Ar.~\textit{Eq.}~750 οὐκ ἂν καθιζοίμην ἐν ἄλλῳ χωρίῳ ; Ar.~\textit{Eccl.}~98–99 ἢν δ᾽ ἐγκαθεζώμεσθα πρότεραι, λήσομεν / ξυστειλάμεναι θαἰμάτια ; Alex.~fr.~42~K.-A. (=~41~Kock) (Poll.~9.44) ἐνταῦθα περὶ τὴν ἐσχάτην δεῖ κερκίδα / ὑμᾶς καθιζούσας θεωρεῖν ὡς ξένας ; Eur.~\textit{IA}~1106 Λήδας γένεθλον, ἐν καλῷ σ᾽ ἔξω δόμων ; Aesch.~\textit{Supp.}~223 σέβεσθ᾽· ἐν ἁγνῷ δ᾽ ἑσμὸς ὡς πελειάδων ; X.~\textit{HG}~2.1.25 οὐκ ἐν καλῷ ἔφη αὐτοὺς ὁρμεῖν ; Luc.~\textit{Phal.}~1.13 ἀναθέντες δὲ τὸν ταῦρον ἐν καλῷ τοῦ ἱεροῦ ; Σ Luc.~\textit{Phal.}~1.13 ἐν καλῷ τοῦ ἱεροῦ] ἤγουν ἐν καλῷ ἱερῷ ; Hsch.~γ~252 γέγονεν ἐν καλῷ· ἐν καλλίστῳ ⟨διέπρεψεν⟩ (Cyr. in Esai. 70, 205D) ; Ar.~\textit{Eq.}~783–785 ἐπὶ ταῖσι πέτραις οὐ φροντίζει σκληρῶς σε καθήμενον οὕτως, / … κᾆτα καθίζου μαλακῶς ; Ar.~\textit{Lys.}~910–911 ποῦ γὰρ ἄν τις καὶ τάλαν / δράσειε τοῦθ᾽; — ὅπου; τὸ τοῦ Πανὸς καλόν ; Ar.~\textit{Thesm.}~663–664 εἶά νυν ἴχνευε καὶ μάτευε ταχὺ πάντ᾽, / εἴ τις ἐν τόποις ἑδραῖος ; Ar.~\textit{Eccl.}~152 λέγειν τὰ βέλτισθ᾽, ἵν᾽ ἐκαθήμην ἥσυχος ; Eur.~\textit{Heracl.}~971 οὐκ οὖν ἔτ᾽ ἐστὶν ἐν καλῷ δοῦναι δίκην ; Soph.~\textit{El.}~384 νῦν γὰρ ἐν καλῷ φρονεῖν ; Th.~5.59.4 ἀλλ᾽ ἐν καλῷ ἐδόκει ἡ μάχη ἔσεσθαι ; Plu.~\textit{Caes.}~39.1 προὐκαλεῖτο Πομπήϊον ἱδρυμένον ἐν καλῷ ; D.Chr.~64.16 καὶ τοὐρανοῦ γ̓, ὥς φασιν, ἐστὶν ἐν καλῷ ; Ael.~\textit{NA}~14.25 ἔρχεται δὲ ἔνθα οἱ δοκεῖ καλῶς ἔχειν ἑαυτὸν καθίσαι καὶ ἐν καλῷ τῆς ἄγρας εἶναι πεπίστευκε}
\lp{292 τῶν ῥητόρων]}{\textit{Test.} Phot.~ρ~106 ῥήτορα· καὶ τὴν γυναῖκα. Ἀριστοφάνης (Thesm. 292) ; Ar.~\textit{Eccl.}~244 ἔπειτ᾽ ἀκούουσ᾽ ἐξέμαθον τῶν ῥητόρων ; Hsch.~ρ~281 ῥήτορες· λέκται, οὐχ οἱ συνήγοροι μόνοι ; Ar.~\textit{Ach.}~38 βοᾶν ὑποκρούειν λοιδορεῖν τοὺς ῥήτορας ; Ar.~\textit{Thesm.}~382 ὅπερ ποιοῦσ᾽ οἱ ῥήτορες ; Ar.~\textit{Ach.}~680 ὑπὸ νεανίσκων ἐᾶτε καταγελᾶσθαι ῥητόρων ; Ar.~\textit{Thesm.}~384 λέξουσ᾽ ἀνέστην ὦ γυναῖκες ; Ar.~\textit{Eccl.}~195 τῶν δὲ ῥητόρων}
\lp{292–293 τῶν ῥητόρων ἵν᾽ ἐξακούω]}{Ar.~\textit{Eccl.}~90–91 τί γὰρ ἂν χεῖρον ἀκροῴμην ἄρα / ξαίνουσα ; Ar.~\textit{Eccl.}~588–589 πρὶν ἐπίστασθαι τὴν ἐπίνοιαν καὶ τοῦ φράζοντος ἀκοῦσαι}
\lp{293 ἐξακούω]}{Eur.~\textit{Heracl.}~677 ᾐκάζομεν ταῦτ᾽· οὐ γὰρ ἐξηκούομεν ; Aesch.~\textit{Eum.}~397 πρόσωθεν ἐξήκουσα κληδόνος βοὴν ; Plu.~\textit{Cic.}~26.4 πρὸς τοῦτον ἔν τινι δίκῃ λέγοντα τοῦ Κικέρωνος μὴ ἐξακούειν ; Ar.~\textit{Av.}~1197–1198 ὡς ἐγγὺς ἤδη δαίμονος πεδαρσίου / δίνης πτερωτὸς φθόγγος ἐξακούεται ; Soph.~\textit{Phil.}~676 λόγῳ μὲν ἐξήκουσ᾽, ὄπωπα δ᾽ οὐ μάλα ; X.~\textit{Cyr.}~4.3.3 στὰς ὅπου ἔμελλον πάντες ἀκούσεσθαι τὰ βουλευόμενα ; Soph.~\textit{El.}~553 λυπηρὸν εἶτα σοῦ τάδ᾽ ἐξήκουσ᾽ ὕπο ; Soph.~\textit{Phil.}~378 δηχθεὶς πρὸς ἁξήκουσεν ὧδ᾽ ἠμείψατο ; Soph.~\textit{Phil.}~472 ὅσοισί τ᾽ ἐξήκουσας ἐνναίοντά με ; Hsch.~ε~3526 ἐξάκουστον· περιβόητον (Cyr. in Ioh. 2, 272, 1 . .)}
\lp{293 σὺ δ᾽ ἄπιθ᾽ ὦ Θρᾷττ᾽ ἐκποδών]}{Ar.~\textit{Eccl.}~134 ἄπιθ᾽ ἐκποδών· τοιαῦτ᾽ ἂν ἡμᾶς ἠργάσω ; Ar.~\textit{Ach.}~240 ἐκποδών· θύσων γὰρ ἁνὴρ ὡς ἔοικ᾽ ἐξέρχεται ; Ar.~\textit{Thesm.}~36 ἀλλ᾽ ἐκποδὼν πτήξωμεν, ὡς ἐξέρχεται ; Hsch.~ε~1594 ἐκ πάτου· ἐκποδών ; Hsch.~ε~1635 ἐκποδών· ἐκ μέσου (Eur. Hec. 1282) ⟨ … ⟩ ἐξαλείφεσθαι ; Zen.~4.33 Θύραζε Κᾶρες, οὐκ ἔτ᾽ Ἀνθεστήρια: οἱ μὲν διὰ πλῆθος οἰκετῶν Καρικῶν εἰρῆσθαι φασὶν, ὡς ἐν τοῖς Ἀνθεστηρίοις εὐωχουμένων αὐτῶν καὶ οὐκ ἐργαζομένων ; \textit{idem} Diogenian.~5.23 (διὰ πλῆθος οἰκετῶν ; cetera om.) ; Ar.~\textit{Eq.}~1151 ἄπαγ᾽ ἐς μακαρίαν ἐκποδών ; Zen.~1.90 Ἀνδρὸς κακῶς πράσσοντος ἐκποδὼν φίλοι}
\lp{293 ἄπιθ᾽ … ἐκποδών (δούλοις, v. 294)]}{Ar.~\textit{Thesm.}~537 αὐταί τε καὶ τὰ δουλάρια τέφραν ποθὲν λαβοῦσαι ; Is.~6.50 ἣν οὔτε παρελθεῖν εἴσω τοῦ ἱεροῦ ἔδει οὔτ᾽ ἰδεῖν τῶν ἔνδον οὐδέν ; Ar.~\textit{Thesm.}~541 παρρησίας κἀξὸν λέγειν ὅσαι πάρεσμεν ἀσταί ; Ar.~\textit{Thesm.}~609 ἔχουσα; — τίτθη νὴ Δί᾽ ἐμή. — διοίχομαι ; Ar.~\textit{Thesm.}~329–330 τελέως δ᾽ / ἐκκλησιάσαιμεν Ἀθηναίων / εὐγενεῖς γυναῖκες ; Ar.~\textit{Thesm.}~728 ἴωμεν ἐπὶ τὰς κληματίδας ὦ Μανία ; Ar.~\textit{Ran.}~354–356 εὐφημεῖν χρὴ κἀξίστασθαι τοῖς ἡμετέροισι χοροῖσιν, / ὅστις ἄπειρος τοιῶνδε λόγων ἢ γνώμῃ μὴ καθαρεύει}
